% Leçon 4 — Éléments socioculturels : processus d’hérédité au Cameroun et en Chine — Chinois Tle A — Cours signé OnBuch+
% Programme officiel MINESEC. Compiler avec : tectonic ZH04-socio-processus-dheredite-au-cameroun-et-en-chine.tex
\newcommand{\DOCMATIERE}{Chinois}
\newcommand{\DOCNIVEAU}{Tle A}
\newcommand{\DOCMODULE}{Module 1 : Vie familiale et intégration sociale}
\newcommand{\DOCLECON}{ZH04}
\newcommand{\DOCTITRE}{Leçon 4 — Éléments socioculturels : processus d’hérédité au Cameroun et en Chine}
% OnBuch+ — préambule « cours signés » ENRICHI (compilé avec tectonic / XeLaTeX)
% Système visuel « Pop » d'OnBuch+ : fond crème (jamais blanc), bordures encre
% épaisses, ombres « dures » décalées, titres Archivo Black en capitales.
% Version MATHÉMATIQUES : graphes (pgfplots/tikz), boîtes pédagogiques
% (définition, propriété, méthode, attention, le savais-tu, exercice résolu,
% exercice, TP, bilan), page de garde et signature OnBuch+ ancrée en bas.
%
% Macros attendues dans main.tex AVANT \input{preamble} :
%   \DOCMATIERE  \DOCNIVEAU  \DOCMODULE  \DOCLECON (numéro)  \DOCTITRE
\documentclass[11pt]{article}

\usepackage{fontspec}
\usepackage[french]{babel} % français = langue principale ; le chinois via \zh{..} / textebox (xeCJK)
\usepackage{xeCJK}
\usepackage{pifont} % \ding{51} \ding{55} utilisés par les modèles
\usepackage[a4paper,margin=2cm,top=3.4cm,bottom=2.8cm,headheight=1.4cm,footskip=1.3cm]{geometry}
\usepackage{amsmath,amssymb,amsfonts}
\usepackage{mathtools} % \xmapsto, \coloneqq... souvent utilisés par les modèles
\usepackage{cancel} % simplifications barrées (bilans de forces)
% \abs{z} (module, valeur absolue) et \norm{u} (norme), utilisés par les modèles
\providecommand{\abs}{}\let\abs\relax\DeclarePairedDelimiter{\abs}{\lvert}{\rvert}
\providecommand{\norm}{}\let\norm\relax\DeclarePairedDelimiter{\norm}{\lVert}{\rVert}
\usepackage[table]{xcolor} % [table] : fournit aussi \rowcolors (alternance de lignes)
\usepackage{pagecolor}
\usepackage{tikz}
\usetikzlibrary{arrows.meta,positioning,calc,shapes.geometric,shapes.misc,shapes.arrows,decorations.pathmorphing,decorations.pathreplacing,decorations.markings,patterns,shadows,mindmap,trees,fit,backgrounds,matrix,intersections,angles,quotes}
\usepackage{pgfplots}
\usepackage[version=4]{mhchem} % équations nucléaires \ce{^{235}_{92}U}
\usepackage[european,siunitx]{circuitikz} % schémas électriques
\pgfplotsset{compat=1.18}
\usepgfplotslibrary{fillbetween,groupplots} % aires entre courbes (calcul intégral)
\usepackage{enumitem}
\usepackage{fancyhdr}
% [most] charge toutes les bibliothèques tcolorbox (listings, minted, raster,
% documentation...) et gonfle inutilement la mémoire TeX sur un document long
% (~300 boîtes) : on ne charge que ce qui est réellement utilisé.
\usepackage[skins,breakable]{tcolorbox}
\usepackage{graphicx}
\usepackage{colortbl}
\usepackage{booktabs}
\usepackage{tabularx}
\usepackage{diagbox} % case d'en-tête diagonale des tableaux à double entrée
% Colonnes extensibles centrée / gauche / droite (C, L, R), souvent utilisées par les modèles
\newcolumntype{C}{>{\centering\arraybackslash}X}
\newcolumntype{L}{>{\raggedright\arraybackslash}X}
\newcolumntype{R}{>{\raggedleft\arraybackslash}X}
\usepackage{array}
\usepackage{multicol}
\usepackage{multirow}
\usepackage{siunitx}
% parse-numbers=false : accepte aussi \SI{2,5 \times 10^{-3}}{...}
\sisetup{locale=FR,output-decimal-marker={,},group-minimum-digits=4,parse-numbers=false}
\DeclareSIUnit\ohm{\ensuremath{\symup{\Omega}}} % U+2126 absent de la police
\DeclareSIUnit\division{div} % réglages d'oscilloscope (ms/div, V/div)
\DeclareSIUnit\tr{tr} % tours (vitesse de rotation, ex. \SI{1500}{\tr\per\minute})
\DeclareSIUnit\molar{M} % concentration molaire (ex. \SI{3}{\molar}, \SI{1}{\micro\molar})
\DeclareSIUnit\an{an} % année (le modèle confond parfois avec \year, un compteur TeX) — ex. \SI{5}{\kilo\meter\per\an}
\DeclareSIUnit\FCFA{FCFA} % franc CFA (ex. \SI{500}{\FCFA\per\kilo\gram})
\DeclareSIUnit\degreeBrix{\textdegree Brix} % teneur en sucre d'un sirop/jus (réfractométrie)
\DeclareSIUnit\month{mois} % le modèle utilise parfois \month, un compteur TeX, comme unité de durée
\usepackage{qrcode}
% \degree : commande fréquemment utilisée par les modèles pour un angle ou une
% température en dehors de \SI{}{} (non fournie par les paquets chargés).
\providecommand{\degree}{^{\circ}}
% \qed (fin de démonstration), utilisé par les modèles sans amsthm
\providecommand{\qed}{\hfill$\square$}
% \barre : double barre des valeurs interdites dans un tableau de variations (tkz-tab)
\providecommand{\barre}{\ensuremath{\|}}
% environnement proof (amsthm non chargé) utilisé par les modèles
\ifdefined\proof\else\newenvironment{proof}[1][Démonstration]{\par\noindent\textit{#1.}\ }{\qed\par}\fi
\usepackage{lastpage}
\usepackage{hyperref}
\hypersetup{hidelinks}

% ============================================================
% Palette « Pop » OnBuch+ — mêmes hex que la classe P (plus_ui.dart)
% ============================================================
\definecolor{poporange}{HTML}{F59321}
\definecolor{popdark}{HTML}{E07A0C}
\definecolor{popcream}{HTML}{FFF6E8}
\definecolor{popink}{HTML}{1C1714}
\definecolor{poppurple}{HTML}{7A5AE0}
\definecolor{popgreen}{HTML}{2E9E6A}
\definecolor{poppink}{HTML}{E0457B}
\definecolor{popblue}{HTML}{2F7DE1}
\definecolor{popgold}{HTML}{C98A12}
\definecolor{popmuted}{HTML}{8A7F76}
\definecolor{popline}{HTML}{EADFD2}
\definecolor{popgray}{HTML}{8A7F76}
\colorlet{popgrey}{popgray}
% Alias tolérants (le modèle nomme parfois les couleurs différemment)
\colorlet{popwhite}{white}
\colorlet{popblack}{popink}
\colorlet{popred}{poppink}
\colorlet{popyellow}{popgold}
% Teintes claires (fonds de tableaux/encadrés) — le modèle nomme parfois
% les couleurs avec un suffixe L (ex. popblueL) sans les avoir définies.
\colorlet{poporangeL}{poporange!20}
\colorlet{popdarkL}{popdark!20}
\colorlet{poppurpleL}{poppurple!20}
\colorlet{popgreenL}{popgreen!20}
\colorlet{poppinkL}{poppink!20}
\colorlet{popblueL}{popblue!20}
\colorlet{popgoldL}{popgold!20}
\colorlet{popredL}{popred!20}
\colorlet{popgrayL}{popgray!20}
% Teintes claires pour fonds de tableaux / zones
\colorlet{popblueL}{popblue!10!white}
\colorlet{poporangeL}{poporange!14!white}
\colorlet{popgreenL}{popgreen!12!white}
\colorlet{poppurpleL}{poppurple!12!white}
\colorlet{poppinkL}{poppink!12!white}
\colorlet{popgoldL}{popgold!14!white}

\pagecolor{popcream}

% ============================================================
% Typographie
% ============================================================
\defaultfontfeatures{Ligatures=TeX}
\setmainfont{PlusJakartaSans-Medium.ttf}[Path=../../../fonts/,BoldFont=PlusJakartaSans-Bold.ttf]
% Caractères chinois : WenQuanYi Zen Hei (sous-ensemble GB2312 + pinyin), gras/italique simulés
\setCJKmainfont{WenQuanYiZenHei-subset.ttf}[Path=../../../fonts/,AutoFakeBold=3,AutoFakeSlant=0.2]
\xeCJKsetup{CJKspace=false}
% Pinyin : voyelles à caron (ǎ ǐ ǒ ǔ ǖ ǘ ǚ ǜ) absentes de la police principale -> caractères actifs
\newfontface\pyfont{WenQuanYiZenHei-subset.ttf}[Path=../../../fonts/]
\newcommand\pyactive[1]{\catcode#1=13 \begingroup\lccode`\~=#1 \lowercase{\endgroup\def~}{{\pyfont\char#1}}}
\pyactive{"01CD}\pyactive{"01CE}\pyactive{"01CF}\pyactive{"01D0}\pyactive{"01D1}\pyactive{"01D2}\pyactive{"01D3}\pyactive{"01D4}\pyactive{"01D5}\pyactive{"01D6}\pyactive{"01D7}\pyactive{"01D8}\pyactive{"01D9}\pyactive{"01DA}\pyactive{"01DB}\pyactive{"01DC}
% Maths en sans-serif (Fira Math) avec chiffres et lettres droites de Plus
% Jakarta, pour que les formules se fondent dans le texte.
\usepackage[math-style=ISO,bold-style=ISO]{unicode-math}
\setmathfont{FiraMath-Regular.otf}[Path=../../../fonts/]
\setmathfont{PlusJakartaSans-Medium.ttf}[Path=../../../fonts/,range={up/num,up/{Latin,latin},\mathup}]
\usepackage{icomma} % virgule décimale sans espace en mode math
% Caractères mathématiques Unicode absents des polices de texte (Plus Jakarta,
% Archivo Black) : rendus par leur équivalent mathématique, en texte comme en
% formule — sinon ils s'affichent en carré vide (ex. « Arithmétique dans ℤ »).
\usepackage{newunicodechar}
\newunicodechar{ℝ}{\ensuremath{\mathbb{R}}}
\newunicodechar{ℤ}{\ensuremath{\mathbb{Z}}}
\newunicodechar{ℂ}{\ensuremath{\mathbb{C}}}
\newunicodechar{ℕ}{\ensuremath{\mathbb{N}}}
\newunicodechar{ℚ}{\ensuremath{\mathbb{Q}}}
\newunicodechar{𝔻}{\ensuremath{\mathbb{D}}}
\newunicodechar{α}{\ensuremath{\alpha}}
\newunicodechar{β}{\ensuremath{\beta}}
\newunicodechar{γ}{\ensuremath{\gamma}}
\newunicodechar{δ}{\ensuremath{\delta}}
\newunicodechar{ε}{\ensuremath{\varepsilon}}
\newunicodechar{θ}{\ensuremath{\theta}}
\newunicodechar{λ}{\ensuremath{\lambda}}
\newunicodechar{μ}{\ensuremath{\mu}}
\newunicodechar{σ}{\ensuremath{\sigma}}
\newunicodechar{τ}{\ensuremath{\tau}}
\newunicodechar{φ}{\ensuremath{\varphi}}
\newunicodechar{ψ}{\ensuremath{\psi}}
\newunicodechar{ω}{\ensuremath{\omega}}
\newunicodechar{Σ}{\ensuremath{\Sigma}}
% majuscules produites par \MakeUppercase dans les titres (α-aminés -> Α)
\newunicodechar{Α}{\ensuremath{\alpha}}
\newunicodechar{Β}{\ensuremath{\beta}}
\newunicodechar{Γ}{\ensuremath{\Gamma}}
\newunicodechar{Δ}{\ensuremath{\Delta}}
\newunicodechar{Ε}{\ensuremath{\varepsilon}}
\newunicodechar{Θ}{\ensuremath{\Theta}}
\newunicodechar{Λ}{\ensuremath{\Lambda}}
\newunicodechar{Μ}{\ensuremath{\mu}}
\newunicodechar{Τ}{\ensuremath{\tau}}
\newunicodechar{Φ}{\ensuremath{\Phi}}
\newunicodechar{Ψ}{\ensuremath{\Psi}}
\newunicodechar{Ω}{\ensuremath{\Omega}}
\newunicodechar{↦}{\ensuremath{\mapsto}}
\newunicodechar{∈}{\ensuremath{\in}}
\newunicodechar{∉}{\ensuremath{\notin}}
\newunicodechar{∧}{\ensuremath{\wedge}}
\newunicodechar{⇔}{\ensuremath{\Leftrightarrow}}
\newunicodechar{⟺}{\ensuremath{\Longleftrightarrow}}
\newunicodechar{⇒}{\ensuremath{\Rightarrow}}
\newunicodechar{⟹}{\ensuremath{\Longrightarrow}}
\newunicodechar{∘}{\ensuremath{\circ}}
\newunicodechar{∪}{\ensuremath{\cup}}
\newunicodechar{∩}{\ensuremath{\cap}}
\newunicodechar{≡}{\ensuremath{\equiv}}
\newunicodechar{⊂}{\ensuremath{\subset}}
\newunicodechar{⊥}{\ensuremath{\perp}}
\newunicodechar{∃}{\ensuremath{\exists}}
\newunicodechar{∀}{\ensuremath{\forall}}
\newunicodechar{∛}{\ensuremath{\sqrt[3]{\,}}}
\newunicodechar{∜}{\ensuremath{\sqrt[4]{\,}}}
\newunicodechar{⃗}{\ensuremath{{}^{\to}}}
\newunicodechar{ⁿ}{\ifmmode^{n}\else\textsuperscript{\ensuremath{n}}\fi}
\newunicodechar{ˣ}{\ifmmode^{x}\else\textsuperscript{\ensuremath{x}}\fi}
\newunicodechar{ᵇ}{\ifmmode^{b}\else\textsuperscript{\ensuremath{b}}\fi}
\newunicodechar{ʳ}{\ifmmode^{r}\else\textsuperscript{\ensuremath{r}}\fi}
\newunicodechar{ᵘ}{\ifmmode^{u}\else\textsuperscript{\ensuremath{u}}\fi}
\newunicodechar{ʸ}{\ifmmode^{y}\else\textsuperscript{\ensuremath{y}}\fi}
\newunicodechar{ᵃ}{\ifmmode^{a}\else\textsuperscript{\ensuremath{a}}\fi}
\newunicodechar{ᵉ}{\ifmmode^{e}\else\textsuperscript{\ensuremath{e}}\fi}
\newunicodechar{ᵏ}{\ifmmode^{k}\else\textsuperscript{\ensuremath{k}}\fi}
\newunicodechar{ᵖ}{\ifmmode^{p}\else\textsuperscript{\ensuremath{p}}\fi}
\newunicodechar{ᴺ}{\ifmmode^{N}\else\textsuperscript{\ensuremath{N}}\fi}
\newunicodechar{ᶜ}{\ifmmode^{c}\else\textsuperscript{\ensuremath{c}}\fi}
\newunicodechar{ʰ}{\ifmmode^{h}\else\textsuperscript{\ensuremath{h}}\fi}
\newunicodechar{ᵠ}{\ifmmode^{q}\else\textsuperscript{\ensuremath{q}}\fi}
\newunicodechar{ₙ}{\ifmmode_{n}\else\textsubscript{\ensuremath{n}}\fi}
\newunicodechar{ₐ}{\ifmmode_{a}\else\textsubscript{\ensuremath{a}}\fi}
\newunicodechar{ᵢ}{\ifmmode_{i}\else\textsubscript{\ensuremath{i}}\fi}
\newunicodechar{ᵣ}{\ifmmode_{r}\else\textsubscript{\ensuremath{r}}\fi}
\newunicodechar{ₖ}{\ifmmode_{k}\else\textsubscript{\ensuremath{k}}\fi}
\newunicodechar{ᵦ}{\ifmmode_{\beta}\else\textsubscript{\ensuremath{\beta}}\fi}
\newunicodechar{ₓ}{\ifmmode_{x}\else\textsubscript{\ensuremath{x}}\fi}
\newfontfamily\popdisplay{ArchivoBlack-Regular.ttf}[Path=../../../fonts/]
\newfontfamily\poplogo{SpaceGrotesk-Bold.ttf}[Path=../../../fonts/]
\newfontfamily\popsemi{PlusJakartaSans-SemiBold.ttf}[Path=../../../fonts/]
\newfontfamily\popxbold{PlusJakartaSans-ExtraBold.ttf}[Path=../../../fonts/]
\newfontfamily\popxboldls{PlusJakartaSans-ExtraBold.ttf}[Path=../../../fonts/,LetterSpace=3.0]

\setlist[enumerate]{leftmargin=1.7em,itemsep=5pt,topsep=4pt}
\setlist[itemize]{leftmargin=1.7em,itemsep=4pt,topsep=4pt,label={\color{poporange}\textbullet}}
\setlength{\parindent}{0pt}
\setlength{\parskip}{5pt}
\renewcommand{\arraystretch}{1.25}

% pgfplots : style Pop par défaut
\pgfplotsset{
  every axis/.append style={axis line style={popink,line width=1pt},tick style={popink},
    grid style={popline,line width=0.5pt},label style={font=\small},tick label style={font=\footnotesize}},
  popcurve/.style={poporange,line width=1.8pt,no markers,smooth},
}
% Même style disponible dans un \draw TikZ simple (hors pgfplots)
\tikzset{popcurve/.style={poporange,line width=1.8pt}}
\tikzset{popgrid/.style={popline,very thin}}
\colorlet{popcurve}{poporange} % le nom du style est parfois utilisé comme couleur

% ============================================================
% Pill / squiggle
% ============================================================
\newcommand{\poppill}[3]{%
  \tikz[baseline=-0.55ex]\node[draw=popink,line width=1.2pt,rounded corners=2.6pt,%
    fill=#2,inner xsep=6.5pt,inner ysep=3pt]{%
    \popxboldls\fontsize{7}{7}\selectfont\color{#3}\MakeUppercase{#1}};%
}
\newcommand{\popsquiggle}[1]{%
  \tikz[baseline=-0.4ex]{
    \draw[#1,line width=2.6pt,line cap=round]
      plot [smooth] coordinates {(0,0.09) (0.34,0.01) (0.68,0.21) (1.02,0.01) (1.36,0.10)};
  }%
}
% Mot-clé surligné (marqueur orange sous le texte)
\newcommand{\cle}[1]{{\popxbold\color{popdark}#1}}

% ============================================================
% En-tête / pied
% ============================================================
\pagestyle{fancy}
\fancyhf{}
\renewcommand{\headrulewidth}{0pt}
\renewcommand{\footrulewidth}{0pt}
\lhead{\raisebox{-0.15ex}{\poplogo\fontsize{13}{13}\selectfont\color{popink}OnBuch\color{poporange}+}}
\rhead{\poppill{\DOCMATIERE\ \textbullet\ \DOCNIVEAU\ \textbullet\ Leçon \DOCLECON}{white}{popink}}
\lfoot{\popxbold\fontsize{7.6}{7.6}\selectfont\color{popmuted}\MakeUppercase{Cours signé OnBuch+}\ \textbullet\ {\popsemi\DOCTITRE}}
\rfoot{\tikz[baseline=-0.6ex]\node[rounded corners=8.5pt,fill=popink,minimum height=17pt,minimum width=17pt,inner xsep=5pt,inner sep=0]{\popxbold\fontsize{8}{8}\selectfont\color{popcream}\thepage\,/\,\pageref*{LastPage}};}

% ============================================================
% Titres
% ============================================================
\newcommand{\chaptitre}[2]{%
  \begin{tcolorbox}[enhanced,colback=poporange,colframe=popink,boxrule=2.6pt,arc=11pt,
      drop shadow=popink,left=15pt,right=15pt,top=12pt,bottom=11pt,before skip=3pt,after skip=18pt]
    \poppill{#2}{popink}{popcream}\par\vspace{9pt}
    {\raggedright\hyphenpenalty=10000\exhyphenpenalty=10000\popdisplay\fontsize{22}{24}\selectfont\color{popcream}\MakeUppercase{#1}\par}
  \end{tcolorbox}
}
% Section numérotée : pastille orange avec le numéro + titre Archivo
\newcounter{popsec}
\newcommand{\coursec}[1]{%
  \stepcounter{popsec}\setcounter{popsub}{0}%
  \par\vspace{16pt}\noindent
  \begin{minipage}{\linewidth}
  \tikz[baseline=-0.7ex]\node[draw=popink,line width=1.6pt,fill=poporange,rounded corners=4pt,minimum size=22pt,inner sep=0]{\popdisplay\fontsize{12}{12}\selectfont\color{popcream}\thepopsec};%
  \hspace{8pt}{\popdisplay\fontsize{15}{17}\selectfont\color{popink}\MakeUppercase{#1}}\par
  \vspace{4pt}\noindent{\color{poporange}\rule{36pt}{3.6pt}}
  \end{minipage}\par\nopagebreak\vspace{8pt}
}
\newcounter{popsub}
\newcommand{\courssub}[1]{%
  \stepcounter{popsub}%
  \par\vspace{9pt}\noindent{\popxbold\fontsize{11.5}{13}\selectfont\color{popink}{\color{poporange}\thepopsec.\thepopsub}\hspace{6pt}#1}\par\nopagebreak\vspace{4pt}
}

% ============================================================
% Boîtes pédagogiques — même recette : fond blanc, bordure épaisse colorée,
% ombre dure encre, badge plein. Titre optionnel [#1].
% ============================================================
\tcbset{popbase/.style={breakable,enhanced,colback=white,boxrule=2.3pt,arc=9pt,
  shadow={3pt}{-3pt}{0pt}{popink},left=14pt,right=14pt,top=11pt,bottom=11pt,before skip=11pt,after skip=15pt,
  fontupper=\popsemi\fontsize{10.6}{14.5}\selectfont\color{popink},
  fontlower=\popsemi\fontsize{10.6}{14.5}\selectfont\color{popink}}}
% Coche dessinée (le glyphe ✓ n'existe pas dans les polices du document)
\newcommand{\popcheck}[1]{\tikz[baseline=-0.5ex]\draw[#1,line width=1.6pt,line cap=round,line join=round](-0.13,0.0)--(-0.04,-0.09)--(0.13,0.1);}
\providecommand{\coche}{\popcheck{popgreen}}
\providecommand{\resultat}[1]{\par\smallskip\noindent\fcolorbox{popgreen}{popgreenL}{\parbox{\dimexpr\linewidth-2\fboxsep-2\fboxrule\relax}{\textbf{Résultat.} #1}}\par\smallskip}
\newcommand{\popboxhead}[3]{\poppill{#1}{#2}{white}\ifx\relax#3\relax\else\hspace{6pt}{\popxbold\fontsize{10.6}{12}\selectfont\color{popink}#3}\fi\par\vspace{7pt}}

\newtcolorbox{aretenir}[1][]{popbase,colframe=popgreen,before upper={\popboxhead{À retenir}{popgreen}{#1}}}
% Texte en chinois (document de lecture, dialogue, extrait) : cadre
% d'étude. Un titre optionnel (source, auteur) s'affiche dans l'en-tête.
\newtcolorbox{textebox}[1][]{popbase,colframe=popblue,before upper={\popboxhead{文本 · Texte}{popblue}{#1}},after upper={}}
% Marques ✓ / ✗ (forme correcte / fautive) souvent utilisées par les modèles
\providecommand{\cross}{\textcolor{poppink}{\ensuremath{\times}}}
\providecommand{\xmark}{\cross}\providecommand{\crossmark}{\cross}\providecommand{\cmark}{\ensuremath{\checkmark}}
% Ligne de réponse / justification à compléter (inventée par les modèles)
\providecommand{\justification}[1]{\par\noindent\textit{Justification~:} \dotfill\par}
% \textipa (package tipa non chargé) : la police ne couvre pas l'API -> simple passage
\providecommand{\textipa}[1]{#1}
% Soulignement (ulem non chargé) : \uline, \ul
\providecommand{\uline}[1]{\underline{#1}}\providecommand{\ul}[1]{\underline{#1}}
% \glossaire{\item ...} inventée par les modèles : liste de glossaire
\providecommand{\glossaire}[1]{\begin{itemize}#1\end{itemize}}
% \alert (beamer) inventée par les modèles : texte mis en évidence
\providecommand{\alert}[1]{\textbf{\textcolor{poppink}{#1}}}
% variantes ASCII d'ordinaux écrites en macro
\providecommand{\iere}{\textsuperscript{re}}\providecommand{\ieme}{\textsuperscript{e}}\providecommand{\ere}{\textsuperscript{re}}\providecommand{\eme}{\textsuperscript{e}}\providecommand{\er}{\textsuperscript{er}}\providecommand{\ie}{\textsuperscript{e}}
% Ordinaux français écrits en macro par les modèles : 1\ière, 3\ième
\newcommand{\ière}{\textsuperscript{re}}\newcommand{\ième}{\textsuperscript{e}}
% \exemple (inventée par les modèles) : étiquette « Exemple : »
\providecommand{\exemple}{\textbf{Exemple~:}\ }
% \square utilisable aussi hors mode math (cases à cocher)
\AtBeginDocument{\let\squareMath\square\renewcommand{\square}{\ensuremath{\squareMath}}}
% Mot ou phrase en chinois à l'intérieur d'un paragraphe français
\newcommand{\zh}[1]{\ifmmode\text{#1}\else{#1}\fi}
\let\eng\zh \let\esp\zh \let\all\zh % alias tolérés
% flèches d'intonation inventées par les modèles
\providecommand{\textsearrow}{}\renewcommand{\textsearrow}{\ensuremath{\searrow}}\providecommand{\textnearrow}{}\renewcommand{\textnearrow}{\ensuremath{\nearrow}}
% \fr{..} : mot français (inventé par les modèles)
\providecommand{\fr}[1]{\textit{#1}}
% zhbox : encadré de texte chinois inventé par les modèles
\newenvironment{zhbox}{\begin{quote}}{\end{quote}}
% \blank : case à compléter inventée par les modèles
\providecommand{\blank}[1][]{\underline{\hspace{1.6em}}}
\newtcolorbox{exemplebox}[1][]{popbase,colframe=poporange,before upper={\popboxhead{Exemple}{poporange}{#1}}}
\newtcolorbox{definition}[1][]{popbase,colframe=popblue,before upper={\popboxhead{Définition}{popblue}{#1}}}
\newtcolorbox{propriete}[1][]{popbase,colframe=poppurple,before upper={\popboxhead{Propriété}{poppurple}{#1}}}
\newtcolorbox{methode}[1][]{popbase,colframe=popgold,before upper={\popboxhead{Méthode}{popgold}{#1}}}
\newtcolorbox{attention}[1][]{popbase,colframe=poppink,before upper={\popboxhead{Attention}{poppink}{#1}}}
\newtcolorbox{experience}[1][]{popbase,colframe=popblue,colback=popblueL,before upper={\popboxhead{Expérience}{popblue}{#1}}}
% « Le savais-tu ? » : carte violette pleine (contexte, culture, Cameroun)
\newtcolorbox{savaistu}[1][]{popbase,colback=poppurple,colframe=popink,
  fontupper=\popsemi\fontsize{10.4}{14.2}\selectfont\color{white},
  before upper={\poppill{Le savais-tu\,?}{white}{poppurple}\ifx\relax#1\relax\else\hspace{6pt}{\popxbold\color{white}#1}\fi\par\vspace{7pt}}}
% Exercice résolu : énoncé en haut, \tcblower puis la solution sur fond crème
\newcounter{popexo}\newcounter{popexr}
\newtcolorbox{exoresolu}[1][]{popbase,colframe=poporange,colbacklower=popcream,
  segmentation style={popink,line width=1.2pt,dashed},
  before upper={\stepcounter{popexr}\popboxhead{Exercice résolu \thepopexr}{poporange}{#1}},
  before lower={\poppill{Solution}{popink}{popcream}\par\vspace{6pt}}}
% Exercice (non résolu en place) — la correction est dans la partie Corrigés
\newtcolorbox{exercice}[1][]{popbase,colframe=popink,
  before upper={\stepcounter{popexo}\popboxhead{Exercice \thepopexo}{popink}{#1}}}
% Corrigé
\newtcolorbox{corrige}[1][]{popbase,colframe=popgreen,colback=popgreenL,
  before upper={\popboxhead{Corrigé}{popgreen}{#1}}}
% Objectifs / prérequis / bilan
\newtcolorbox{objectifs}{popbase,colframe=popink,colback=poporangeL,
  before upper={\popboxhead{Objectifs de la leçon}{popink}{}}}
\newtcolorbox{prerequis}{popbase,colframe=popink,colback=popblueL,
  before upper={\popboxhead{Prérequis}{popblue}{}}}
\newtcolorbox{bilan}{popbase,colframe=popink,colback=white,boxrule=2.8pt,
  before upper={\popboxhead{Fiche bilan}{poporange}{L'essentiel à connaître pour l'examen}}}
% Figure encadrée avec légende
\newtcolorbox{popfigure}[1][]{enhanced,breakable=false,colback=white,colframe=popline,boxrule=1.4pt,arc=8pt,
  left=10pt,right=10pt,top=10pt,bottom=8pt,before skip=10pt,after skip=12pt,halign=center,
  lower separated=false,halign lower=center,
  fontlower=\popsemi\fontsize{9}{11.5}\selectfont\color{popmuted},
  #1}
\newcommand{\legende}[1]{\par\vspace{6pt}{\popsemi\fontsize{9}{11.5}\selectfont\color{popmuted}#1}}

% ============================================================
% Signature OnBuch+
% ============================================================
\newcommand{\poplogoinline}[1]{{\poplogo\fontsize{#1}{#1}\selectfont\color{popink}OnBuch\color{poporange}+}}
\newcommand{\signatureonbuch}{%
  \begin{tcolorbox}[enhanced,colback=popink,colframe=popink,boxrule=2.2pt,arc=10pt,
      drop shadow=poporange,left=14pt,right=14pt,top=10pt,bottom=10pt,before skip=0pt,after skip=0pt]
    \tikz[baseline=-0.7ex]\node[circle,fill=popgreen,draw=popcream,line width=1.3pt,minimum size=22pt,inner sep=0]{\popcheck{white}};%
    \hspace{9pt}\begin{minipage}[c]{0.62\linewidth}
      {\popxbold\fontsize{12}{13}\selectfont\color{popcream}Signé\ \textbullet\ {\poplogo OnBuch\color{poporange}+}}\par\vspace{2pt}
      {\raggedright\popsemi\fontsize{8}{10.5}\selectfont\color{popline}\MakeUppercase{Équipe pédagogique OnBuch+}\\\MakeUppercase{Conforme au programme officiel MINESEC}\par}
    \end{minipage}\hfill
    {\poplogo\fontsize{22}{22}\selectfont\color{popcream}OnBuch\color{poporange}+}
  \end{tcolorbox}%
}

% ============================================================
% Page de garde
% #1 titre · #2 matière · #3 niveau · #4 module · #5 n° leçon · #6 durée ·
% #7 description / objectifs (texte ou itemize)
% ============================================================
\newcommand{\pagegarde}[7]{%
  \thispagestyle{empty}
  \newgeometry{a4paper,margin=1.7cm,top=1.6cm,bottom=1.4cm}%
  \begin{tikzpicture}[remember picture,overlay]
    \fill[poporange!18] ([xshift=-3.2cm,yshift=-3.6cm]current page.north east) circle (4.6cm);
    \fill[poppurple!14] ([xshift=2.4cm,yshift=7.6cm]current page.south west) circle (3.8cm);
  \end{tikzpicture}%
  {\poplogo\fontsize{30}{30}\selectfont\color{popink}OnBuch\color{poporange}+}\hfill
  \raisebox{0.6ex}{\poppill{Cours signé \textbullet\ Premium}{popink}{popcream}}\par
  \vspace{1pt}\popsquiggle{poppurple}\par
  \vspace{8pt}
  \begin{tcolorbox}[enhanced,colback=poporange,colframe=popink,boxrule=2.8pt,arc=13pt,
      drop shadow=popink,left=18pt,right=18pt,top=12pt,bottom=13pt,after skip=10pt]
    \poppill{#2 \textbullet\ #3}{popink}{popcream}\hspace{5pt}\poppill{#4}{white}{popink}\par\vspace{8pt}
    {\popdisplay\fontsize{14}{14}\selectfont\color{popink}LEÇON #5}\par\vspace{4pt}
    {\raggedright\hyphenpenalty=10000\exhyphenpenalty=10000\popdisplay\fontsize{25}{27}\selectfont\color{popcream}\MakeUppercase{#1}\par}
    \vspace{8pt}
    {\popxbold\fontsize{9.5}{9.5}\selectfont\color{popink}\MakeUppercase{Durée conseillée : #6}}
  \end{tcolorbox}
  \begin{tcolorbox}[enhanced,colback=white,colframe=popink,boxrule=2.2pt,arc=10pt,
      drop shadow=popink,left=16pt,right=16pt,top=10pt,bottom=10pt,after skip=8pt]
    \poppill{Dans cette leçon}{poppurple}{white}\par\vspace{5pt}
    \popsemi\fontsize{9.3}{12.6}\selectfont\color{popink}#7
  \end{tcolorbox}
  \noindent\begin{minipage}[t]{0.487\linewidth}
    \begin{tcolorbox}[enhanced,colback=popgreen,colframe=popink,boxrule=2.2pt,arc=10pt,
        drop shadow=popink,left=11pt,right=11pt,top=10pt,bottom=10pt,height=3.1cm,valign=center]
      \tcbox[colback=white,colframe=popink,boxrule=1.3pt,arc=5pt,boxsep=0pt,nobeforeafter,
        left=3pt,right=3pt,top=3pt,bottom=3pt,valign=center]{\qrcode[height=1.9cm]{https://play.google.com/store/apps/details?id=com.luvvix.onbuch&hl=fr}}%
      \hspace{8pt}\begin{minipage}[c]{0.5\linewidth}
        {\popdisplay\fontsize{11}{12.5}\selectfont\color{white}\MakeUppercase{Télécharge OnBuch}}\par\vspace{4pt}
        {\raggedright\popsemi\fontsize{8.4}{11}\selectfont\color{white}Gratuit \textbullet\ Google Play\\Tuteur IA, annales, concours, résultats\par}
      \end{minipage}
    \end{tcolorbox}
  \end{minipage}\hfill
  \begin{minipage}[t]{0.487\linewidth}
    \begin{tcolorbox}[enhanced,colback=poppurple,colframe=popink,boxrule=2.2pt,arc=10pt,
        drop shadow=popink,left=11pt,right=11pt,top=10pt,bottom=10pt,height=3.1cm,valign=center]
      \tikz[baseline=-0.7ex]\node[circle,fill=white,draw=popink,line width=1.3pt,minimum size=1.6cm,inner sep=0]{\popdisplay\fontsize{24}{24}\selectfont\color{poppurple}+};%
      \hspace{8pt}\begin{minipage}[c]{0.6\linewidth}
        {\popdisplay\fontsize{11}{12.5}\selectfont\color{white}\MakeUppercase{Passe à OnBuch+}}\par\vspace{4pt}
        {\raggedright\popsemi\fontsize{8.4}{11}\selectfont\color{white}Tous les cours signés, Tuteur IA illimité, fiches PDF — onglet OnBuch+ de l'app\par}
      \end{minipage}
    \end{tcolorbox}
  \end{minipage}
  \vspace{8pt}
  \popcontenu
  \vfill
  \signatureonbuch
  \restoregeometry
  \newpage
}

% Rangée « Ce cours contient » (page de garde)
\newcommand{\poptile}[3]{% #1 couleur #2 gros texte #3 légende
  \begin{tcolorbox}[enhanced,colback=#1,colframe=popink,boxrule=2pt,arc=9pt,shadow={2.5pt}{-2.5pt}{0pt}{popink},
      left=6pt,right=6pt,top=9pt,bottom=9pt,halign=center,valign=center,height=1.75cm,width=\linewidth,nobeforeafter]
    {\popdisplay\fontsize{12.5}{13}\selectfont\color{white}\MakeUppercase{#2}}\par\vspace{3pt}
    {\centering\popsemi\fontsize{7.8}{9.5}\selectfont\color{white}#3\par}
  \end{tcolorbox}}
\newcommand{\popcontenu}{%
  \par\noindent{\popxboldls\fontsize{7.5}{7.5}\selectfont\color{popmuted}CE COURS SIGNÉ CONTIENT}\par\vspace{7pt}
  \noindent\begin{minipage}[t]{0.235\linewidth}\poptile{popblue}{Cours}{complet, illustré, pas à pas}\end{minipage}\hfill
  \begin{minipage}[t]{0.235\linewidth}\poptile{poporange}{Méthodes}{et exercices résolus}\end{minipage}\hfill
  \begin{minipage}[t]{0.235\linewidth}\poptile{poppink}{Exercices}{type examen + corrigés}\end{minipage}\hfill
  \begin{minipage}[t]{0.235\linewidth}\poptile{popgreen}{Fiche bilan}{l'essentiel à retenir}\end{minipage}\par}

% Sommaire
\newtcolorbox{sommaire}{enhanced,colback=white,colframe=popink,boxrule=2.2pt,arc=10pt,
  drop shadow=popink,left=18pt,right=18pt,top=13pt,bottom=13pt,before skip=6pt,after skip=20pt,
  before upper={\poppill{Sommaire}{poppurple}{white}\par\vspace{9pt}\popsemi\fontsize{10.6}{14.5}\selectfont\color{popink}}}

% Signature de fin de cours — poussée tout en bas de la dernière page
\newcommand{\findecours}{%
  \par\vfill\nopagebreak
  \begin{center}\popsquiggle{poporange}\end{center}\vspace{4pt}
  \signatureonbuch
}
\AtBeginDocument{\def\llbracket{\mathopen{[\mkern-3.5mu[}}\def\rrbracket{\mathclose{]\mkern-3.5mu]}}} % intervalles d'entiers (glyphe ⟦ absent de la police)
% Glyphes absents de Fira Math : redéfinis à partir de glyphes disponibles
\AtBeginDocument{%
  \def\setminus{\mathbin{\backslash}}\let\smallsetminus\setminus
  \def\vdots{\mathord{\vbox{\baselineskip3.2pt\lineskiplimit0pt\kern4pt\hbox{.}\hbox{.}\hbox{.}}}}%
  \def\ddots{\mathinner{\mkern1mu\raise6.5pt\hbox{.}\mkern2mu\raise3.6pt\hbox{.}\mkern2mu\raise0.7pt\hbox{.}\mkern1mu}}%
  \def\triangle{\mathord{\scalebox{0.9}{$\symup{\Delta}$}}}%
  \def\nabla{\mathord{\raisebox{\depth}{\scalebox{1}[-1]{$\symup{\Delta}$}}}}%
  \def\checkmark{\popcheck{popdark}}%
  \def\star{\mathbin{\ast}}\def\bigstar{\mathord{\ast}}%
  \def\diamond{\mathbin{\scalebox{0.6}{$\Diamond$}}}%
  \def\bigcup{\mathop{\mathchoice{\raisebox{-0.3em}{\scalebox{1.7}{$\cup$}}}{\scalebox{1.25}{$\cup$}}{\cup}{\cup}}\displaylimits}%
  \def\bigcap{\mathop{\mathchoice{\raisebox{-0.3em}{\scalebox{1.7}{$\cap$}}}{\scalebox{1.25}{$\cap$}}{\cap}{\cap}}\displaylimits}%
  \def\lesseqqgtr{\mathrel{\lessgtr}}\def\bigoplus{\mathop{\oplus}\displaylimits}%
}
\providecommand{\ssi}{si et seulement si} % abréviation fréquente
\AtBeginDocument{\providecommand{\wideparen}{\overparen}} % arc de cercle

\begin{document}

\pagegarde{Leçon 4 — Éléments socioculturels : processus d’hérédité au Cameroun et en Chine}{Chinois}{Tle A}{Module 1 : Vie familiale et intégration sociale}{ZH04}{2 h}{Cette leçon socioculturelle compare de façon factuelle et respectueuse les processus d'hérédité au Cameroun et en Chine : traditions coutumières, droit moderne, place des fils et des filles. Elle fournit le vocabulaire chinois associé et prépare à un mini-exposé argumenté.
\vspace{4pt}
\begin{multicols}{2}\small
\begin{itemize}
\item À la fin de cette leçon, je sais définir ce qu'est l'hérédité et la succession dans les deux cultures
\item Je sais décrire les pratiques traditionnelles camerounaises en matière d'héritage
\item Je sais décrire l'évolution historique de l'héritage en Chine (tradition confucéenne, loi de 1985)
\item Je sais comparer les deux systèmes avec respect et esprit critique
\item Je sais employer le vocabulaire socioculturel du thème en caractères et en pinyin
\item Je sais exprimer un point de vue simple en chinois avec 我觉得, 我认为
\end{itemize}\end{multicols}}

\begin{sommaire}
\begin{multicols}{2}
\begin{enumerate}
\item Introduction : qu'est-ce que l'hérédité ?
\item L'hérédité au Cameroun
\item L'hérédité en Chine
\item Comparaison Cameroun / Chine
\item Vocabulaire et écriture des caractères
\item Réflexion et mini-exposé
\item Activité d'intégration
\item Méthodes et exercices résolus
\item Exercices
\item Corrigés des exercices
\item Fiche bilan
\end{enumerate}
\end{multicols}
\end{sommaire}

\begin{objectifs}
\begin{itemize}
\item À la fin de cette leçon, je sais définir ce qu'est l'hérédité et la succession dans les deux cultures
\item Je sais décrire les pratiques traditionnelles camerounaises en matière d'héritage
\item Je sais décrire l'évolution historique de l'héritage en Chine (tradition confucéenne, loi de 1985)
\item Je sais comparer les deux systèmes avec respect et esprit critique
\item Je sais employer le vocabulaire socioculturel du thème en caractères et en pinyin
\item Je sais exprimer un point de vue simple en chinois avec 我觉得, 我认为
\item Je sais présenter un mini-exposé de deux minutes sur l'héritage dans mon pays
\item Je sais répondre à des questions de réflexion sur l'égalité filles-garçons dans la succession
\end{itemize}
\end{objectifs}

\begin{prerequis}
\begin{itemize}
\item Vocabulaire de la famille (家人, 父亲, 母亲, 儿子, 女儿) vu en Première
\item Structures de comparaison avec 比 et 和……一样
\item Expression de l'opinion : 我觉得 / 我认为
\item Notions sur la famille et le mariage en Chine (leçons socio précédentes)
\end{itemize}
\end{prerequis}

\newpage

\coursec{Leçon 4 — Éléments socioculturels : processus d’hérédité au Cameroun et en Chine}

\courssub{Introduction : qu'est-ce que l'hérédité ?}

\courssub{Situation d'accroche : deux familles, une même question}

À Yaoundé, après le décès d'un chef de famille, toute la famille élargie se réunit pour décider qui héritera de la concession et des champs : l'aîné des fils ? Les enfants ensemble ? La veuve ? Les oncles du défunt ont souvent leur mot à dire, et les discussions peuvent durer des jours. En Chine aussi, la question de l'héritage a longtemps suivi des règles précises, souvent en faveur des fils, qui portaient le nom de la famille et entretenaient les tombes des ancêtres. Aujourd'hui, les lois modernes des deux pays reconnaissent des droits à tous les enfants, garçons et filles, mais les traditions restent vivaces, au village comme en ville.

\textbf{Problématique :} comment définir l'hérédité, distinguer héritage matériel et immatériel, et surtout, comment parler de tout cela en chinois ?

\courssub{Observation : un premier texte pour découvrir le mot 遗产}

Avant toute définition, lisons une phrase simple qui présente l'idée centrale de la leçon.

\begin{textebox}[Qu'est-ce que l'héritage ?]
遗产是人去世以后留下的东西，比如房子、土地和钱。\\
\emph{Yíchǎn shì rén qùshì yǐhòu liúxià de dōngxi, bǐrú fángzi, tǔdì hé qián.}\\
« L'héritage, ce sont les choses laissées après le décès d'une personne, comme la maison, la terre et l'argent. »
\end{textebox}

Observons la structure de cette phrase :

\begin{itemize}
\item \zh{是} (\emph{shì}, être) relie le sujet \zh{遗产} à sa définition, comme en français « l'héritage, c'est… ».
\item \zh{去世以后} (\emph{qùshì yǐhòu}) signifie « après le décès » : \zh{以后} (après) se place \textbf{après} le mot qu'il complète, contrairement au français « après le décès » où « après » précède.
\item \zh{比如} (\emph{bǐrú}) introduit des exemples : « par exemple, comme ».
\end{itemize}

\begin{definition}[遗产 yíchǎn : l'héritage]
\zh{遗产} (\emph{yíchǎn}) est un nom qui désigne l'\cle{héritage} : l'ensemble des biens laissés par une personne après son décès. Exemple : \zh{父亲留下了很多遗产。} (\emph{Fùqin liúxià le hěn duō yíchǎn.}) « Le père a laissé beaucoup de biens en héritage. »
\end{definition}

\courssub{Définition générale : qu'est-ce que l'hérédité ?}

\begin{definition}[L'hérédité]
L'\cle{hérédité}, au sens social et juridique, est la \textbf{transmission} des biens, des terres, des titres et parfois du statut social d'une personne après son décès, à des personnes appelées \textbf{héritiers}. En chinois, on parle de \zh{继承} (\emph{jìchéng}, hériter, transmettre) et de \zh{遗产} (\emph{yíchǎn}, héritage).
\end{definition}

Il faut distinguer deux grandes catégories d'héritage :

\begin{itemize}
\item L'\cle{héritage matériel} : ce qu'on peut toucher, mesurer, partager — la maison \zh{房子} (\emph{fángzi}), les champs et les terres \zh{土地} (\emph{tǔdì}), l'argent \zh{钱} (\emph{qián}), le bétail, les meubles.
\item L'\cle{héritage immatériel} : ce qu'on ne peut pas peser mais qui a une grande valeur — le nom de famille \zh{姓} (\emph{xìng}), le titre de chef \zh{头衔} (\emph{tóuxián}), les traditions \zh{传统} (\emph{chuántǒng}), les savoir-faire, les histoires familiales.
\end{itemize}

Au Cameroun, dans beaucoup de communautés, le titre de chef ou de notable se transmet selon des règles coutumières précises, parfois au fils aîné, parfois à un neveu choisi par la famille. En Chine, le nom de famille \zh{姓} (\emph{xìng}) se transmet traditionnellement par le père, et le respect des ancêtres fait partie de l'héritage immatériel le plus important.

\begin{definition}[继承 et 继承人]
\zh{继承} (\emph{jìchéng}) est un verbe : \cle{hériter}, recevoir en héritage, mais aussi « poursuivre » une tradition. Structure : \textbf{Sujet + 继承 + complément}. Exemple : \zh{儿子继承父亲的土地。} (\emph{Érzi jìchéng fùqin de tǔdì.}) « Le fils hérite des terres du père. »\\
\zh{继承人} (\emph{jìchéngrén}) est un nom : l'\cle{héritier}, la personne qui reçoit l'héritage. Exemple : \zh{她是唯一的继承人。} (\emph{Tā shì wéiyī de jìchéngrén.}) « Elle est la seule héritière. »
\end{definition}

\begin{definition}[遗嘱 yízhǔ : le testament]
\zh{遗嘱} (\emph{yízhǔ}) désigne le \cle{testament} : le document ou la déclaration par lesquels une personne dit, avant sa mort, comment ses biens doivent être partagés. Exemple : \zh{爷爷写了遗嘱。} (\emph{Yéye xiě le yízhǔ.}) « Grand-père a écrit un testament. » Dans beaucoup de familles camerounaises, le partage se fait oralement devant la famille élargie, sans testament écrit : c'est une différence de pratique intéressante à observer.
\end{definition}

\begin{savaistu}[Le caractère 遗 : ce qui « passe »]
En chinois, le caractère \zh{遗} (\emph{yí}, laisser après soi) contient la clé \zh{辶} (\emph{chuò}), la clé du \textbf{mouvement}, qu'on retrouve dans \zh{过} (\emph{guò}, passer), \zh{进} (\emph{jìn}, entrer), \zh{送} (\emph{sòng}, offrir, accompagner). L'héritage, c'est littéralement ce qui « se déplace », ce qui « passe » d'une génération à l'autre. La partie droite \zh{贵} (\emph{guì}, précieux) rappelle que ce qui est transmis a de la valeur. Belle image pour mémoriser le caractère : un trésor en mouvement !
\end{savaistu}

\courssub{Texte support : deux familles, deux traditions}

Lisons maintenant un texte qui met en scène une famille camerounaise et une famille chinoise. Repérez les mots-clés du thème : \zh{遗产}, \zh{继承}, \zh{继承人}, \zh{遗嘱}.

\begin{textebox}[两个家庭的故事 — L'histoire de deux familles]
在喀麦隆的雅温得，有一个大家庭。爷爷去世以后，全家人聚在一起开会。爷爷没有写遗嘱，所以家里人要一起决定：谁继承房子和土地？按照传统，大儿子继承土地，但是女儿们也能分到一些东西。妈妈说：“遗产是全家人的，不是一个人的。”\\
在中国上海，王先生的父亲也去世了。父亲留下了遗嘱：房子给儿子和女儿，钱给妻子。王先生说：“我们尊重父亲的遗嘱。”他的妹妹也是继承人，她和哥哥一样继承父亲的遗产。\\
两个国家，两个家庭，但是问题是一样的：人去世以后，谁继承什么？传统和法律，有时候一样，有时候不一样。\\
\emph{Zài Kāmàilóng de Yǎwēndé, yǒu yí ge dà jiātíng. Yéye qùshì yǐhòu, quánjiā rén jù zài yìqǐ kāihuì. Yéye méiyǒu xiě yízhǔ, suǒyǐ jiālǐ rén yào yìqǐ juédìng : shéi jìchéng fángzi hé tǔdì ? Ànzhào chuántǒng, dà érzi jìchéng tǔdì, dànshì nǚ'érmen yě néng fēndào yìxiē dōngxi. Māma shuō : “Yíchǎn shì quánjiā rén de, bú shì yí ge rén de.”}\\\emph{Zài Zhōngguó Shànghǎi, Wáng xiānsheng de fùqin yě qùshì le. Fùqin liúxià le yízhǔ : fángzi gěi érzi hé nǚ'ér, qián gěi qīzi. Wáng xiānsheng shuō : “Wǒmen zūnzhòng fùqin de yízhǔ.” Tā de mèimei yě shì jìchéngrén, tā hé gēge yíyàng jìchéng fùqin de yíchǎn.}\\\emph{Liǎng ge guójiā, liǎng ge jiātíng, dànshì wèntí shì yíyàng de : rén qùshì yǐhòu, shéi jìchéng shénme ? Chuántǒng hé fǎlǜ, yǒu shíhou yíyàng, yǒu shíhou bù yíyàng.}\\
« À Yaoundé, au Cameroun, il y a une grande famille. Après le décès du grand-père, toute la famille se réunit en conseil. Le grand-père n'a pas écrit de testament, donc la famille doit décider ensemble : qui hérite de la maison et des terres ? Selon la tradition, le fils aîné hérite des terres, mais les filles peuvent aussi recevoir une part. La maman dit : “L'héritage appartient à toute la famille, pas à une seule personne.”\\
À Shanghai, en Chine, le père de M. Wang est décédé lui aussi. Le père a laissé un testament : la maison pour le fils et la fille, l'argent pour son épouse. M. Wang dit : “Nous respectons le testament de notre père.” Sa petite sœur est aussi héritière : comme son frère, elle hérite des biens du père.\\
Deux pays, deux familles, mais la question est la même : après le décès d'une personne, qui hérite de quoi ? Tradition et loi coïncident parfois, parfois non. »
\end{textebox}

\textbf{Glossaire du texte}

\begin{center}
\begin{tabularx}{\linewidth}{|X|X|X|X|}
\hline
\rowcolor{popblueL} Caractères & Pinyin & Nature & Traduction \\
\hline
遗产 & yíchǎn & nom & héritage, biens laissés \\
\hline
继承 & jìchéng & verbe & hériter, recevoir en héritage \\
\hline
继承人 & jìchéngrén & nom & héritier, héritière \\
\hline
遗嘱 & yízhǔ & nom & testament \\
\hline
去世 & qùshì & verbe & décéder (terme respectueux) \\
\hline
留下 & liúxià & verbe & laisser (après soi) \\
\hline
土地 & tǔdì & nom & terre, terrain \\
\hline
按照 & ànzhào & préposition & selon, conformément à \\
\hline
传统 & chuántǒng & nom & tradition \\
\hline
尊重 & zūnzhòng & verbe & respecter \\
\hline
法律 & fǎlǜ & nom & loi, législation \\
\hline
分 & fēn & verbe & partager, diviser \\
\hline
\end{tabularx}
\end{center}

\textbf{Questions de compréhension}

\begin{enumerate}
\item \zh{爷爷有没有遗嘱？} (\emph{Yéye yǒu méiyǒu yízhǔ ?}) Le grand-père de Yaoundé a-t-il laissé un testament ?
\item \zh{按照传统，谁继承土地？} Selon la tradition camerounaise décrite, qui hérite des terres ?
\item \zh{王先生的妹妹是不是继承人？} La petite sœur de M. Wang est-elle héritière ?
\item Quelle phrase de la maman camerounaise résume une idée importante sur l'héritage ? Traduisez-la.
\end{enumerate}

\courssub{Carte mentale du vocabulaire : autour de 遗产}

\begin{popfigure}
\begin{center}
\begin{tikzpicture}[mindmap, grow cyclic, every node/.style={concept, align=center}, concept color=popblue!40, level 1/.append style={level distance=3.2cm, sibling angle=72}, text=popdark]
\node[align=center]{\zh{遗产}\\ \emph{yíchǎn}\\ héritage}
  child [concept color=poporange!50] { node[align=center]{\zh{土地}\\ \emph{tǔdì}\\ terres} }
  child [concept color=popgreen!50] { node[align=center]{\zh{房子}\\ \emph{fángzi}\\ maison} }
  child [concept color=popgold!60] { node[align=center]{\zh{钱}\\ \emph{qián}\\ argent} }
  child [concept color=poppurple!40] { node[align=center]{\zh{头衔}\\ \emph{tóuxián}\\ titre de chef} }
  child [concept color=poppink!50] { node[align=center]{\zh{传统}\\ \emph{chuántǒng}\\ traditions} };
\end{tikzpicture}
\end{center}
\legende{Carte mentale : les cinq grandes branches de l'héritage — trois matérielles (terres, maison, argent), deux immatérielles (titre, traditions).}
\end{popfigure}

Observez la carte : les trois premières branches (\zh{土地}, \zh{房子}, \zh{钱}) forment l'héritage \textbf{matériel} ; les deux dernières (\zh{头衔}, \zh{传统}) forment l'héritage \textbf{immatériel}. Retenez ce classement, il servira dans tout le module.

\courssub{Structures utiles pour parler d'héritage}

\begin{propriete}[Dire « qui hérite de quoi »]
Structure de base : \textbf{Sujet (héritier) + 继承 + complément (bien hérité)}.\\
\zh{大儿子继承土地。} (\emph{Dà érzi jìchéng tǔdì.}) « Le fils aîné hérite des terres. »\\
Pour dire « hériter de quelqu'un », le chinois utilise souvent le possessif avec \zh{的} : \zh{继承父亲的遗产} « hériter des biens du père ». On ne dit pas une préposition « de » comme en français : le lien se fait par \zh{的} entre le possesseur et le bien.
\end{propriete}

\begin{propriete}[Donner à quelqu'un : 给]
Pour décrire le partage d'un testament : \textbf{Donneur + 给 + bénéficiaire + bien} ou \textbf{bien + 给 + bénéficiaire}.\\
\zh{房子给儿子和女儿。} (\emph{Fángzi gěi érzi hé nǚ'ér.}) « La maison (revient) au fils et à la fille. »\\
\zh{父亲把钱给妻子。} (\emph{Fùqin bǎ qián gěi qīzi.}) « Le père donne l'argent à son épouse. »
\end{propriete}

\begin{center}
\begin{tabularx}{\linewidth}{|X|X|X|}
\hline
\rowcolor{popblueL} Structure & Exemple en chinois + pinyin & Traduction \\
\hline
Sujet + 继承 + bien & \zh{她继承母亲的房子。} \emph{Tā jìchéng mǔqin de fángzi.} & Elle hérite de la maison de sa mère. \\
\hline
Sujet + 是 + 继承人 & \zh{他是继承人。} \emph{Tā shì jìchéngrén.} & Il est l'héritier. \\
\hline
Sujet + 留下 + bien & \zh{爷爷留下了土地。} \emph{Yéye liúxià le tǔdì.} & Grand-père a laissé des terres. \\
\hline
Bien + 给 + bénéficiaire & \zh{钱给妻子。} \emph{Qián gěi qīzi.} & L'argent (revient) à l'épouse. \\
\hline
按照 + règle + , + phrase & \zh{按照传统，儿子继承土地。} \emph{Ànzhào chuántǒng, érzi jìchéng tǔdì.} & Selon la tradition, le fils hérite des terres. \\
\hline
\end{tabularx}
\end{center}

\begin{attention}[Ne pas confondre 遗产 et 传统]
\zh{遗产} (\emph{yíchǎn}) désigne l'\textbf{héritage}, les biens laissés après un décès ; \zh{传统} (\emph{chuántǒng}) désigne la \textbf{tradition}, les coutumes transmises de génération en génération. Les deux peuvent « se transmettre », mais pas avec les mêmes verbes : on dit \zh{继承遗产} (hériter de biens) et plutôt \zh{继承传统} dans le sens de « perpétuer une tradition », mais on ne dit jamais \zh{遗产} pour parler d'une coutume. Erreur fréquente des francophones : traduire « héritage culturel » par \zh{遗产} seul — le chinois précise \zh{文化遗产} (\emph{wénhuà yíchǎn}, patrimoine culturel).
\end{attention}

\begin{attention}[La place de 以后]
En français, « après » précède son complément (« après le décès »). En chinois, \zh{以后} se place \textbf{après} : \zh{去世以后} « décès + après ». Ne dites pas \zh{以后去世} pour « après le décès » : cela signifierait « après, décéder ».
\end{attention}

\courssub{Exercice résolu et entraînement}

\begin{exoresolu}[Traduction guidée]
Énoncé : traduisez en chinois la phrase suivante : « Après le décès du père, la fille aînée hérite de la maison, et le fils hérite des terres. »
\tcblower
Solution détaillée, étape par étape :
\begin{enumerate}
\item « Après le décès du père » : le complément de temps avec \zh{以后} se place en tête de phrase, et \zh{以后} suit son complément : \zh{父亲去世以后}.
\item « la fille aînée » : \zh{大女儿} (\emph{dà nǚ'ér}) ; « hérite de la maison » : \zh{继承房子}.
\item « et » entre deux phrases : \zh{而} (\emph{ér}) ou simplement la virgule chinoise ， ; ici on peut utiliser \zh{儿子继承土地} après la virgule.
\end{enumerate}
Phrase complète : \zh{父亲去世以后，大女儿继承房子，儿子继承土地。}\\
\emph{Fùqin qùshì yǐhòu, dà nǚ'ér jìchéng fángzi, érzi jìchéng tǔdì.}\\
Points de vigilance : bien placer \zh{以后} après \zh{去世} ; ne pas ajouter de préposition « de » entre \zh{继承} et le bien ; la virgule chinoise ，sépare les deux propositions.
\end{exoresolu}

\begin{exercice}[À vous de jouer]
\begin{enumerate}
\item Traduisez en français : \zh{妈妈留下了一个遗嘱。} (\emph{Māma liúxià le yí ge yízhǔ.})
\item Traduisez en chinois : « Selon la loi, tous les enfants sont héritiers. » (utilisez \zh{按照}, \zh{法律}, \zh{都}, \zh{继承人})
\item Classez ces mots en deux colonnes « héritage matériel / immatériel » : \zh{钱}, \zh{头衔}, \zh{房子}, \zh{传统}, \zh{土地}.
\item Répondez en chinois avec une phrase complète : \zh{在你的家庭里，谁继承土地？} (\emph{Zài nǐ de jiātíng lǐ, shéi jìchéng tǔdì ?})
\end{enumerate}
\end{exercice}

\begin{corrige}[Exercice « À vous de jouer »]
\begin{enumerate}
\item « Maman a laissé un testament. »
\item \zh{按照法律，所有的孩子都是继承人。} (\emph{Ànzhào fǎlǜ, suǒyǒu de háizi dōu shì jìchéngrén.})
\item Matériel : \zh{钱}, \zh{房子}, \zh{土地} ; immatériel : \zh{头衔}, \zh{传统}.
\item Réponse possible : \zh{在我的家庭里，大儿子继承土地。} (\emph{Zài wǒ de jiātíng lǐ, dà érzi jìchéng tǔdì.}) « Dans ma famille, le fils aîné hérite des terres. » Toute phrase complète avec \zh{继承} est acceptée.
\end{enumerate}
\end{corrige}

\begin{aretenir}[L'essentiel de l'introduction]
\begin{itemize}
\item L'hérédité = transmission des biens, terres, titres et statut après un décès : \zh{遗产} (héritage), \zh{继承} (hériter), \zh{继承人} (héritier), \zh{遗嘱} (testament).
\item On distingue l'héritage \textbf{matériel} (\zh{房子}, \zh{土地}, \zh{钱}) et l'héritage \textbf{immatériel} (\zh{头衔}, \zh{传统}, nom de famille).
\item Structure clé : \textbf{Sujet + 继承 + bien} ; le lien « de quelqu'un » se fait avec \zh{的}.
\item Piège : \zh{以后} se place après son complément ; \zh{遗产} n'est pas \zh{传统}.
\end{itemize}
\end{aretenir}

\courssub{L'hérédité au Cameroun : entre coutume et loi moderne}

\courssub{Texte support : l'héritage au Cameroun}

\begin{textebox}[喀麦隆的继承习俗与法律 — Héritage au Cameroun : coutumes et loi]
喀麦隆有两百多个民族，继承习俗不一样。在很多草原地区（如巴米莱克、巴穆恩），传统上实行\zh{长子继承制} (\emph{zhǎngzǐ jìchéngzhì})：大儿子继承父亲的土地、房子和头衔。女儿传统上不继承土地，但可以得到钱或物品。寡妇 (\emph{guāfù}, veuve) 过去常被家人“继承”（\zh{继娶} \emph{jìqǔ}, lévirat), 现在法律禁止。\\
现代法律 (\emph{现代法律} \emph{xiàndài fǎlǜ}, Code des personnes et de la famille, 2011) 规定：男女平等 (\emph{男女平等} \emph{nánnǚ píngděng}), 所有子女 (\emph{所有子女} \emph{suǒyǒu zǐnǚ}) 都是法定继承人 (\emph{法定继承人} \emph{fǎdìng jìchéngrén}). 但是在农村，习惯法 (\emph{习惯法} \emph{xíguànfǎ}) 依然很强。家庭会议 (\emph{家庭会议} \emph{jiātíng huìyì}) 往往按传统分配。\\
\emph{Kāmàilóng yǒu liǎngbǎi duō ge mínzú, jìchéng xísú bù yíyàng. Zài hěnduō cǎoyuán dìqū (rú Bāmǐlèikè, Bāmù'ēn), chuántǒng shàng shíxíng \zh{长子继承制} : dà érzi jìchéng fùqin de tǔdì, fángzi hé tóuxián. Nǚ'ér chuántǒng shàng bù jìchéng tǔdì, dàn kěyǐ dédào qián huò wùpǐn. Guāfù guòqù cháng bèi jiārén “jìchéng” (\zh{继娶} \emph{jìqǔ}), xiànzài fǎlǜ jìnzhǐ.}\\\emph{Xiàndài fǎlǜ (Code des personnes et de la famille, 2011) guīdìng : nánnǚ píngděng, suǒyǒu zǐnǚ dōu shì fǎdìng jìchéngrén. Dànshì zài nóngcūn, xíguànfǎ yīrán hěn qiáng. Jiātíng huìyì wǎngwǎng ànzhào chuántǒng fēnpèi.}\\
« Le Cameroun compte plus de 200 ethnies, les coutumes successorales varient. Dans de nombreuses régions des Grassfields (Bamiléké, Bamoun), la tradition applique la \cle{primogéniture} (\zh{长子继承制}) : le fils aîné hérite des terres, de la maison et du titre du père. Les filles traditionnellement n'héritent pas des terres, mais peuvent recevoir de l'argent ou des objets. La veuve (\zh{寡妇} \emph{guāfù}) était autrefois souvent "héritée" par la famille (\zh{继娶} \emph{jìqǔ}, lévirat), la loi l'interdit aujourd'hui.\\
La loi moderne (Code des personnes et de la famille, 2011) stipule : l'égalité homme-femme (\zh{男女平等}), tous les enfants (\zh{所有子女}) sont héritiers légaux (\zh{法定继承人}). Mais en milieu rural, le droit coutumier (\zh{习惯法}) reste fort. Le conseil de famille (\zh{家庭会议}) partage souvent selon la tradition. »
\end{textebox}

\textbf{Vocabulaire spécifique : Cameroun}

\begin{center}
\begin{tabularx}{\linewidth}{|X|X|X|X|}
\hline
\rowcolor{popblueL} Caractères & Pinyin & Nature & Traduction \\
\hline
长子继承制 & zhǎngzǐ jìchéngzhì & nom & primogéniture (droit d'aînesse) \\
\hline
寡妇 & guāfù & nom & veuve \\
\hline
继娶 & jìqǔ & verbe & lévirat (épouser la veuve du frère/défunt) \\
\hline
习惯法 & xíguànfǎ & nom & droit coutumier \\
\hline
法定继承人 & fǎdìng jìchéngrén & nom & héritier légal / légitime \\
\hline
男女平等 & nánnǚ píngděng & locution & égalité homme-femme \\
\hline
家庭会议 & jiātíng huìyì & nom & conseil de famille \\
\hline
分配 & fēnpèi & verbe & répartir, distribuer (parts) \\
\hline
\end{tabularx}
\end{center}

\begin{attention}[Le lévirat : 继娶 vs 继承]
Attention : \zh{继娶} (\emph{jìqǔ}) ne signifie pas « hériter » au sens de biens (\zh{继承}), mais « épouser la veuve d'un parent décédé » (lévirat). C'est une pratique coutumière ancienne, aujourd'hui interdite par la loi camerounaise (Code de la famille, art. 73) et chinoise. Ne confondez pas les caractères : \zh{承} (recevoir, porter) vs \zh{娶} (prendre pour femme, clé \zh{女}).
\end{attention}

\courssub{L'hérédité en Chine : de la tradition patriarcale au Code civil}

\courssub{Texte support : l'héritage en Chine}

\begin{textebox}[中国的继承传统与民法典 — Chine : tradition et Code civil]
中国传统上实行\zh{宗法制} (\emph{zōngfǎzhì}, système patriarcal de lignage). 核心是\zh{长子继承制} (\emph{zhǎngzǐ jìchéngzhì})：嫡长子 (\emph{dí zhǎngzǐ}, fils aîné de l'épouse principale) 继承家产、家族头衔和\zh{祭祖} (\emph{jìzǔ}, culte des ancêtres) 的责任。女儿出嫁后 (\emph{出嫁} \emph{chūjià}), 一般不分家产，被视为“泼出去的水” (\emph{pō chūqù de shuǐ}, eau versée hors du foyer).\\
1985年\zh{继承法} 颁布，确立男女平等继承权。2021年\zh{民法典} (\emph{mínfǎdiǎn}, Code civil) 取代继承法，细化遗嘱形式：\zh{公证遗嘱} (\emph{gōngzhèng yízhǔ}, testament notarié), \zh{自书遗嘱} (\emph{zìshū yízhǔ}, testament olographe), \zh{代书遗嘱} (\emph{dàishū yízhǔ}, testament par tiers), \zh{录音录像遗嘱} (\emph{lùyīn lùxiàng yízhǔ}, testament audio/vidéo), \zh{口头遗嘱} (\emph{kǒutóu yízhǔ}, testament oral - urgence). 遗嘱自由 (\emph{遗嘱自由} \emph{yízhǔ zìyóu}) 受尊重，但必须保留必要份额给无劳动能力继承人 (\emph{无劳动能力继承人} \emph{wú láodòng nénglì jìchéngrén}).\\
\emph{Zhōngguó chuántǒng shàng shíxíng \zh{宗法制}. Héxīn shì \zh{长子继承制} : dí zhǎngzǐ jìchéng jiāchǎn, jiāzú tóuxián hé \zh{祭祖} de zérèn. Nǚ'ér chūjià hòu, yībān bù fēn jiāchǎn, bèi shìwéi “pō chūqù de shuǐ”.}\\\emph{Yījiǔbāwǔ nián \zh{继承法} bānbu, quànlì nánnǚ píngděng jìchéngquán. Èrlíngèryī nián \zh{民法典} tìdài \zh{继承法}, xìhuà yízhǔ xíngshì : \zh{公证遗嘱}, \zh{自书遗嘱}, \zh{代书遗嘱}, \zh{录音录像遗嘱}, \zh{口头遗嘱}. Yízhǔ zìyóu shòu zūnzhòng, dàn bìxū bǎoliú bìyào fèn'é gěi wú láodòng nénglì jìchéngrén.}\\
« La Chine appliquait traditionnellement le \cle{système patriarcal de lignage} (\zh{宗法制}). Le cœur est la \cle{primogéniture} (\zh{长子继承制}) : le fils aîné légitime (\zh{嫡长子}) hérite du patrimoine familial, du titre de chef de lignage et de la responsabilité du \cle{culte des ancêtres} (\zh{祭祖}). La fille, une fois mariée (\zh{出嫁}), ne partage généralement pas le patrimoine, considérée comme de « l'eau versée hors du foyer » (\zh{泼出去的水}).\\
La \cle{Loi sur les successions} de 1985 établit l'égalité des droits successoraux hommes-femmes. Le \cle{Code civil} de 2021 la remplace, détaillant les formes de testament : \cle{testament notarié} (\zh{公证遗嘱}), \cle{testament olographe} (\zh{自书遗嘱}), \cle{testament par tiers} (\zh{代书遗嘱}), \cle{testament audio/vidéo} (\zh{录音录像遗嘱}), \cle{testament oral} (\zh{口头遗嘱} - urgence). La \cle{liberté testamentaire} (\zh{遗嘱自由}) est respectée, mais une part réservée doit être laissée aux héritiers sans capacité de travail (\zh{无劳动能力继承人}). »
\end{textebox}

\textbf{Vocabulaire spécifique : Chine}

\begin{center}
\begin{tabularx}{\linewidth}{|X|X|X|X|}
\hline
\rowcolor{popblueL} Caractères & Pinyin & Nature & Traduction \\
\hline
宗法制 & zōngfǎzhì & nom & système patriarcal de lignage \\
\hline
嫡长子 & dí zhǎngzǐ & nom & fils aîné légitime (épouse principale) \\
\hline
祭祖 & jìzǔ & verbe/nom & honorer les ancêtres / culte ancestral \\
\hline
出嫁 & chūjià & verbe & se marier (pour une fille, quitter le foyer natal) \\
\hline
民法典 & mínfǎdiǎn & nom & Code civil (2021) \\
\hline
公证遗嘱 & gōngzhèng yízhǔ & nom & testament notarié \\
\hline
自书遗嘱 & zìshū yízhǔ & nom & testament olographe (écrit main) \\
\hline
无劳动能力 & wú láodòng nénglì & adj. & sans capacité de travail (invalide, âgé) \\
\hline
\end{tabularx}
\end{center}

\begin{savaistu}[Le caractère 宗 : l'ancêtre et le toit]
\zh{宗} (\emph{zōng}) signifie « ancêtre », « lignage principal », « respecter ». En haut : \zh{宀} (\emph{mián}, toit, maison) ; en bas : \zh{示} (\emph{shì}, autel, montrer, rite). L'ancêtre est celui qu'on honore sous le toit familial par des rites. \zh{宗法制} (\emph{zōngfǎzhì}) est le système juridique et social organisé autour du lignage agnatique (par les hommes).
\end{savaistu}

\courssub{Comparaison Cameroun / Chine : faits et regards croisés}

Le tableau ci-dessous synthétise les points clés pour une analyse factuelle et respectueuse.

\begin{center}
\begin{tabularx}{\linewidth}{|X|X|X|}
\hline
\rowcolor{popblueL} Critère & Cameroun (aperçu général) & Chine (aperçu général) \\
\hline
\textbf{Fondement traditionnel} & Droit coutumier (\zh{习惯法}) variable selon ethnies (patrilinéaire ou matrilinéaire). Primogéniture fréquente (Grassfields). & Système patriarcal de lignage (\zh{宗法制}). Primogéniture stricte (\zh{长子继承制}) : fils aîné légitime. \\
\hline
\textbf{Héritiers traditionnels} & Fils aîné (terres, titres). Filles : parts mobilières. Veuve : lévirat (\zh{继娶}) ou entretien par la famille. & Fils aîné légitime (patrimoine + culte ancêtres \zh{祭祖}). Autres fils : parts mineures. Filles : exclues (mariées \zh{出嫁}). Veuve : reste dans la famille, entretient autel. \\
\hline
\textbf{Loi moderne (égalité)} & Code des personnes et de la famille (2011) : égalité totale hommes/femmes, tous enfants héritiers légaux (\zh{法定继承人}). & Code civil (2021, ex-Loi 1985) : égalité totale, liberté testamentaire (\zh{遗嘱自由}), part réservée aux vulnérables. \\
\hline
\textbf{Testament} & Rare en milieu coutumier. Oral devant conseil de famille (\zh{家庭会议}). Écrit admis par la loi. & Très codifié (5 formes légales depuis 2021 : notarié \zh{公证}, olographe \zh{自书}, etc.). Priorité au testament notarié (sauf dernier). \\
\hline
\textbf{Héritage immatériel} & Titres de chefferie, noms de famille, masques, danses, secrets de lignage. Transmission souvent au neveu (sœur du défunt) ou fils aîné. & Nom de clan (\zh{姓}), registre généalogique (\zh{族谱}), culte ancêtres (\zh{祭祖}), préceptes familiaux (\zh{家训}). Transmission patrilinéaire. \\
\hline
\textbf{Pratique actuelle} & Pluralisme juridique : loi officielle vs coutume vivante. Conflits fréquents (terres urbaines). Tendance à l'écrit. & Uniformisation par le Code civil. Urbanisation : biens immobiliers (appartements) = enjeu principal. Testaments vidéo en hausse. \\
\hline
\end{tabularx}
\end{center}

\begin{aretenir}[Points de convergence et divergence]
\begin{itemize}
\item \textbf{Convergence :} Les deux pays ont connu un système traditionnel favorisant les fils aînés (\zh{长子继承制}). Les deux ont adopté des lois modernes garantissant l'\textbf{égalité homme-femme} (\zh{男女平等}) et les droits de \textbf{tous les enfants} (\zh{所有子女}).
\item \textbf{Divergence :} Au Cameroun, le \textbf{pluralisme juridique} (loi + coutume) crée des situations hybrides ; le conseil de famille (\zh{家庭会议}) reste central. En Chine, le \textbf{Code civil unifié} (2021) prime, avec une codification très technique des testaments.
\item \textbf{Immatériel :} Au Cameroun, transmission souvent \textbf{matrilinéaire} pour les titres (neveu = fils de la sœur) dans certains groupes ; en Chine, strictement \textbf{patrilinéaire} (fils aîné, registre \zh{族谱}).
\end{itemize}
\end{aretenir}

\courssub{Vocabulaire approfondi et écriture des caractères}

\courssub{Tableau récapitulatif du vocabulaire de la leçon (HSK 3-4)}

\begin{center}
\begin{tabularx}{\linewidth}{|X|X|X|X|}
\hline
\rowcolor{popblueL} Caractères & Pinyin & Nature & Traduction \\
\hline
遗产 & yíchǎn & nom & héritage, succession \\
\hline
继承 & jìchéng & verbe & hériter, succéder, perpétuer \\
\hline
继承人 & jìchéngrén & nom & héritier, héritière \\
\hline
遗嘱 & yízhǔ & nom & testament \\
\hline
去世 & qùshì & verbe & décéder (euphémisme respectueux) \\
\hline
留下 & liúxià & verbe & laisser, léguer \\
\hline
土地 & tǔdì & nom & terre, terrain, foncier \\
\hline
房子 & fángzi & nom & maison, bâtiment \\
\hline
头衔 & tóuxián & nom & titre, fonction, qualité \\
\hline
传统 & chuántǒng & nom & tradition, coutume \\
\hline
法律 & fǎlǜ & nom & loi, droit, législation \\
\hline
按照 & ànzhào & prép. & selon, conformément à, d'après \\
\hline
分配 & fēnpèi & verbe & répartir, attribuer des parts \\
\hline
长子继承制 & zhǎngzǐ jìchéngzhì & nom & primogéniture, droit d'aînesse \\
\hline
法定继承人 & fǎdìng jìchéngrén & nom & héritier légal / réservataire \\
\hline
男女平等 & nánnǚ píngděng & loc. & égalité entre hommes et femmes \\
\hline
习惯法 & xíguànfǎ & nom & droit coutumier \\
\hline
宗法制 & zōngfǎzhì & nom & système patriarcal de lignage \\
\hline
祭祖 & jìzǔ & verbe/nom & vénérer les ancêtres / culte ancestral \\
\hline
民法典 & mínfǎdiǎn & nom & Code civil \\
\hline
公证遗嘱 & gōngzhèng yízhǔ & nom & testament notarié \\
\hline
自书遗嘱 & zìshū yízhǔ & nom & testament olographe \\
\hline
\end{tabularx}
\end{center}

\courssub{Écriture des caractères : ordre des traits et clés (radicaux)}

Maîtriser l'écriture des caractères clés du thème. Respectez l'ordre : \textbf{haut → bas, gauche → droite, horizontal avant vertical, trait central avant les ailes, cadre avant contenu, fermeture en dernier}.

\begin{exemplebox}[Analyse graphique de 5 caractères clés]
\begin{enumerate}
\item \textbf{遗  (yí)} — 13 traits. Clé : \zh{辶} (\emph{chuò}, marche/passage, 3 traits). Partie phonétique/sémantique : \zh{贵} (\emph{guì}, précieux). Ordre : écrire \zh{贵} (haut \zh{竹} simplifié, bas \zh{贝}), puis la clé \zh{辶} (point, horizontal court, grand trait sinueux).
\item \textbf{产  (chǎn)} — 6 traits. Clé : \zh{生} (\emph{shēng}, naître/vie, 5 traits) — *note : clé souvent classée sous \zh{立} ou \zh{生}*. Structure : \zh{立} (debout) + \zh{二} + \zh{厂} (cliff). Ordre : \zh{立} (horizontal, vertical, horizontal, horizontal), puis \zh{二} (deux horizontaux), puis \zh{厂} (horizontal, vertical crochet).
\item \textbf{继  (jì)} — 10 traits. Clé : \zh{纟} (\emph{sī}, fil de soie, 3 traits). Partie droite : \zh{米} (riz) + \zh{儿} (enfant). Ordre : clé \zh{纟} (diagonale, diagonale, horizontal), puis \zh{米} (deux diagonales, horizontal, deux traits points), puis \zh{儿} (horizontal, crochet vertical descendant).
\item \textbf{承  (chéng)} — 8 traits. Clé : \zh{手} (\emph{shǒu}, main, ici \zh{扌} 3 traits). Partie phonétique : \zh{丙} (\emph{bǐng}). Ordre : clé \zh{扌} (horizontal, crochet vertical, point), puis \zh{丙} (horizontal, vertical, horizontal, horizontal, point final long).
\item \textbf{嘱  (zhǔ)} — 15 traits. Clé : \zh{口} (\emph{kǒu}, bouche, 3 traits). Partie phonétique : \zh{主} (\emph{zhǔ}, maître) + \zh{寸} (\emph{cùn}, pouce/mesure). Ordre : \zh{口} (vertical, horizontal+crochet, horizontal), puis \zh{主} (point, horizontal, vertical, horizontal), puis \zh{寸} (horizontal, point, horizontal, vertical crochet).
\end{enumerate}
\end{exemplebox}

\begin{exercice}[Entraînement à l'écriture]
Écrivez chaque caractère ci-dessus 5 fois en respectant l'ordre des traits. Pour chacun, identifiez la clé (radical) et donnez un mot composé utilisant ce caractère (autre que ceux vus en classe).
\end{exercice}

\courssub{Réflexion, débat et mini-exposé}

\courssub{Questions de réflexion (préparation orale / écrite)}

\begin{enumerate}
\item \textbf{Analyse comparative :} « La loi moderne a tout changé. » Êtes-vous d'accord ? Nuancez en citant le rôle du \zh{家庭会议} au Cameroun et la \zh{遗嘱自由} en Chine.
\item \textbf{Héritage immatériel :} Pourquoi le \zh{祭祖} (culte des ancêtres) est-il indissociable de l'héritage en Chine traditionnelle ? Quel équivalent trouvez-vous au Cameroun (ex: \zh{头衔} de chef, masques, \zh{族谱} vs généalogie orale) ?
\item \textbf{Femmes et héritage :} Comparez l'évolution du statut de la \zh{寡妇} (veuve) et de la \zh{女儿} (fille) dans les deux pays. Citez les articles de loi si possible (Code camerounais 2011, Code civil chinois 2021).
\item \textbf{Conflits :} Comment la rédaction d'un \zh{遗嘱} (testament) peut-elle prévenir ou au contraire déclencher des conflits familiaux ? Donnez un exemple fictif plausible.
\end{enumerate}

\begin{experience}[Mini-exposé oral : « Héritage et modernité »]
\textbf{Consigne :} En binôme, préparez un exposé de 3 minutes (1 min 30 par élève) en chinois, niveau HSK 3-4. Vous présenterez \textbf{un} des deux sujets au choix :
\begin{enumerate}
\item \textbf{Sujet A :} \zh{我的家庭继承故事} (L'histoire d'héritage de ma famille — réel ou inventé). Utilisez : \zh{去世}, \zh{留下}, \zh{遗产}, \zh{继承}, \zh{按照}, \zh{传统}, \zh{法律}, \zh{公平} (\emph{gōngpíng}, équitable).
\item \textbf{Sujet B :} \zh{喀麦隆与中国：继承法的比较} (Cameroun vs Chine : comparaison des droits successoraux). Utilisez : \zh{相同点} (\emph{xiāngtóngdiǎn}, points communs), \zh{不同点} (\emph{bùtóngdiǎn}, différences), \zh{男女平等}, \zh{习惯法}, \zh{民法典}, \zh{遗嘱}.
\end{enumerate}
\textbf{Critères d'évaluation :} Prononciation (tons), exactitude grammaticale (structures vues), richesse lexicale (vocabulaire du tableau), cohérence du propos, contact visuel.
\textbf{Aide au démarrage (amorces) :}
\begin{itemize}
\item \zh{大家好，今天我想谈谈……} (Bonjour, aujourd'hui je veux parler de…)
\item \zh{首先，在喀麦隆…… 其次，在中国……} (Premièrement, au Cameroun… Deuxièmement, en Chine…)
\item \zh{我认为…… 因为……} (Je pense que… parce que…)
\item \zh{谢谢大家。} (Merci à tous.)
\end{itemize}
\end{experience}

\begin{methode}[Réussir son mini-exposé socioculturel en chinois]
\begin{enumerate}
\item \textbf{Structurer :} Introduction (sujet) → Développement (2-3 points avec connecteurs \zh{首先}、\zh{其次}、\zh{最后}) → Conclusion (avis personnel \zh{我认为}).
\item \textbf{Vocabulaire :} Sélectionnez 8-10 mots du tableau de la leçon et forcez-vous à les placer.
\item \textbf{Phrases-types :} Préparez 3-4 phrases « passe-partout » (\zh{这是一个很重要的问题}, \zh{情况比较复杂}, \zh{随着法律的发展……}).
\item \textbf{Entraînement :} Enregistrez-vous, écoutez, corrigez les tons (surtout \zh{继承} \emph{jìchéng} vs \zh{经济} \emph{jīngjì}, \zh{土地} \emph{tǔdì} vs \zh{土豆} \emph{tǔdòu}).
\end{enumerate}
\end{methode}

\begin{aretenir}[Synthèse finale de la leçon ZH04]
\begin{itemize}
\item \textbf{Vocabulaire clé :} \zh{遗产}, \zh{继承}, \zh{继承人}, \zh{遗嘱}, \zh{去世}, \zh{留下}, \zh{土地}, \zh{房子}, \zh{头衔}, \zh{传统}, \zh{法律}, \zh{按照}, \zh{分配}, \zh{长子继承制}, \zh{法定继承人}, \zh{男女平等}, \zh{习惯法}, \zh{宗法制}, \zh{祭祖}, \zh{民法典}.
\item \textbf{Grammaire :} \zh{以后} (après) placé \textbf{après} le verbe/événement ; \zh{继承} + bien (sans préposition) ; \zh{给} pour le bénéficiaire ; \zh{按照} + règle + virgule.
\item \textbf{Culture :} Cameroun = pluralisme (coutume \zh{习惯法} + loi moderne), conseil de famille \zh{家庭会议}, diversité ethnique. Chine = unification par \zh{民法典}, codification testaments, tradition \zh{宗法制}/\zh{祭祖} patrilinéaire.
\item \textbf{Compétences :} Comparer factuellement, utiliser le lexique spécialisé, exposer un point de vue structuré en chinois.
\end{itemize}
\end{aretenir}

\coursec{L'hérédité au Cameroun}

Dans la section précédente, nous avons défini l'hérédité (\zh{继承} jìchéng) comme la transmission des biens, des droits et des fonctions d'une personne disparue à ses descendants. Nous avons aussi vu que chaque société organise cette transmission selon ses propres règles. Nous allons maintenant observer concrètement comment cela fonctionne au Cameroun, un pays où coexistent plusieurs systèmes de succession. Observons d'abord un court texte, puis dégagons les notions essentielles.

\courssub{Texte de découverte : la succession chez les Mballa}

\begin{textebox}[Texte : La succession chez les Mballa]
在喀麦隆，很多地方儿子继承父亲的土地和房子。\\
Zài Kāmàilóng, hěn duō dìfang érzi jìchéng fùqin de tǔdì hé fángzi.\\
« Au Cameroun, dans beaucoup de régions, les fils héritent de la terre et de la maison du père. »\\[4pt]
Mballa 先生是雅温得附近一个村子的农民。他去世以后，他的大家庭开了一个会。家里的老人和酋长一起讨论：谁继承什么？按照他们的习俗，长子继承父亲的房子和最大的土地，其他的儿子也得到一些地。女儿们得到钱和东西。寡妇可以住在丈夫的房子里，也可以继承一部分财产。\\
Mballa xiānsheng shì Yǎwēndé fùjìn yí ge cūnzi de nóngmín. Tā qùshì yǐhòu, tā de dà jiātíng kāi le yí ge huì. Jiā lǐ de lǎorén hé qiúzhǎng yìqǐ tǎolùn : shéi jìchéng shénme ? Ànzhào tāmen de xísú, zhǎngzǐ jìchéng fùqin de fángzi hé zuì dà de tǔdì, qítā de érzi yě dédào yìxiē dì. Nǚ'ér men dédào qián hé dōngxi. Guǎfu kěyǐ zhù zài zhàngfu de fángzi lǐ, yě kěyǐ jìchéng yí bùfen cáichǎn.\\
« Monsieur Mballa était un paysan d'un village près de Yaoundé. Après son décès, sa famille élargie a tenu une réunion. Les anciens de la famille et le chef traditionnel ont discuté ensemble : qui hérite de quoi ? Selon leur coutume, le fils aîné hérite de la maison du père et du plus grand terrain, les autres fils reçoivent aussi des terres. Les filles reçoivent de l'argent et des biens. La veuve peut habiter dans la maison de son mari et peut aussi hériter d'une partie des biens. »\\[4pt]
可是，在喀麦隆西北部的一些地方，习俗不一样：在那里，孩子继承妈妈的家庭，舅舅（妈妈的兄弟）把财产给外甥。\\
Kěshì, zài Kāmàilóng xīběibù de yìxiē dìfang, xísú bù yíyàng : zài nàli, háizi jìchéng māma de jiātíng, jiùjiu (māma de xiōngdì) bǎ cáichǎn gěi wàisheng.\\
« Cependant, dans certaines régions du Nord-Ouest du Cameroun, la coutume est différente : là-bas, les enfants relèvent de la famille de la mère, et l'oncle maternel (le frère de la mère) transmet ses biens à son neveu. »
\end{textebox}

\textbf{Questions de compréhension}
\begin{enumerate}
\item Qui participe à la réunion de famille après le décès de M. Mballa ?
\item Que reçoit le fils aîné selon la coutume décrite ?
\item Quelle différence existe-t-il dans le Nord-Ouest du Cameroun ?
\end{enumerate}

\courssub{Observation et vocabulaire clé}

\begin{definition}[习俗 xísú]
\zh{习俗} (xísú) : nom, « coutume, pratique traditionnelle ». Le caractère \zh{习} évoque l'habitude, la répétition (on le retrouve dans \zh{学习} xuéxí, « étudier ») ; \zh{俗} évoque ce qui est commun au peuple. Une coutume est donc littéralement une « habitude du peuple ». Exemple : \zh{每个地方有不同的习俗。} (Měi ge dìfang yǒu bùtóng de xísú.) « Chaque région a des coutumes différentes. »
\end{definition}

\begin{center}
\begin{tabularx}{\linewidth}{|X|X|X|X|}
\hline
\rowcolor{popblueL} Caractères & Pinyin & Nature & Traduction \\
\hline
习俗 & xísú & nom & coutume \\
\hline
酋长 & qiúzhǎng & nom & chef traditionnel \\
\hline
长子 & zhǎngzǐ & nom & fils aîné \\
\hline
寡妇 & guǎfu & nom & veuve \\
\hline
大家庭 & dà jiātíng & nom & famille élargie \\
\hline
继承 & jìchéng & verbe & hériter, succéder à \\
\hline
财产 & cáichǎn & nom & biens, patrimoine \\
\hline
土地 & tǔdì & nom & terre, terrain \\
\hline
舅舅 & jiùjiu & nom & oncle maternel (frère de la mère) \\
\hline
外甥 & wàisheng & nom & neveu (fils de la sœur) \\
\hline
去世 & qùshì & verbe & décéder (terme respectueux) \\
\hline
按照 & ànzhào & préposition & selon, conformément à \\
\hline
\end{tabularx}
\end{center}

\begin{attention}[家 : famille ou maison ?]
Le caractère \zh{家} (jiā) signifie à la fois « famille » et « maison / foyer ». Dans \zh{大家庭} (dà jiātíng), il s'agit bien de la \cle{famille élargie} (grands-parents, oncles, tantes, cousins), et non d'une « grande maison ». Pour parler du bâtiment, on dirait \zh{大房子} (dà fángzi). Ne traduisez donc jamais \zh{大家庭} par « grande maison » !
\end{attention}

\courssub{Deux grands systèmes de filiation au Cameroun}

\begin{propriete}[Filiation patrilinéaire et matrilinéaire]
Au Cameroun, il n'existe pas un seul système d'hérédité, mais plusieurs :
\begin{itemize}
\item La \cle{filiation patrilinéaire} : l'appartenance au lignage et l'héritage passent par le père. C'est le système le plus répandu, notamment chez de nombreux groupes du Centre, du Sud et de l'Ouest. Les biens (terres, maison) vont du père aux fils, souvent avec une place privilégiée pour le fils aîné (\zh{长子}).
\item La \cle{filiation matrilinéaire} : l'appartenance au lignage passe par la mère. Chez certains peuples du Nord-Ouest, comme les Kom, l'héritage ne va pas du père au fils, mais de l'oncle maternel (\zh{舅舅}, frère de la mère) au neveu (\zh{外甥}, fils de la sœur).
\end{itemize}
\end{propriete}

\begin{propriete}[Structure \zh{把} (bǎ) : action sur un objet défini]
La structure \zh{把} (bǎ) permet de placer l'objet \emph{avant} le verbe pour insister sur le résultat de l'action ou sa manipulation. Elle est obligatoire avec certains verbes (donner, mettre, prendre) quand l'objet est défini.
\medskip
\noindent \textbf{Schéma :} Sujet + \zh{把} + Objet + Verbe + Complément de résultat / direction / particule \zh{了} / etc.
\medskip
\noindent \textbf{Exemple du texte :} \zh{舅舅把财产给外甥。} (Jiùjiu bǎ cáichǎn gěi wàisheng.) « L'oncle maternel donne les biens au neveu. »
\medskip
\noindent \textbf{Attention :} On ne dit pas \zh{舅舅给外甥财产} (structure SVO simple) si l'on veut insister sur la \emph{transmission effective} du patrimoine. \zh{把} marque la disposition de l'objet.
\end{propriete}

\begin{popfigure}
\begin{tikzpicture}[
  node distance=0.9cm and 1.6cm,
  boite/.style={draw=popblue, fill=popblueL, rounded corners, align=center, minimum width=2.6cm, minimum height=0.8cm, font=\small},
  boiteM/.style={draw=poppurple, fill=poppurpleL, rounded corners, align=center, minimum width=2.6cm, minimum height=0.8cm, font=\small},
  fleche/.style={-{Stealth[length=3mm]}, thick, popdark},
  titre/.style={font=\bfseries\small, align=center}
]
% Patrilinéaire
\node[titre, popblue, align=center] (t1) at (0,0){Système patrilinéaire\\父系 fùxì};
\node[boite, below=of t1, align=center] (pere){父亲 fùqīn\\père};
\node[boite, below left=0.9cm and 0.4cm of pere, align=center] (aine){长子 zhǎngzǐ\\fils aîné};
\node[boite, below right=0.9cm and 0.4cm of pere, align=center] (autres){其他的儿子\\autres fils};
\draw[fleche] (pere) -- (aine);
\draw[fleche] (pere) -- (autres);
% Matrilinéaire
\node[titre, poppurple, align=center] (t2) at (8.5,0){Système matrilinéaire\\母系 mǔxì};
\node[boiteM, below=of t2, align=center] (mere){母亲 mǔqīn\\mère};
\node[boiteM, below left=0.9cm and 0.4cm of mere, align=center] (oncle){舅舅 jiùjiu\\frère de la mère};
\node[boiteM, below right=0.9cm and 0.4cm of mere, align=center] (neveu){外甥 wàisheng\\neveu};
\draw[fleche] (mere) -- (oncle);
\draw[fleche] (oncle) -- (neveu);
\end{tikzpicture}
\legende{Arbres de succession : à gauche, transmission du père vers les fils ; à droite, transmission de l'oncle maternel vers le neveu.}
\end{popfigure}

\courssub{Le rôle de la famille élargie et des notables}

Au Cameroun, la succession n'est pas une affaire purement individuelle : c'est la \cle{famille élargie} (\zh{大家庭}) qui se réunit, souvent en présence des anciens et du chef traditionnel (\zh{酋长}), pour organiser le partage. Cette réunion, parfois appelée « conseil de famille », décide de la répartition des terres, de la maison, de l'argent, mais aussi de la prise en charge des veuves et des enfants.

\begin{exemplebox}[Exemples traduits]
\zh{大家庭一起决定。}\\
\emph{Dà jiātíng yìqǐ juédìng.}\\
« La famille élargie décide ensemble. »\\[4pt]
\zh{寡妇也可以继承。}\\
\emph{Guǎfu yě kěyǐ jìchéng.}\\
« La veuve aussi peut hériter. »\\[4pt]
\zh{按照习俗，长子继承父亲的房子。}\\
\emph{Ànzhào xísú, zhǎngzǐ jìchéng fùqin de fángzi.}\\
« Selon la coutume, le fils aîné hérite de la maison du père. »\\[4pt]
\zh{酋长和老人一起开会。}\\
\emph{Qiúzhǎng hé lǎorén yìqǐ kāihuì.}\\
« Le chef traditionnel et les anciens tiennent une réunion ensemble. »\\[4pt]
\zh{舅舅把财产给外甥。}\\
\emph{Jiùjiu bǎ cáichǎn gěi wàisheng.}\\
« L'oncle maternel transmet les biens au neveu. »
\end{exemplebox}

Remarquez deux structures chinoises essentielles :
\begin{itemize}
\item \cle{Sujet + 一起 + Verbe} (« faire ensemble ») : \zh{大家庭一起决定}。
\item \cle{按照 + Nom + ，+ Phrase} (« selon... ») : la préposition \zh{按照} se place toujours \emph{avant} le verbe principal, en tête de proposition, contrairement au français où « selon la coutume » peut se placer à la fin.
\item \cle{Sujet + 把 + Objet + Verbe} : structure de disposition (voir boîte de règle ci-dessus).
\end{itemize}

\courssub{Place de l'aîné, des fils, des veuves et des filles}

Traditionnellement, dans les systèmes patrilinéaires, le \cle{fils aîné} (\zh{长子}) occupe une place centrale : il reçoit souvent la maison familiale et la part la plus importante, mais il hérite aussi de \emph{responsabilités} (s'occuper de sa mère, de ses frères et sœurs cadets, gérer les affaires du lignage). Les filles, selon certaines coutumes, reçoivent des biens mobiliers ou de l'argent (\zh{钱和东西}) plutôt que la terre immobilière, car elles sont censées rejoindre la famille de leur mari. Les veuves (\zh{寡妇}), selon les régions et les époques, pouvaient être « héritées » avec le patrimoine (lévirat) ou simplement autorisées à rester dans la maison (\zh{住在丈夫的房子里}) ; aujourd'hui, ces pratiques évoluent fortement et le droit moderne protège mieux les droits des veuves et des filles.

\courssub{Droit coutumier et droit moderne : une coexistence}

Le Cameroun connaît une \cle{coexistence} entre le droit coutumier (\zh{习俗法} xísú fǎ) et le droit moderne (\zh{现代法律} xiàndài fǎlǜ : Code civil, droit de la famille). En pratique :
\begin{itemize}
\item Dans beaucoup de villages, le partage se fait d'abord selon la coutume, devant la famille et les notables (\zh{酋长}, \zh{老人}).
\item En ville, et en cas de conflit, les tribunaux appliquent le droit écrit, qui reconnaît des droits successoraux égaux à l'épouse et aux enfants, filles comme garçons (\zh{男女平等} nánnǚ píngděng).
\item Il arrive donc qu'une même succession soit discutée deux fois : une fois « à la maison » selon la coutume, une fois devant le juge selon la loi.
\end{itemize}

\courssub{Héritage immatériel : la succession au trône}

L'hérédité ne concerne pas seulement les biens matériels : elle concerne aussi les \cle{fonctions}. Dans les chefferies et royaumes bamilékés de l'Ouest, comme dans les sultanats et lamidats du Nord, le successeur du chef (\zh{酋长}) est désigné selon des règles coutumières précises, au sein d'une famille royale ou d'un lignage particulier. Le nouveau chef hérite du trône, des insignes et des devoirs rituels de son prédécesseur.

\begin{savaistu}[Un successeur choisi en secret]
Dans certaines chefferies de l'Ouest camerounais, le successeur du chef est choisi \emph{en secret} parmi les fils du défunt par un collège de notables, et son nom n'est révélé qu'après les funérailles. Cette discrétion protège le futur chef et préserve l'unité de la chefferie pendant le deuil.
\end{savaistu}

\courssub{Tableau récapitulatif}

\begin{center}
\begin{tabularx}{\linewidth}{|X|X|X|}
\hline
\rowcolor{popblueL} Notion & Chinois + pinyin & Réalité camerounaise \\
\hline
Coutume & 习俗 xísú & règles traditionnelles de partage, variables selon les régions \\
\hline
Filiation patrilinéaire & 父系继承 fùxì jìchéng & Centre, Sud, Ouest : père $\rightarrow$ fils \\
\hline
Filiation matrilinéaire & 母系继承 mǔxì jìchéng & Nord-Ouest (ex. Kom) : oncle maternel $\rightarrow$ neveu \\
\hline
Famille élargie & 大家庭 dà jiātíng & conseil de famille qui décide du partage \\
\hline
Chef traditionnel & 酋长 qiúzhǎng & arbitre la succession ; succession royale dans les chefferies \\
\hline
Fils aîné & 长子 zhǎngzǐ & part privilégiée + responsabilités \\
\hline
Veuve & 寡妇 guǎfu & droit d'habiter la maison ; droits renforcés par le droit moderne \\
\hline
Droit moderne & 现代法律 xiàndài fǎlǜ & Code civil, protection des épouses et des filles \\
\hline
Structure \zh{把} & 把字句 bǎzìjù & action sur objet défini : \zh{舅舅把财产给外甥} \\
\hline
\end{tabularx}
\end{center}

\courssub{Exercice résolu et pièges à éviter}

\begin{exoresolu}[Traduisez en chinois]
Énoncé : Traduisez en chinois : « Selon la coutume, la famille élargie décide ensemble, et la veuve peut aussi hériter. »
\tcblower
Solution détaillée :
\begin{enumerate}
\item « Selon la coutume » se place en tête : \zh{按照习俗} (ànzhào xísú), suivi d'une virgule chinoise \zh{，}.
\item « La famille élargie décide ensemble » : sujet \zh{大家庭} + \zh{一起} + verbe \zh{决定} : \zh{大家庭一起决定}。
\item « Et la veuve peut aussi hériter » : \zh{寡妇也可以继承} (l'adverbe \zh{也} « aussi » se place avant le verbe modal \zh{可以}).
\end{enumerate}
Phrase complète : \zh{按照习俗，大家庭一起决定，寡妇也可以继承。}\\
\emph{Ànzhào xísú, dà jiātíng yìqǐ juédìng, guǎfu yě kěyǐ jìchéng.}
\end{exoresolu}

\begin{attention}[Erreurs fréquentes des francophones]
\begin{itemize}
\item Ne placez pas \zh{按照习俗} à la fin de la phrase comme en français (« ...selon la coutume ») : en chinois, ce complément se met avant le verbe ou en tête de phrase.
\item Attention aux tons de \zh{长子} : \zh{长} se lit ici \zh{zhǎng} (3\up{e} ton, « aîné »), et non \zh{cháng} (« long »).
\item \zh{寡妇} : le second caractère \zh{妇} se lit en ton léger (\zh{fu}), comme dans beaucoup de mots de parenté (\zh{夫人} fūrén, \zh{媳妇} xífu).
\item N'oubliez pas la virgule pleine chasse \zh{，} entre les propositions, et le point chinois \zh{。} en fin de phrase.
\item Avec la structure \zh{把}, l'objet doit être \emph{défini} (connu, spécifique) : \zh{把财产} (les biens), pas \zh{把一笔财产} (une somme de biens — indéfini).
\end{itemize}
\end{attention}

\begin{exercice}[À vous de jouer]
\begin{enumerate}
\item Traduisez en français : \zh{在西北，舅舅把财产给外甥。}
\item Traduisez en chinois : « Le fils aîné hérite de la terre du père. »
\item Répondez en une phrase chinoise simple : dans votre région, qui décide du partage d'une succession ?
\end{enumerate}
\end{exercice}

\begin{corrige}[Exercice « À vous de jouer »]
\begin{enumerate}
\item « Dans le Nord-Ouest, l'oncle maternel transmet les biens au neveu. » Structure \zh{把} : Sujet \zh{舅舅} + \zh{把} + Objet \zh{财产} + Verbe \zh{给} + Destinataire \zh{外甥}。
\item \zh{长子继承父亲的土地。} (Zhǎngzǐ jìchéng fùqin de tǔdì.)
\item Réponse possible : \zh{大家庭和酋长一起决定。} (Dà jiātíng hé qiúzhǎng yìqǐ juédìng.) « La famille élargie et le chef décident ensemble. »
\end{enumerate}
\end{corrige}

\begin{aretenir}[L'essentiel de la section]
Au Cameroun, l'hérédité repose sur des systèmes variés : patrilinéaire dans de nombreuses régions (père $\rightarrow$ fils, place privilégiée de l'aîné \zh{长子}), matrilinéaire chez certains peuples du Nord-Ouest (oncle maternel \zh{舅舅} $\rightarrow$ neveu \zh{外甥}). La famille élargie (\zh{大家庭}) et les notables (\zh{酋长}) organisent le partage, dans un dialogue constant entre droit coutumier (\zh{习俗}) et droit moderne (\zh{现代法律}). L'hérédité concerne aussi les fonctions, comme la succession au trône dans les chefferies. La structure \zh{把} (Sujet + \zh{把} + Objet + Verbe) est indispensable pour exprimer la transmission effective d'un patrimoine défini. La section suivante examinera comment ces questions se posent en Chine.
\end{aretenir}

\coursec{Leçon 4 — Éléments socioculturels : processus d'hérédité au Cameroun et en Chine}

\courssub{Introduction : qu'est-ce que l'hérédité ?}

L'hérédité (\zh{继承}, jìchéng) désigne la transmission du patrimoine (biens, titres, noms, valeurs) d'une génération à l'autre. Elle repose sur deux piliers qui varient selon les sociétés : la \cle{coutume} (\zh{习俗}, xísú) et la \cle{loi écrite} (\zh{成文法}, chéngwénfǎ). Au Cameroun comme en Chine, on observe une tension historique entre une tradition forte (souvent patrilinéaire) et une législation moderne qui proclame l'égalité. Comprendre ces processus, c'est comprendre la place de l'individu dans la famille et la société.

\begin{definition}[Vocabulaire de base : la succession]
\begin{itemize}
\item \zh{遗产} (yíchǎn) : héritage, patrimoine, biens laissés par un défunt.
\item \zh{继承人} (jìchéngrén) : héritier, personne qui a droit à la succession.
\item \zh{立遗嘱人} (lì yízhǔ rén) : testateur, celui qui fait un testament.
\item \zh{法定继承} (fǎdìng jìchéng) : succession légale (sans testament), selon l'ordre prévu par la loi.
\item \zh{习惯法} (xíguànfǎ) : droit coutumier, ensemble de règles non écrites transmises oralement.
\end{itemize}
\end{definition}

\begin{textebox}[Dialogue : deux élèves comparent]
--- \zh{小张，你们喀麦隆人怎么分遗产？} (Xiǎo Zhāng, nǐmen Kāmàilóng rén zěnme fēn yíchǎn?) « Petit Zhang, comment partagez-vous l'héritage au Cameroun ? »\\
--- \zh{这很复杂。有习惯法，也有民法典。传统上，儿子继承，但现在女儿也有权利。} (Zhè hěn fùzá. Yǒu xíguànfǎ, yě yǒu Mínfǎdiǎn. Chuántǒng shàng, érzi jìchéng, dàn xiànzài nǚ'ér yě yǒu quánlì.) « C'est complexe. Il y a le droit coutumier et le Code civil. Traditionnellement, les fils héritent, mais maintenant les filles ont aussi des droits. »\\
--- \zh{中国也一样。以前重男轻女，现在法律规定男女平等。} (Zhōngguó yě yīyàng. Yǐqián zhòng nán qīng nǚ, xiànzài fǎlǜ guīdìng nán nǚ píngděng.) « La Chine c'est pareil. Autrefois on valorisait les garçons et on méprisait les filles, maintenant la loi établit l'égalité hommes-femmes. »
\end{textebox}

\courssub{L'hérédité au Cameroun : entre coutume et droit moderne}

Le Cameroun est un pays \cle{bijuridique} : le droit coutumier (variable selon les ethnies : Bamiléké, Bassa, Béti, Peul, etc.) et le droit moderne (Code civil, Ordinance 81-02 du 29 juin 1981) coexistent. Cette dualité crée des situations concrètes que les élèves doivent savoir décrire en chinois.

\begin{textebox}[L'héritage au Cameroun : coutume et loi]
喀麦隆有两种法律：习惯法和现代法。\\
Kāmàilóng yǒu liǎng zhǒng fǎlǜ: xíguànfǎ hé xiàndài fǎ.\\
« Le Cameroun a deux sortes de lois : le droit coutumier et le droit moderne. »\\[0.4em]
传统上，很多民族实行父系继承，只有儿子继承土地和房子。女儿嫁人后，通常不分遗产。\\
Chuántǒng shàng, hěn duō mínzú shíxíng fùxì jìchéng, zhǐyǒu érzi jìchéng tǔdì hé fángzi. Nǚ'ér jiàrén hòu, tōngcháng bù fēn yíchǎn.\\
« Traditionnellement, beaucoup de groupes ethniques pratiquent la succession patrilinéaire : seuls les fils héritent de la terre et des maisons. Après leur mariage, les filles ne partagent généralement pas l'héritage. »\\[0.4em]
但是，现代法律《民法典》规定男女平等，女儿也有继承权。在城市，很多家庭现在写遗嘱，公平分配财产。\\
Dànshì, xiàndài fǎlǜ « Mínfǎdiǎn » guīdìng nán nǚ píngděng, nǚ'ér yě yǒu jìchéngquán. Zài chéngshì, hěn duō jiātíng xiànzài xiě yízhǔ, gōngpíng fēnpèi cáichǎn.\\
« Cependant, le droit moderne (Code civil) établit l'égalité hommes-femmes, les filles ont aussi des droits de succession. En ville, beaucoup de familles font maintenant un testament pour répartir équitablement les biens. »
\end{textebox}

\begin{propriete}[Structure : 虽然...但是... (Suīrán... dànshì...) — Concession]
Pour exprimer une opposition (coutume vs loi), on utilise \zh{虽然} (bien que) en début de proposition et \zh{但是} / \zh{可是} (mais) au début de la suivante. Le sujet peut être répété ou omis s'il est identique.\\[0.3em]
Schéma : \textbf{虽然 + Proposition 1, 但是 + Proposition 2.}\\[0.3em]
Exemple : \zh{虽然习惯法不让女儿继承，但是现代法规定平等。} (Suīrán xíguànfǎ bù ràng nǚ'ér jìchéng, dànshì xiàndài fǎ guīdìng píngděng.) « Bien que le droit coutumier n'autorise pas les filles à hériter, la loi moderne établit l'égalité. »
\end{propriete}

\begin{attention}[Erreur fréquente : confusion « mais » / « et »]
En français, on peut dire « Bien qu'il soit pauvre, \emph{et} il est heureux » (incorrect en français standard mais structure possible en chinois avec 而且). En chinois, \zh{虽然} \textbf{appelle obligatoirement} \zh{但是} / \zh{可是} / \zh{却}. Ne mettez jamais \zh{而且} (et de plus) après \zh{虽然}.
\end{attention}

\begin{savaistu}[Diversité camerounaise : le cas matrilinéaire]
Si la majorité des groupes (Béti, Bamiléké, Peuls) sont \cle{patrilinéaires} (héritage par les fils, \zh{父系继承}), certains groupes comme les \cle{Maka} ou \cle{Kaka} (région Est) pratiquent une succession \cle{matrilinéaire} (\zh{母系继承}, mǔxì jìchéng) : l'héritier principal est le neveu (fils de la sœur du défunt), non le fils. En chinois : \zh{侄子继承舅舅的遗产} (Zhízǐ jìchéng jiùjiu de yíchǎn). Cette diversité montre qu'il n'existe pas « une » coutume camerounaise unique.
\end{savaistu}

\courssub{L'hérédité en Chine : tradition et modernité}

Après avoir étudié les pratiques successorales au Cameroun — où le droit coutumier et le droit moderne coexistent —, nous traversons maintenant l'Asie pour observer la situation chinoise. Nous allons voir que la Chine connaît une évolution comparable : une tradition très ancienne fondée sur la primauté du fils, puis une loi moderne qui établit l'égalité entre fils et filles. Comparer les deux pays nous aidera à mieux comprendre notre propre société.

\courssub{Observation : un texte pour découvrir}

Lisons d'abord ce court texte. Soulignez mentalement les mots qui expriment le temps (autrefois, aujourd'hui) et ceux qui désignent les membres de la famille.

\begin{textebox}[L'héritage en Chine : hier et aujourd'hui]
以前在中国，只有儿子能继承遗产。女儿结婚以后，就住在丈夫的家里，不能继承父母的遗产。\\
Yǐqián zài Zhōngguó, zhǐyǒu érzi néng jìchéng yíchǎn. Nǚ'ér jiéhūn yǐhòu, jiù zhù zài zhàngfu de jiā lǐ, bù néng jìchéng fùmǔ de yíchǎn.\\
« Autrefois en Chine, seuls les fils pouvaient hériter des biens. Après leur mariage, les filles vivaient dans la famille de leur mari et ne pouvaient pas hériter des biens de leurs parents. »\\[0.4em]
现在法律规定儿子和女儿平等。儿子和女儿都可以继承父母的遗产。丈夫死了以后，妻子也可以继承他的遗产。\\
Xiànzài fǎlǜ guīdìng érzi hé nǚ'ér píngděng. Érzi hé nǚ'ér dōu kěyǐ jìchéng fùmǔ de yíchǎn. Zhàngfu sǐ le yǐhòu, qīzi yě kěyǐ jìchéng tā de yíchǎn.\\
« Aujourd'hui, la loi établit l'égalité entre fils et filles. Les fils et les filles peuvent tous hériter des biens de leurs parents. Quand le mari meurt, l'épouse peut aussi hériter de ses biens. »\\[0.4em]
很多中国家庭还有家谱，家谱里有祖先的名字。每年清明节，全家人一起祭祖。\\
Hěn duō Zhōngguó jiātíng hái yǒu jiāpǔ, jiāpǔ lǐ yǒu zǔxiān de míngzi. Měi nián Qīngmíng Jié, quán jiā rén yìqǐ jìzǔ.\\
« Beaucoup de familles chinoises possèdent encore un livre généalogique, dans lequel figurent les noms des ancêtres. Chaque année, à la fête de Qingming, toute la famille honore ensemble les ancêtres. »
\end{textebox}

\textbf{Questions de compréhension}\par
\begin{enumerate}
\item Autrefois, qui pouvait hériter en Chine ?
\item Que dit la loi aujourd'hui ?
\item Qu'est-ce qu'un \zh{家谱} ?
\item Que font les familles à la fête de Qingming (\zh{清明节}) ?
\end{enumerate}

\courssub{La tradition confucéenne : la primauté du fils}

Pendant plus de deux mille ans, la société chinoise a été organisée selon les idées du philosophe Confucius (\zh{孔子}, Kǒngzǐ). Dans cette tradition, la famille est au centre de tout, et le respect des parents et des ancêtres (\zh{祖先}, zǔxiān) est une valeur fondamentale.

Dans le système traditionnel :

\begin{itemize}
\item le \cle{nom de famille} (\zh{姓}, xìng) se transmet par les fils ; une fille qui se marie « sort » de sa famille natale (\zh{出嫁}, chūjià) pour entrer dans celle de son mari ;
\item les biens de la famille (terres, maison, argent) sont partagés \cle{entre les fils} à la mort du père ; le fils aîné (\zh{大儿子}, dà érzi) joue souvent un rôle de chef de famille (\zh{当家}, dāngjiā) ;
\item les filles sont traditionnellement \cle{exclues de la succession}, car on considère qu'elles appartiennent désormais à la famille de leur époux (\zh{泼出去的水}, pō chūqù de shuǐ — proverbe : « l'eau versée ne revient pas ») ;
\item le culte des ancêtres (\zh{祭祖}, jìzǔ) est assuré par les fils, qui entretiennent les tombes (\zh{坟墓}, fénmù) et les tablettes des ancêtres (\zh{神主牌}, shénzhǔpái).
\end{itemize}

\begin{savaistu}[Le livre généalogique 家谱]
Certaines familles chinoises conservent un \zh{家谱} (jiāpǔ), un livre généalogique qui enregistre les noms des membres de la famille, génération après génération, parfois sur plus de vingt générations ! Le caractère \zh{家} (jiā) signifie « maison, famille » (clé \zh{宀} toit + \zh{豕} cochon : le cochon sous le toit = prospérité) et \zh{谱} (pǔ) signifie « registre, partition » (clé \zh{讠} parole + \zh{普} universel). Au Cameroun, la mémoire des ancêtres se transmet plutôt oralement, par les récits des anciens (\zh{口述历史}, kǒushù lìshǐ) lors des veillées : deux façons différentes mais complémentaires de garder la mémoire familiale.
\end{savaistu}

\courssub{La loi de 1985 : l'égalité entre fils et filles}

En 1949, avec la fondation de la République populaire de Chine, de profonds changements sociaux commencent. Mais c'est surtout la \cle{loi sur la succession de 1985} (\zh{继承法}, jìchéngfǎ), révisée en 2021 (Code civil), qui fixe clairement les nouvelles règles :

\begin{itemize}
\item \cle{égalité légale} (\zh{法定平等}, fǎdìng píngděng) entre les fils et les filles : tous les enfants héritent de la même manière (premier ordre : conjoint, enfants, parents) ;
\item le \cle{conjoint survivant} (\zh{配偶者}, pèiǒuzhě) fait partie des héritiers de premier ordre, au même titre que les enfants et les parents du défunt ;
\item le \cle{testament} (\zh{遗嘱}, yízhǔ) est reconnu sous plusieurs formes (écrit, enregistré, vidéo, oral en urgence) : chacun peut décider par écrit de la répartition de ses biens. Autrefois rare, le testament est aujourd'hui de plus en plus utilisé dans les grandes villes comme Pékin (\zh{北京}, Běijīng) ou Shanghai (\zh{上海}, Shànghǎi).
\end{itemize}

\begin{definition}[Deux mots essentiels]
\zh{法律} (fǎlǜ), nom : la loi, l'ensemble des règles écrites d'un pays. Exemple : \zh{中国有一部继承法。} (Zhōngguó yǒu yí bù jìchéngfǎ.) « La Chine a une loi sur la succession. »\\[0.3em]
\zh{平等} (píngděng), adjectif/nom : égal, l'égalité. Exemple : \zh{儿子和女儿平等。} (Érzi hé nǚ'ér píngděng.) « Les fils et les filles sont égaux. »\\[0.3em]
\zh{配偶者} (pèiǒuzhě), nom : conjoint (époux ou épouse), terme juridique neutre.
\end{definition}

\begin{popfigure}
\begin{tikzpicture}[font=\small]
\draw[line width=1.2pt, popblue, -{Stealth[length=3mm]}] (0,0) -- (13.5,0);
\foreach \x/\date/\txt/\col in {
  1.2/{Tradition}/{Primauté du fils aîné ;\\filles exclues}/{poporange},
  4.9/{1949}/{Fondation de la\\Rép. populaire de Chine}/{poppurple},
  8.6/{1985}/{Loi sur la succession :\\égalité fils--filles}/{popgreen},
  12.1/{2021/ Aujourd'hui}/{Code civil : testaments\\vidéo, numériques}/{popblue}}{
  \fill[\col] (\x,0) circle (0.12);
  \node[above, align=center, text=\col] at (\x,0.15) {\textbf{\date}};
  \node[below, align=center, text width=2.9cm] at (\x,-0.15) {\txt};
}
\end{tikzpicture}
\legende{Frise chronologique : l'évolution de l'héritage en Chine}
\end{popfigure}

\courssub{Point de grammaire : 以前 et 现在}

Dans le texte, observez la place des mots de temps :

\zh{以前在中国，只有儿子能继承遗产。} (Yǐqián zài Zhōngguó, zhǐyǒu érzi néng jìchéng yíchǎn.) « Autrefois en Chine, seuls les fils pouvaient hériter. »

\zh{现在法律规定儿子和女儿平等。} (Xiànzài fǎlǜ guīdìng érzi hé nǚ'ér píngděng.) « Aujourd'hui, la loi établit l'égalité entre fils et filles. »

\begin{propriete}[La place de 以前 et 现在]
Les compléments de temps comme \zh{以前} (yǐqián, autrefois) et \zh{现在} (xiànzài, maintenant) se placent \cle{avant le verbe}, soit au tout début de la phrase (Thème), soit \cle{juste après le sujet} — mais \textbf{jamais à la fin} de la phrase comme souvent en français.\\[0.3em]
Structure : \textbf{Sujet + 以前 / 现在 + Verbe + ...}\\[0.3em]
Exemple : \zh{我以前住在北京。} (Wǒ yǐqián zhù zài Běijīng.) « J'habitais à Pékin autrefois. »\\
Exemple : \zh{我现在学习汉语。} (Wǒ xiànzài xuéxí Hànyǔ.) « Maintenant, j'apprends le chinois. »\\
Exemple (début de phrase) : \zh{以前没有电脑，现在人人有手机。} (Yǐqián méiyǒu diànnǎo, xiànzài rénrén yǒu shǒujī.) « Autrefois pas d'ordinateurs, maintenant tout le monde a un smartphone. »
\end{propriete}

\begin{attention}[Erreur fréquente des francophones]
En français, on dit « J'habitais à Pékin \emph{autrefois} », avec le mot de temps à la fin. En chinois, on ne peut pas dire \zh{我住在北京以前} : la phrase est incorrecte. Dites toujours \zh{我以前住在北京。} Le complément de temps précède le verbe.
\end{attention}

\courssub{Comparaison Cameroun / Chine : faits comparés}

Le tableau ci-dessous synthétise les points clés pour votre expression orale et écrite. Utilisez les connecteurs logiques : \zh{相同点} (points communs), \zh{不同点} (différences), \zh{一方面...另一方面...} (d'un côté... de l'autre...).

\begin{center}
\begin{tabularx}{\linewidth}{|X|X|X|}
\hline
\rowcolor{popblueL} \textbf{Critère} & \textbf{Cameroun} & \textbf{Chine} \\
\hline
\textbf{Base traditionnelle} & Droit coutumier variable (patrilinéaire majoritaire, matrilinéaire minoritaire). Transmission orale. & Tradition confucéenne unifiée : patrilinéaire stricte, primogéniture, exclusion des filles. Écrits : \zh{家谱}. \\
\hline
\textbf{Loi moderne} & Code civil (Ordonnance 81-02, 2016) : égalité fils/filles, conjoint héritier. Pluralisme juridique (coutume vs loi). & Loi sur la succession 1985 / Code civil 2021 : égalité stricte, testament diversifié. Unité juridique nationale. \\
\hline
\textbf{Conjoint survivant} & Héritier (droit moderne), mais souvent exclu par la coutume (biens aux enfants/neveux). & Héritier de premier ordre (même part que les enfants), protégé par la loi. \\
\hline
\textbf{Héritage immatériel} & Rites funéraires, veillées, transmission orale des noms/titres, totems. & Culte des ancêtres (\zh{祭祖}), fête de Qingming (\zh{清明节}), généalogies (\zh{家谱}), transmission du nom (\zh{姓}). \\
\hline
\textbf{Évolution actuelle} & Urbanisation : testaments plus fréquents, conflits coutume/loi devant les tribunaux. & Urbanisation : testaments numériques, vidéo, planification successorale, femmes propriétaires. \\
\hline
\end{tabularx}
\end{center}

\begin{methode}[Comment rédiger un paragraphe de comparaison (Expression écrite)]
\begin{enumerate}
\item \textbf{Introduire le thème} : \zh{中喀两国在继承问题上既有相同点，也有不同点。} (Zhōng Kā liǎng guó zài jìchéng wèntí shàng jì yǒu xiāngtóngdiǎn, yě yǒu bùtóngdiǎn.) « La Chine et le Cameroun ont à la fois des points communs et des différences sur la question de l'héritage. »
\item \textbf{Les similitudes} : \zh{相同点是：传统上都重视儿子，现代法律都规定男女平等。} (Xiāngtóngdiǎn shì: chuántǒng shàng dōu zhòngshì érzi, xiàndài fǎlǜ dōu guīdìng nán nǚ píngděng.)
\item \textbf{Les différences} : \zh{不同点是：喀麦隆有习惯法和现代法并存，中国法律统一。} (Bùtóngdiǎn shì: Kāmàilóng yǒu xíguànfǎ hé xiàndài fǎ bìngcún, Zhōngguó fǎlǜ tǒngyī.)
\item \textbf{Conclure} : \zh{总的来说，两国都在从传统走向现代法治。} (Zǒngde lái shuō, liǎng guó dōu zài cóng chuántǒng zǒuxiàng xiàndài fǎzhì.) « En résumé, les deux pays passent de la tradition à l'État de droit moderne. »
\end{enumerate}
\end{methode}

\courssub{Vocabulaire, écriture des caractères et révision}

\begin{center}
\begin{tabularx}{\linewidth}{|X|X|X|X|}
\hline
\rowcolor{popblueL} Caractères & Pinyin & Nature & Traduction \\
\hline
法律 & fǎlǜ & nom & loi \\
\hline
平等 & píngděng & adjectif/nom & égal, égalité \\
\hline
儿子 & érzi & nom & fils \\
\hline
女儿 & nǚ'ér & nom & fille \\
\hline
祖先 & zǔxiān & nom & ancêtres \\
\hline
家谱 & jiāpǔ & nom & arbre généalogique, livre généalogique \\
\hline
继承 & jìchéng & verbe & hériter, succéder \\
\hline
遗产 & yíchǎn & nom & héritage, patrimoine \\
\hline
遗嘱 & yízhǔ & nom & testament \\
\hline
祭祖 & jìzǔ & verbe & honorer les ancêtres (culte) \\
\hline
规定 & guīdìng & verbe & stipuler, prescrire (loi, règlement) \\
\hline
以前 & yǐqián & nom de temps & autrefois, avant \\
\hline
现在 & xiànzài & nom de temps & maintenant, actuellement \\
\hline
习惯法 & xíguànfǎ & nom & droit coutumier \\
\hline
民法典 & Mínfǎdiǎn & nom propre & Code civil \\
\hline
父系/母系 & fùxì / mǔxì & adj/nom & patrilinéaire / matrilinéaire \\
\hline
配偶者 & pèiǒuzhě & nom & conjoint (juridique) \\
\hline
清明节 & Qīngmíng Jié & nom propre & Fête de la Pure Clarté (début avril) \\
\hline
\end{tabularx}
\end{center}

\begin{propriete}[Écriture des caractères clés : ordre des traits et clés]
\begin{itemize}
\item \zh{法} (fǎ, loi) : 8 traits. Clé \zh{氵} (eau, 3 traits) + \zh{去} (qù). Ordre : point, point, trait montant (eau), puis \zh{去} (haut→bas : \zh{土} puis \zh{厶}). Mnémo : « L'eau qui part = la règle qui s'impose ».
\item \zh{平} (píng, égal) : 5 traits. Clé \zh{干} (gān). Ordre : horizontal, vertical, horizontal, horizontal, horizontal final long. Symétrie parfaite = égalité.
\item \zh{继} (jì, succéder) : 10 traits. Clé \zh{纟} (soie, 2 traits) + \zh{米} (riz) + \zh{\&} (partie phonétique). Ordre : soie (point, pli), puis riz (point, horizontal, vertical, deux points), puis le reste.
\item \zh{遗} (yí, laisser/héritage) : 12 traits. Clé \zh{辶} (marche, 3 traits, s'écrit EN DERNIER). Haut : \zh{贵} (précieux) sans le point du bas. « Ce qui est précieux et qu'on laisse derrière ».
\item \zh{祖} (zǔ, ancêtre) : 9 traits. Clé \zh{礻} (shì, autel, 4 traits : point, horizontal, vertical, point) + \zh{且} (qiě). Ordre : autel d'abord, puis \zh{且}.
\end{itemize}
\end{propriete}

\begin{attention}[Ne confondez pas 儿子 et 女儿 — Piège tonal]
\zh{女儿} (nǚ'ér, fille) et \zh{儿子} (érzi, fils) partagent \zh{儿} (enfant/fils).
\begin{itemize}
\item \zh{儿} seul = \emph{ér} (2\textsuperscript{e} ton, montant).
\item Dans \zh{儿子} : \emph{érzi} (2 + ton neutre).
\item Dans \zh{女儿} : \emph{nǚ'ér} (3\textsuperscript{e} ton + 2\textsuperscript{e} ton avec \emph{ér} qui garde son ton montant car c'est un suffixe nominalisé). \textbf{Attention} : \zh{女} = \emph{nǚ} (3\textsuperscript{e} ton, bas-remontant), pas \emph{nú} ni \emph{nù}.
\end{itemize}
Exercice : Lisez à voix haute 3 fois : \zh{儿子、女儿、儿子、女儿} (érzi, nǚ'ér, érzi, nǚ'ér).
\end{attention}

\courssub{Exercice résolu}

\begin{exoresolu}[Traduction, structure 虽然...但是... et place du temps]
Traduisez en chinois : « Autrefois, au Cameroun, la coutume n'autorisait pas les filles à hériter, mais maintenant la loi établit l'égalité. »
\tcblower
\textbf{Analyse :} Deux propositions opposées (coutume vs loi) $\rightarrow$ \zh{虽然...但是...}. Mots de temps : « autrefois » (\zh{以前}), « maintenant » (\zh{现在}). Sujets : « la coutume » (\zh{习惯法} ou \zh{习俗}), « la loi » (\zh{法律}). Verbes : « autoriser » (\zh{允许} / \zh{让}), « établir » (\zh{规定}).\\[0.3em]
\textbf{Construction :}\\
1. Proposition 1 : \zh{虽然以前习惯法不让女儿继承遗产} (Suīrán yǐqián xíguànfǎ bù ràng nǚ'ér jìchéng yíchǎn).\\
2. Proposition 2 : \zh{但是现在法律规定男女平等} (Dànshì xiànzài fǎlǜ guīdìng nán nǚ píngděng).\\
\textbf{Résultat :} \zh{虽然以前习惯法不让女儿继承遗产，但是现在法律规定男女平等。} (Suīrán yǐqián xíguànfǎ bù ràng nǚ'ér jìchéng yíchǎn, dànshì xiànzài fǎlǜ guīdìng nán nǚ píngděng.)
\end{exoresolu}

\courssub{Réflexion et mini-exposé : préparer sa prise de parole}

Cette section vise la compétence \textbf{Expression orale continue} (exposé de 2-3 minutes) et \textbf{Interaction} (débat). Vous devez mobiliser le vocabulaire, la grammaire (\zh{虽然...但是...}, \zh{以前/现在}, \zh{相同点/不同点}) et les faits socioculturels.

\begin{experience}[Activité orale : Débat structuré « Coutume vs Loi »]
\textbf{Consigne :} En binôme. L'un défend le point de vue « Tradition / Coutume », l'autre « Modernité / Loi ». Utilisez \textbf{au moins 5 mots du vocabulaire} et la structure \zh{虽然...但是...}.
\begin{itemize}
\item \textbf{Rôle A (Tradition) :} \zh{我认为习惯法很重要，因为它保护家族团结和祖先崇拜。虽然法律平等，但习俗是根。} (Wǒ rènwéi xíguànfǎ hěn zhòngyào, yīnwèi tā bǎohù jiāzú tuánjié hé zǔxiān chóngbài. Suīrán fǎlǜ píngděng, dàn xísú shì gēn.) « Je pense que le droit coutumier est important car il protège l'unité du clan et le culte des ancêtres. Bien que la loi soit égale, la coutume est la racine. »
\item \textbf{Rôle B (Modernité) :} \zh{我支持现代法，因为男女平等是人权。虽然传统有价值，但是法律必须公平。女儿也有继承权。} (Wǒ zhīchí xiàndài fǎ, yīnwèi nán nǚ píngděng shì rénquán. Suīrán chuántǒng yǒu jiàzhí, dànshì fǎlǜ bìxū gōngpíng. Nǚ'ér yě yǒu jìchéngquán.) « Je soutiens la loi moderne car l'égalité est un droit de l'homme. Bien que la tradition ait de la valeur, la loi doit être juste. Les filles ont aussi des droits. »
\end{itemize}
\textbf{Prolongement :} Échangez les rôles. L'enseignant note la prononciation des tons (surtout \zh{女儿}, \zh{平等}, \zh{继承}) et la fluidité des connecteurs.
\end{experience}

\begin{methode}[Préparer son mini-exposé écrit (150-200 caractères) — Évaluation]
\begin{enumerate}
\item \textbf{Planifier} : Remplissez un tableau « Cameroun / Chine » (Similitudes / Différences) en pinyin/caractères.
\item \textbf{Rédiger} : 4 paragraphes courts.
    \begin{enumerate}
    \item Intro : \zh{中喀两国继承制度比较} (Comparaison des systèmes successoraux Chine-Cameroun).
    \item Tradition : \zh{传统上，两国都重男轻女，实行父系继承。} (Traditionnellement, les deux valorisaient les garçons, succession patrilinéaire.)
    \item Droit moderne : \zh{现在，两国法律都规定男女平等，女儿有继承权。} (Maintenant, les deux lois établissent l'égalité, filles ont droits.)
    \item Différence clé : \zh{不同的是，喀麦隆习惯法与现代法并存，中国法律统一。} (Différence : Cameroun coutume+loi coexistent, Chine loi unifiée.)
    \item Conclusion personnelle : \zh{我觉得法律进步保护了妇女权益，但祭祖传统值得保留。} (Je pense que le progrès légal protège les droits des femmes, mais la tradition du culte des ancêtres mérite d'être gardée.)
    \end{enumerate}
\item \textbf{Relire} : Vérifier l'ordre des mots (Sujet + Temps + Verbe), les tons (\zh{平等}, \zh{女儿}), la ponctuation chinoise 。
\item \textbf{S'enregistrer} : Lire à voix haute pour l'oral.
\end{enumerate}
\end{methode}

\begin{exercice}[Entraînement complet : Écrit, Traduction, Oral]
\begin{enumerate}
\item \textbf{Traduction (Version) :} Traduisez en français : \zh{虽然中国以前只有儿子继承遗产，但是现在法律规定儿子女儿平等，妻子也可以继承丈夫的遗产。}
\item \textbf{Traduction (Thème) :} Traduisez en chinois : « Au Cameroun, la coutume et la loi moderne coexistent. Autrefois, seuls les fils héritaient de la terre. Maintenant, les filles ont aussi le droit d'hériter. »
\item \textbf{Expression écrite :} Répondez en chinois (phrase complète) : \zh{你觉得“祭祖”这个传统重要吗？为什么？} (Pensez-vous que la tradition du « culte des ancêtres » est importante ? Pourquoi ?)
\item \textbf{Production orale (à préparer chez soi) :} Préparez un exposé de 2 minutes : « L'héritage immatériel chez moi et en Chine : le nom de famille, les généalogies, les rites ». Utilisez \zh{家谱}, \zh{祭祖}, \zh{清明节}, \zh{口述历史}, \zh{传承} (chuánchéng, transmettre).
\end{enumerate}
\end{exercice}

\begin{corrige}[Exercice]
\begin{enumerate}
\item « Bien qu'autrefois en Chine seuls les fils héritassent des biens, maintenant la loi établit l'égalité entre fils et filles, et l'épouse peut aussi hériter des biens de son mari. »
\item \zh{喀麦隆习惯法和现代法并存。以前只有儿子继承土地，现在女儿也有继承权。} (Kāmàilóng xíguànfǎ hé xiàndài fǎ bìngcún. Yǐqián zhǐyǒu érzi jìchéng tǔdì, xiànzài nǚ'ér yě yǒu jìchéngquán.) — \textbf{Attention} : \zh{并存} (coexister), \zh{土地} (terre/terrain), \zh{继承权} (droit de succession). Place de \zh{以前}/\zh{现在} après le sujet.
\item Réponse modèle : \zh{我觉得“祭祖”很重要，因为它让我们记住祖先，传承家族文化，增强家庭凝聚力。} (Wǒ juéde « jìzǔ » hěn zhòngyào, yīnwèi tā ràng wǒmen jìzhù zǔxiān, chuánchéng jiāzú wénhuà, zēngqiáng jiātíng níngjùlì.) « Je trouve le culte des ancêtres important car il nous fait souvenir des ancêtres, transmettre la culture familiale, renforcer la cohésion familiale. »
\item (Production libre — évaluation par l'enseignant sur : vocabulaire du thème, structures \zh{虽然...但是...} / \zh{以前/现在}, prononciation tons, cohérence culturelle).
\end{enumerate}
\end{corrige}

\begin{aretenir}[L'essentiel de la leçon 4 — Fiche de révision]
\begin{itemize}
\item \textbf{Concepts} : Héritage (\zh{遗产}), Succession (\zh{继承}), Coutume (\zh{习惯法}), Loi (\zh{法律}/ \zh{民法典}), Testateur (\zh{立遗嘱人}), Conjoint (\zh{配偶者}).
\item \textbf{Cameroun} : Bijuridisme (coutume + loi). Diversité ethnique (patrilinéaire / matrilinéaire). Égalité légale (Code civil) vs pratique coutumière.
\item \textbf{Chine} : Tradition confucéenne (primogéniture mâle, \zh{家谱}, \zh{祭祖}, \zh{清明节}). Loi 1985 / Code civil 2021 : égalité stricte, testament moderne.
\item \textbf{Comparaison} : Même évolution « tradition inégalitaire $\rightarrow$ loi égalitaire ». Différence : pluralisme juridique (Cameroun) vs unification (Chine). Héritage immatériel vivant dans les deux (ancêtres).
\item \textbf{Grammaire} : \zh{以前} / \zh{现在} (Sujet + Temps + Verbe). \zh{虽然...但是...} (Concession). \zh{相同点} / \zh{不同点} (Comparaison).
\item \textbf{Caractères à savoir écrire} : \zh{法, 平, 继, 遗, 祖, 谱, 祭, 明}.
\end{itemize}
\end{aretenir}

\coursec{Comparaison Cameroun / Chine}

Nous venons d'étudier le processus d'hérédité en Chine : une tradition longtemps patrilinéaire, un cadre juridique unique et moderne qui garantit l'égalité entre fils et filles. Nous avons aussi vu, dans les sections précédentes, la situation camerounaise, marquée par une grande diversité coutumière. Le moment est venu de \cle{mettre les deux pays en regard} : qu'ont-ils en commun ? En quoi diffèrent-ils ? Et surtout : comment dire ces comparaisons \emph{en chinois}, avec des structures correctes et une attitude respectueuse ?

\courssub{Observer d'abord : un texte de comparaison}

\begin{textebox}[Deux pays, deux traditions]
喀麦隆和中国都重视家庭，但是两个国家的继承习惯不一样。\\
Kāmàilóng hé Zhōngguó dōu zhòngshì jiātíng, dànshì liǎng ge guójiā de jìchéng xíguàn bù yíyàng.\\
« Le Cameroun et la Chine accordent tous deux de l'importance à la famille, mais les coutumes d'héritage des deux pays sont différentes. »\\[4pt]
在中国，传统上儿子继承父亲的财产；现在在法律上，儿子和女儿有一样的权利。\\
Zài Zhōngguó, chuántǒng shàng érzi jìchéng fùqin de cáichǎn; xiànzài zài fǎlǜ shàng, érzi hé nǚ'ér yǒu yíyàng de quánlì.\\
« En Chine, traditionnellement, le fils héritait des biens du père ; aujourd'hui, devant la loi, fils et filles ont les mêmes droits. »\\[4pt]
在喀麦隆，有的民族是父系的，有的民族是母系的，所以习惯比中国多。现代法律说女儿和儿子一样，但是很多家庭还按照习惯做。\\
Zài Kāmàilóng, yǒude mínzú shì fùxì de, yǒude mínzú shì mǔxì de, suǒyǐ xíguàn bǐ Zhōngguó duō. Xiàndài fǎlǜ shuō nǚ'ér hé érzi yíyàng, dànshì hěn duō jiātíng hái ànzhào xíguàn zuò.\\
« Au Cameroun, certains peuples sont patrilinéaires, d'autres matrilinéaires, donc les coutumes sont plus nombreuses qu'en Chine. La loi moderne dit que la fille est égale au fils, mais beaucoup de familles agissent encore selon la coutume. »
\end{textebox}

\textbf{Glossaire du texte}\par
\begin{center}
\begin{tabularx}{\linewidth}{|X|X|X|X|}
\hline
\rowcolor{popblueL} Caractères & Pinyin & Nature & Traduction \\
\hline
重视 & zhòngshì & verbe & accorder de l'importance à \\
\hline
继承 & jìchéng & verbe & hériter \\
\hline
习惯 & xíguàn & nom & coutume, habitude \\
\hline
财产 & cáichǎn & nom & biens, patrimoine \\
\hline
权利 & quánlì & nom & droit (d'une personne) \\
\hline
父系 & fùxì & nom/adj. & patrilinéaire \\
\hline
母系 & mǔxì & nom/adj. & matrilinéaire \\
\hline
按照 & ànzhào & préposition & selon, conformément à \\
\hline
民族 & mínzú & nom & peuple, ethnie, nationalité \\
\hline
传统 & chuántǒng & nom/adj. & tradition, traditionnel \\
\hline
现代 & xiàndài & adj. & moderne \\
\hline
\end{tabularx}
\end{center}

\textbf{Questions de compréhension}
\begin{enumerate}
\item 两个国家都重视什么？(Liǎng ge guójiā dōu zhòngshì shénme ? — À quoi les deux pays accordent-ils de l'importance ?)
\item 在中国，现在儿子和女儿有一样的权利吗？(Zài Zhōngguó, xiànzài érzi hé nǚ'ér yǒu yíyàng de quánlì ma ? — En Chine, fils et filles ont-ils aujourd'hui les mêmes droits ?)
\item 喀麦隆的习惯为什么比中国多？(Kāmàilóng de xíguàn wèishénme bǐ Zhōngguó duō ? — Pourquoi les coutumes camerounaises sont-elles plus nombreuses que celles de la Chine ?)
\end{enumerate}

\courssub{Points communs et différences : les faits}

\begin{definition}[Points communs]
Dans les \textbf{deux pays} : la \cle{famille} occupe une place centrale dans la société ; historiquement, les \cle{fils} jouaient un rôle privilégié dans la transmission des biens et du nom ; aujourd'hui, le droit moderne évolue vers l'\cle{égalité} entre filles et fils.
\end{definition}

\begin{definition}[Différences principales]
\begin{itemize}
\item \textbf{Diversité coutumière} : le Cameroun connaît des systèmes \cle{patrilinéaires} \emph{et} \cle{matrilinéaires} (chez certains peuples, l'héritage passe par l'oncle maternel) ; la tradition chinoise est essentiellement \cle{patrilinéaire}.
\item \textbf{Cadre juridique} : la Chine applique un \cle{droit unique} (une seule loi nationale sur la succession) ; le Cameroun connaît un \cle{pluralisme juridique} : la coutume et le droit moderne coexistent, parfois en tension.
\end{itemize}
\end{definition}

\begin{savaistu}[Le saviez-vous ?]
En Chine, la loi sur la succession (\zh{继承法}, jìchéngfǎ) affirme clairement l'égalité des filles et des fils depuis des décennies. Au Cameroun, la Constitution affirme aussi l'égalité de tous les citoyens, mais dans beaucoup de villages, c'est encore la coutume locale qui décide de l'héritage. Ce décalage entre \zh{法律} (fǎlǜ, la loi) et \zh{习惯} (xíguàn, la coutume) est un sujet de débat dans les deux sociétés.
\end{savaistu}

\courssub{Dire la comparaison en chinois : les structures}

Observons les phrases du texte : \zh{两个国家都重视家庭} (les deux pays accordent de l'importance à la famille), \zh{习惯比中国多} (les coutumes sont plus nombreuses qu'en Chine), \zh{习惯不一样} (les coutumes sont différentes), \zh{但是…} (mais…). Quatre outils suffisent pour comparer : \cle{都}, \cle{比}, \cle{和……一样 / 不一样}, \cle{但是}.

\begin{propriete}[La comparaison avec 比]
\textbf{Structure :} A + 比 + B + adjectif.\\
\zh{中国的家庭比喀麦隆的家庭小。} \emph{Zhōngguó de jiātíng bǐ Kāmàilóng de jiātíng xiǎo.} « Les familles chinoises sont plus petites que les familles camerounaises. »\\
Emploi : \zh{比} introduit le \emph{terme de comparaison} (ce à quoi on compare). L'adjectif seul suffit : pas besoin de « plus ».\\
\textbf{Attention :} on ne dit jamais \zh{很} devant l'adjectif dans une phrase avec \zh{比} : \emph{pas} \zh{比…很小}, mais \zh{比…小}. Pour insister sur l'écart, on ajoute \zh{多} (un peu plus) ou \zh{得多} (beaucoup plus) \emph{après} l'adjectif : \zh{习惯比中国多得多} (les coutumes sont \emph{beaucoup} plus nombreuses qu'en Chine).\\
Négation : A + 没有 + B + adjectif (A n'est pas aussi… que B).\\
\zh{中国的习惯没有喀麦隆的习惯多。} \emph{Zhōngguó de xíguàn méiyǒu Kāmàilóng de xíguàn duō.} « Les coutumes chinoises ne sont pas aussi nombreuses que les coutumes camerounaises. »
\end{propriete}

\begin{propriete}[L'égalité et la différence : 和……一样 / 不一样]
\textbf{Égalité :} A + 和 + B + 一样 (+ adjectif / verbe).\\
\zh{女儿和儿子一样。} \emph{Nǚ'ér hé érzi yíyàng.} « La fille est égale au fils. »\\
\zh{女儿和儿子有一样的权利。} \emph{Nǚ'ér hé érzi yǒu yíyàng de quánlì.} « La fille et le fils ont les mêmes droits. »\\
\textbf{Différence :} A + 和 + B + 不一样, ou sujet pluriel + 不一样.\\
\zh{两个国家的习惯不一样。} \emph{Liǎng ge guójiā de xíguàn bù yíyàng.} « Les coutumes des deux pays sont différentes. »\\
\textbf{Point commun (ensemble) :} sujet pluriel + 都 + verbe.\\
\zh{两个国家都重视家庭。} \emph{Liǎng ge guójiā dōu zhòngshì jiātíng.} « Les deux pays accordent de l'importance à la famille. »\\
\textbf{Opposition :} 但是 + phrase.\\
\zh{但是习惯不一样。} \emph{Dànshì xíguàn bù yíyàng.} « Mais les coutumes sont différentes. »
\end{propriete}

\begin{exemplebox}[Exemples traduits]
\begin{itemize}
\item \zh{喀麦隆和中国都重视家庭。} \emph{Kāmàilóng hé Zhōngguó dōu zhòngshì jiātíng.} « Le Cameroun et la Chine accordent tous deux de l'importance à la famille. »
\item \zh{喀麦隆的习惯比中国的习惯多。} \emph{Kāmàilóng de xíguàn bǐ Zhōngguó de xíguàn duō.} « Les coutumes camerounaises sont plus nombreuses que les coutumes chinoises. »
\item \zh{在中国，女儿和儿子有一样的权利。} \emph{Zài Zhōngguó, nǚ'ér hé érzi yǒu yíyàng de quánlì.} « En Chine, la fille et le fils ont les mêmes droits. »
\item \zh{两个国家都很重视家庭，但是继承的习惯不一样。} \emph{Liǎng ge guójiā dōu hěn zhòngshì jiātíng, dànshì jìchéng de xíguàn bù yíyàng.} « Les deux pays attachent beaucoup d'importance à la famille, mais les coutumes d'héritage sont différentes. »
\end{itemize}
\end{exemplebox}

\begin{center}
\begin{tabularx}{\linewidth}{|X|X|X|}
\hline
\rowcolor{popblueL} Structure & Exemple (caractères + pinyin) & Traduction \\
\hline
A 和 B 都 + verbe & 两个国家都重视家庭。Liǎng ge guójiā dōu zhòngshì jiātíng. & Les deux pays accordent de l'importance à la famille. \\
\hline
A 比 B + adjectif & 中国的家庭比喀麦隆的家庭小。Zhōngguó de jiātíng bǐ Kāmàilóng de jiātíng xiǎo. & Les familles chinoises sont plus petites que les familles camerounaises. \\
\hline
A 和 B 一样 & 女儿和儿子一样。Nǚ'ér hé érzi yíyàng. & La fille est égale au fils. \\
\hline
A 和 B 不一样 & 两个国家的习惯不一样。Liǎng ge guójiā de xíguàn bù yíyàng. & Les coutumes des deux pays sont différentes. \\
\hline
…，但是… & 都重视家庭，但是习惯不一样。Dōu zhòngshì jiātíng, dànshì xíguàn bù yíyàng. & Tous deux valorisent la famille, mais les coutumes sont différentes. \\
\hline
\end{tabularx}
\end{center}

\begin{attention}[Pièges fréquents des francophones]
\begin{itemize}
\item \textbf{两个, pas 二个} : devant un classificateur, « deux » se dit \zh{两} (liǎng) : \zh{两个国家} (deux pays). \zh{二} (èr) sert à compter (1, 2, 3…) ou dans les nombres composés (12, 20…), jamais directement devant un classificateur.
\item \textbf{Pas de \zh{很} avec \zh{比}} : on dit \zh{习惯比中国多}, jamais \zh{比中国很多}.
\item \textbf{Ordre des mots} : le français dit « plus petit que » (adjectif + que) ; le chinois place le terme de comparaison \emph{avant} l'adjectif : A 比 B + adjectif. Ne traduisez pas mot à mot.
\item \textbf{\zh{一样} ne prend jamais \zh{很}} : on dit \zh{一样} ou \zh{不一样}, point final. Pas de \zh{很一样}.
\item \textbf{Classificateur de \zh{法律}} : à l'écrit formel, on préfère \zh{一部法律} (yí bù fǎlǜ) ; \zh{一个法律} reste acceptable à l'oral.
\end{itemize}
\end{attention}

\courssub{Écriture des caractères : vocabulaire clé}

\begin{definition}[Analyse graphique et ordre des traits]
Les caractères ci-dessous sont des mots-outils ou du vocabulaire spécifique à la leçon. Respectez l'ordre des traits (haut → bas, gauche → droite, horizontal avant vertical, trait central avant les ailes) et identifiez la clé (radical) pour les retrouver dans le dictionnaire.
\end{definition}

\begin{center}
\begin{tabularx}{\linewidth}{|X|X|X|X|X|}
\hline
\rowcolor{popblueL} Caractère & Pinyin & Clé (Radical) & Nb traits & Ordre des traits (résumé) \\
\hline
继 & jì & 纟 (sī, soie) & 10 & 纟 (3) + 米 (6) + 页 (6) → écrire la clé soie, puis 米, puis 页 \\
\hline
承 & chéng & 手 (shǒu, main) & 8 & 手 (3) + 丰 (4) + 一 (1) → main, puis 丰, trait horizontal final \\
\hline
习 & xí & 羽 (yǔ, plume) & 11 & 羽 (6) + 白 (5) → plume (haut), puis blanc (bas) \\
\hline
惯 & guàn & 忄 (xīn, cœur) & 11 & 忄 (3) + 贯 (8) → clé cœur à gauche, puis 贯 (贝 + 毌) \\
\hline
父 & fù & 父 (fù, père) & 4 & 两横 (courtes) + 一长横 + 两捺 (pieds) \\
\hline
系 & xì & 糸 (mì, fil) & 7 & 糸 (6) + 页 (6) → fil (gauche), puis page (droite) \\
\hline
母 & mǔ & 母 (mǔ, mère) & 5 & 两点 + 横折钩 + 两点 (seins) — diffère de \zh{毋} (wú) \\
\hline
法 & fǎ & 氵 (shuǐ, eau) & 8 & 氵 (3) + 去 (5) → eau, puis 去 (上土厶) \\
\hline
律 & lǜ & 彳 (chì, pas) & 9 & 彳 (3) + 聿 (6) → pas, puis pinceau (聿) \\
\hline
权 & quán & 木 (mù, bois) & 6 & 木 (4) + 又 (2) → arbre + main (droit, pouvoir) \\
\hline
利 & lì & 刂 (dāo, couteau) & 7 & 禾 (5) + 刂 (2) → grain, puis couteau à droite \\
\hline
\end{tabularx}
\end{center}

\begin{experience}[Atelier écriture : « Mots croisés » des radicaux]
En binôme : l'un dicte le pinyin (ex. \emph{jìchéng}), l'autre écrit les deux caractères sur ardoise en annonçant la clé et le nombre de traits à haute voix. Inversez les rôles. Vérifiez l'ordre des traits sur \zh{继} (clé soie) et \zh{系} (clé fil) : même clé \zh{糸} mais position différente (bas vs gauche).
\end{experience}

\courssub{Comparer sans juger : la bonne attitude}

\begin{methode}[Comparer deux cultures avec respect]
\begin{enumerate}
\item \textbf{Énoncer le fait} avec une structure neutre : \zh{A 和 B 不一样} (A et B sont différents).
\item \textbf{Préférer} \zh{不一样} (bù yíyàng, « différent ») à \zh{不好} (bù hǎo, « mauvais ») ou \zh{更好} (gèng hǎo, « meilleur ») : on décrit, on ne juge pas.
\item \textbf{Expliquer la raison} avec \zh{因为…所以…} : \zh{因为喀麦隆有很多民族，所以习惯很多。} \emph{Yīnwèi Kāmàilóng yǒu hěn duō mínzú, suǒyǐ xíguàn hěn duō.} « Parce que le Cameroun compte de nombreux peuples, les coutumes sont nombreuses. »
\item \textbf{Conclure sur le point commun} avec \zh{都} : \zh{两个国家都重视家庭。} Cela termine la comparaison sur une note positive et factuelle.
\end{enumerate}
\end{methode}

\courssub{Tableau comparatif à compléter en classe}

\begin{popfigure}
\begin{tikzpicture}[
    cell/.style={draw, minimum height=1cm, align=center, font=\small},
    header/.style={cell, fill=popblueL, draw=popblue, font=\bfseries},
    cam/.style={cell, draw=popgreen, fill=popgreenL},
    chn/.style={cell, draw=poporange, fill=poporangeL},
    label/.style={cell, fill=popblueL, draw=popblue, font=\bfseries}
]
% Grille 4 lignes x 3 colonnes
% Ligne 1 : En-têtes
\node[label] (c11) at (2, 0.5) {Critères / 标准};
\node[header] (c12) at (6, 0.5) {\textbf{喀麦隆} Kāmàilóng};
\node[header] (c13) at (10, 0.5) {\textbf{中国} Zhōngguó};

% Ligne 2 : Tradition
\node[label] (c21) at (2, -0.5) {传统 / chuántǒng};
\node[cam] (c22) at (6, -0.5) {?};
\node[chn] (c23) at (10, -0.5) {?};

% Ligne 3 : Loi actuelle
\node[label] (c31) at (2, -1.5) {法律 / fǎlǜ};
\node[cam] (c32) at (6, -1.5) {?};
\node[chn] (c33) at (10, -1.5) {?};

% Ligne 4 : Droits des filles
\node[label] (c41) at (2, -2.5) {女儿的权利 / nǚ'ér de quánlì};
\node[cam] (c42) at (6, -2.5) {?};
\node[chn] (c43) at (10, -2.5) {?};

% Lignes de grille (redessinées par les styles node, mais on s'assure des bordures)
\draw[popline] (0,1) rectangle (12,-3);
\draw[popline] (4,1) -- (4,-3);
\draw[popline] (8,1) -- (8,-3);
\draw[popline] (0,0) -- (12,0);
\draw[popline] (0,-1) -- (12,-1);
\draw[popline] (0,-2) -- (12,-2);
\end{tikzpicture}
\legende{Tableau comparatif à compléter : pour chaque ligne (tradition, loi actuelle, droits des filles), rédigez une phrase courte en chinois pour chaque pays, en utilisant \zh{比}, \zh{一样} ou \zh{不一样}. Exemple ligne 1 : Cameroun \zh{有的父系，有的母系} ; Chine \zh{传统是父系的} ; Comparaison \zh{喀麦隆的传统和中国的传统不一样。}}
\end{popfigure}

\courssub{Expression orale : Mini-exposé guidé}

\begin{experience}[Débat en binôme : « Deux regards sur l'héritage »]
\begin{enumerate}
\item \textbf{Préparation (5 min)} : Chaque élève choisit un pays (Cameroun ou Chine). Il prépare 3 phrases en chinois (caractères + pinyin) : une sur la tradition, une sur la loi, une sur les droits des filles. Il utilise le vocabulaire de la leçon (\zh{父系}, \zh{母系}, \zh{法律}, \zh{习惯}, \zh{权利}, \zh{平等}).
\item \textbf{Présentation croisée (3 min)} : L'élève A présente le Cameroun, l'élève B la Chine. Ils s'écoutent sans interrompre.
\item \textbf{Synthèse commune (2 min)} : Ensemble, ils formulent \textbf{deux phrases finales} en chinois :
      \begin{itemize}
      \item Une phrase de différence avec \zh{不一样} ou \zh{比}.
      \item Une phrase de point commun avec \zh{都}.
      \end{itemize}
      Exemple : \zh{喀麦隆的习惯比中国的多，但是两个国家都重视家庭平等。}
\item \textbf{Restitution} : Quelques binômes présentent leur synthèse à la classe. Le professeur corrige la prononciation (tons, \zh{ü}, \zh{zh/ch/sh/r}) et la structure.
\end{enumerate}
\end{experience}

\begin{exoresolu}[Construire une phrase de comparaison]
Énoncé : traduisez en chinois « Les coutumes camerounaises sont plus nombreuses que les coutumes chinoises, mais les deux pays valorisent la famille. »
\tcblower
Solution détaillée :
\begin{enumerate}
\item « Plus nombreuses que » appelle la structure A + 比 + B + adjectif : \zh{喀麦隆的习惯比中国的习惯多。} \emph{Kāmàilóng de xíguàn bǐ Zhōngguó de xíguàn duō.}
\item « Mais » se dit \zh{但是} (dànshì).
\item « Les deux pays valorisent la famille » appelle \zh{都} : \zh{两个国家都重视家庭。} — attention : « deux pays » = \zh{两个国家}, avec \zh{两} et le classificateur \zh{个}.
\item Phrase complète : \zh{喀麦隆的习惯比中国的习惯多，但是两个国家都重视家庭。} \emph{Kāmàilóng de xíguàn bǐ Zhōngguó de xíguàn duō, dànshì liǎng ge guójiā dōu zhòngshì jiātíng.}
\end{enumerate}
\end{exoresolu}

\begin{exercice}[À vous de comparer]
\begin{enumerate}
\item Traduisez : « En Chine, il n'y a qu'une seule loi ; au Cameroun, il y a la loi et la coutume. » (utilisez \zh{只有} zhǐyǒu, « il n'y a que », et \zh{和}).
\item Complétez avec \zh{比}, \zh{一样} ou \zh{不一样} : 中国的传统 \rule{1.5cm}{0.4pt} 喀麦隆的传统。
\item Complétez le tableau de la figure ci-dessus : une phrase par case (colonnes Cameroun et Chine), en chinois, avec pinyin.
\item Mini-débat en binôme (voir activité « Expression orale ») : l'un présente le Cameroun, l'autre la Chine ; chacun doit finir par une phrase avec \zh{都} qui montre le point commun.
\end{enumerate}
\end{exercice}

\begin{corrige}[Exercice « À vous de comparer »]
\begin{enumerate}
\item \zh{在中国，只有一部法律；在喀麦隆，有法律，也有习惯。} \emph{Zài Zhōngguó, zhǐyǒu yí bù fǎlǜ; zài Kāmàilóng, yǒu fǎlǜ, yě yǒu xíguàn.} (Classificateur \zh{部} pour la loi ; \zh{一个} accepté à l'oral. \zh{一种法律} possible pour « un type de loi ».)
\item \zh{中国的传统和喀麦隆的传统不一样。} \emph{Zhōngguó de chuántǒng hé Kāmàilóng de chuántǒng bù yíyàng.} « Les traditions chinoises et camerounaises sont différentes. » (\zh{一样} possible si l'élève justifie un point commun précis, mais \zh{不一样} est la réponse attendue pour la différence globale).
\item \textbf{Proposition de correction du tableau} :\\
Ligne « Tradition » : Cameroun — \zh{有的民族是父系的，有的是母系的。}(Yǒude mínzú shì fùxì de, yǒude shì mǔxì de.) ; Chine — \zh{传统上是父系的。}(Chuántǒng shàng shì fùxì de.)\\
Ligne « Loi » : Cameroun — \zh{有现代法律，也有习惯法。}(Yǒu xiàndài fǎlǜ, yě yǒu xíguànfǎ.) ; Chine — \zh{只有一部继承法。}(Zhǐyǒu yí bù jìchéngfǎ.)\\
Ligne « Droits des filles » : Cameroun — \zh{法律上一样，习惯上不一样。}(Fǎlǜ shàng yíyàng, xíguàn shàng bù yíyàng.) ; Chine — \zh{法律上儿子女儿一样。}(Fǎlǜ shàng érzi nǚ'ér yíyàng.)\\
\textbf{Phrase de comparaison globale} : \zh{喀麦隆的习惯比中国的多，但是两个国家的法律都给女儿权利。}
\item Phrase de conclusion attendue : \zh{两个国家都重视家庭，也都给女儿权利。} \emph{Liǎng ge guójiā dōu zhòngshì jiātíng, yě dōu gěi nǚ'ér quánlì.} « Les deux pays accordent de l'importance à la famille et reconnaissent aussi des droits aux filles. »
\end{enumerate}
\end{corrige}

\begin{aretenir}[L'essentiel de la section]
\begin{itemize}
\item \textbf{Faits} : mêmes valeurs familiales et même évolution vers l'égalité, mais diversité coutumière (patrilinéaire + matrilinéaire) et pluralisme juridique au Cameroun ; tradition patrilinéaire unifiée et droit unique en Chine.
\item \textbf{Structures} :
      \begin{itemize}
      \item Supériorité : A + 比 + B + adjectif (+ \zh{多}/\zh{得多}) ; Négation : A + 没有 + B + adjectif.
      \item Égalité / Différence : A + 和 + B + 一样 / 不一样.
      \item Point commun : Sujet pluriel + 都 + Verbe.
      \item Opposition : 但是 + phrase.
      \item Causalité : 因为…所以…
      \end{itemize}
\item \textbf{Vocabulaire clé} : \zh{继承} jìchéng, \zh{习惯} xíguàn, \zh{法律} fǎlǜ, \zh{权利} quánlì, \zh{父系} fùxì, \zh{母系} mǔxì, \zh{平等} píngděng, \zh{民族} mínzú.
\item \textbf{Attitude} : on compare avec \zh{不一样} (différent), jamais avec un jugement de valeur (\zh{不好}, \zh{更好}) ; on dit \zh{两个国家}, jamais \zh{二个国家} ; pas de \zh{很} avec \zh{比} ni avec \zh{一样}.
\item \textbf{Écriture} : maîtriser l'ordre des traits et les clés de \zh{继承习惯父系母系法律权利} (clés : 纟, 羽, 忄, 父, 糸, 母, 氵, 彳, 木, 刂).
\end{itemize}
\end{aretenir}

\coursec{Vocabulaire et écriture des caractères}

Après avoir comparé les pratiques successorales au Cameroun et en Chine, nous devons maintenant maîtriser le lexique précis qui permet d'en parler en chinois. Cette section structure le vocabulaire thématique, analyse la composition des caractères clés pour en faciliter la mémorisation et l'écriture, et vous entraîne à éviter les confusions fréquentes.

\courssub{Glossaire thématique : l'hérédité et la famille}

Le tableau suivant regroupe les termes indispensables pour décrire le processus d'hérédité, les acteurs et le cadre légal ou coutumier. Observez la formation des mots composés (ex. \zh{继承人} = hériter + personne) et notez la nature grammaticale de chaque entrée.

\begin{center}
\begin{tabularx}{\linewidth}{|X|X|X|X|}
\hline
\rowcolor{popblueL} \textbf{Caractères} & \textbf{Pinyin} & \textbf{Nature} & \textbf{Traduction} \\
\hline
遗产 & yíchǎn & n. & héritage, succession, biens laissés par un défunt \\
\hline
继承 & jìchéng & v. & hériter, succéder à ; continuer (une tradition) \\
\hline
继承人 & jìchéngrén & n. & héritier, successeur \\
\hline
遗嘱 & yízhǔ & n. & testament, dernières volontés \\
\hline
法律 & fǎlǜ & n. & loi, législation \\
\hline
平等 & píngděng & adj. & égal, égalitaire \\
\hline
习俗 & xísú & n. & coutume, mœurs, tradition sociale \\
\hline
酋长 & qiúzhǎng & n. & chef (coutumier, tribal) \\
\hline
长子 & zhǎngzi & n. & fils aîné \\
\hline
寡妇 & guǎfu & n. & veuve \\
\hline
祖先 & zǔxiān & n. & ancêtres, aïeux \\
\hline
家谱 & jiāpǔ & n. & généalogie, arbre généalogique, registre de famille \\
\hline
\end{tabularx}
\end{center}

\begin{textebox}[Texte support : La succession dans deux contextes]
\zh{在我喀麦隆的老家，酋长管理土地继承。按照习俗，长子往往得到主要遗产，但寡妇也有居住权。法律现在规定男女平等，保护寡妇权利。} \\
\zh{Wǒ zài Kāmàilóng de lǎojiā, qiúzhǎng guǎnlǐ tǔdì jìchéng. Ànzhào xísú, zhǎngzi wǎngwǎng dédào zhǔyào yíchǎn, dàn guǎfu yě yǒu jūzhùquán. Fǎlǜ xiànzài guīdìng nánnǚ píngděng, bǎohù guǎfu quánlì.} \\
« Dans mon village natal au Cameroun, le chef coutumier gère l'héritage de la terre. Selon la coutume, le fils aîné reçoit souvent l'héritage principal, mais la veuve a aussi un droit de résidence. La loi stipule désormais l'égalité hommes-femmes et protège les droits des veuves. » \\

\zh{中国古代，家谱记录家族血脉，长子继承制很普遍。现代法律规定：遗产由第一顺序继承人继承，包括配偶、子女、父母，男女平等。立遗嘱可以自由分配财产。} \\
\zh{Zhōngguó gǔdài, jiāpǔ jìlù jiāzú xuémài, zhǎngzi jìchéngzhì hěn pǔbiàn. Xiàndài fǎlǜ guīdìng: yíchǎn yóu dì yī shùnxù jìchéngrén jìchéng, bāokuò pèi'ǒu, zǐnǚ, fùmǔ, nánnǚ píngděng. Lì yízhǔ kěyǐ zìyóu fēnpèi cáichǎn.} \\
« Dans la Chine ancienne, la généalogie enregistrait la lignée du clan, et le système d'héritage par le fils aîné était très répandu. La loi moderne stipule : l'héritage revient aux héritiers du premier ordre (conjoint, enfants, parents), hommes et femmes égaux. Faire un testament permet de répartir librement les biens. »
\end{textebox}

\begin{exercice}[Compréhension du texte]
\begin{enumerate}
\item Selon le texte, qui gère l'héritage de la terre au village camerounais ?
\item Quels sont les héritiers du premier ordre selon la loi chinoise moderne citée ?
\item Quel terme chinois désigne « l'égalité hommes-femmes » et où apparaît-il dans les deux paragraphes ?
\end{enumerate}
\end{exercice}

\begin{corrige}[Exercice Compréhension du texte]
\begin{enumerate}
\item C'est le \zh{酋长} (chef coutumier).
\item \zh{配偶、子女、父母} (conjoint, enfants, parents).
\item \zh{男女平等} (nánnǚ píngděng) : apparaît à la fin du premier paragraphe (\zh{法律现在规定男女平等}) et à la fin du second (\zh{男女平等}).
\end{enumerate}
\end{corrige}

\courssub{Analyse morphologique et clés d'écriture (Radicaux)}

La plupart des caractères de cette leçon sont des caractères composés \textbf{phono-sémantiques} (形声字) : une clé sémantique (radical) indique le domaine de sens, et une composante phonétique suggère la prononciation. Reconnaître ces clés permet de deviner le sens et de classer les caractères dans le dictionnaire.

\begin{definition}[Clés sémantiques (Radicaux) de la leçon]
\begin{itemize}
\item \textbf{\zh{辶} (chuò) — « marche, cheminement »} : présente dans \zh{遗} (yí). Elle évoque l'idée de transmission, de passage d'une génération à l'autre (le bien « chemine » vers l'héritier).
\item \textbf{\zh{纟} (sī) — « fil de soie, filer »} : présente dans \zh{继} (jì). Sens originel : raccorder des fils de soie coupés → prolonger, succéder, continuer.
\item \textbf{\zh{氵} (shuǐ) — « eau » (trois points d'eau)} : présente dans \zh{法} (fǎ). Dans l'ancienne Chine, la loi était comparée à l'eau : impartiale, qui trouve son niveau, qui lave les fautes.
\item \textbf{\zh{宀} (mián) — « toit, maison »} : présente dans \zh{寡} (guǎ) et \zh{家} (jiā). Indique le domaine domestique, familial.
\item \textbf{\zh{女} (nǚ) — « femme »} : présente dans \zh{妇} (fù). Composante sémantique évidente pour « femme mariée », « épouse ».
\end{itemize}
\end{definition}

\begin{methode}[Mémoriser \zh{遗} (yí) — « laisser en héritage, perdre »]
\begin{enumerate}
\item \textbf{Décomposer} : \zh{辶} (cheminer) + \zh{贵} (guì, précieux, noble).
\item \textbf{Imaginer} : Un objet \textbf{précieux} (\zh{贵}) qui \textbf{chemine} (\zh{辶}) hors de la maison, ou qui est transmis le long du chemin des générations.
\item \textbf{Ancrage} : \zh{遗产} (héritage) = le bien précieux transmis ; \zh{遗嘱} (testament) = les mots précieux laissés avant de partir.
\item \textbf{Écrire} : Ordre des traits de \zh{辶} : point, horizontal court, pli vertical (3 traits) ; puis \zh{贵} (9 traits). Total : 12 traits.
\end{enumerate}
\end{methode}

\begin{popfigure}
\begin{tikzpicture}[
    node distance=1.2cm and 0.8cm,
    box/.style={draw=popblue, thick, rounded corners=4pt, fill=popblueL, minimum width=2.2cm, minimum height=1.2cm, align=center, font=\sffamily},
    arrow/.style={-Stealth, thick, popdark},
    radical/.style={draw=poporange, thick, rounded corners=4pt, fill=poporangeL, minimum width=1.8cm, minimum height=1cm, align=center, font=\sffamily\small}
]
% Nœud central
\node[box, align=center] (ji){\Huge \zh{继}\\[2pt] jì \\ hériter, continuer};

% Composants
\node[radical, align=center] (si) [left=of ji, yshift=0.6cm]{\zh{纟}\\(sī) \\ fil de soie};
\node[radical, align=center] (mi) [left=of ji, yshift=-0.6cm]{\zh{米}\\(mǐ) \\ riz / grains};
\node[radical, align=center] (yi) [below=of ji]{\zh{乚}\\(yǐ) \\ crochet / raccorder};

% Flèches
\draw[arrow] (si) -- node[auto, swap, font=\small, text=popdark] {sémantique} (ji.west);
\draw[arrow] (mi) -- node[auto, font=\small, text=popdark] {phonétique (forme simplifiée)} (ji.west);
\draw[arrow] (yi) -- node[auto, font=\small, text=popdark] {structure} (ji.south);

% Annotations
\node[align=center, text=popmuted, font=\small] at (4.5, 0.6) {Clé \zh{纟} : idée de \\ continuation (filer)};
\node[align=center, text=popmuted, font=\small] at (4.5, -0.6) {Phonétique ancien : \zh{丐} (gài) \\ Forme simplifiée : \zh{米} + \zh{乚}};
\node[align=center, text=popmuted, font=\small] at (0, -2.2) {Sens : raccorder les fils \\ (continuer la lignée)};
\end{tikzpicture}
\legende{Décomposition morphologique du caractère \zh{继} (jì) : clé sémantique \zh{纟} + composants \zh{米} et \zh{乚} (phonétique simplifié de \zh{丐}).}
\end{popfigure}

\begin{propriete}[Ordre des traits : trois caractères à fréquence élevée]
Respecter l'ordre des traits (haut $\to$ bas, gauche $\to$ droite, horizontal $\to$ vertical, cadre avant contenu, fermer en dernier) garantit une écriture lisible, rapide et esthétique.

\begin{center}
\begin{tabularx}{\linewidth}{|X|X|X|}
\hline
\rowcolor{popgreenL} \textbf{\zh{承} (chéng) — 10 traits} & \textbf{\zh{法} (fǎ) — 8 traits} & \textbf{\zh{家} (jiā) — 10 traits} \\
\hline
1. \zh{一} (h) \newline 2. \zh{丨} (v) \newline 3. \zh{一} (h) \newline 4. \zh{丿} (piě) \newline 5. \zh{丶} (diǎn) \newline 6. \zh{一} (h) \newline 7. \zh{丨} (v) \newline 8. \zh{一} (h) \newline 9. \zh{丿} (piě) \newline 10. \zh{丶} (diǎn) & 1. \zh{丶} (diǎn) \newline 2. \zh{丶} (diǎn) \newline 3. \zh{一} (h) \newline 4. \zh{丿} (piě) \newline 5. \zh{一} (h) \newline 6. \zh{丨} (v) \newline 7. \zh{一} (h) \newline 8. \zh{一} (h) & 1. \zh{丶} (diǎn) \newline 2. \zh{一} (h) \newline 3. \zh{丿} (piě) \newline 4. \zh{一} (h) \newline 5. \zh{丨} (v) \newline 6. \zh{一} (h) \newline 7. \zh{丿} (piě) \newline 8. \zh{丶} (diǎn) \newline 9. \zh{一} (h) \newline 10. \zh{一} (h) \\
\hline
\multicolumn{3}{|p{\linewidth}|}{\small \textbf{Note :} \zh{承} = \zh{禾} (hé, riz) + \zh{丙} (bǐng) = 10 traits. \zh{法} = \zh{氵} (eau) + \zh{去} (qù, aller/ôter) = 8 traits. \zh{家} = \zh{宀} (toit) + \zh{豕} (shǐ, cochon) : « sous le toit, le cochon » = la richesse, la famille = 10 traits.}
\end{tabularx}
\end{center}
\end{propriete}

\courssub{Caractères proches : discrimination visuelle et auditive}

Les apprenants francophones confondent souvent des caractères partageant une composante graphique ou phonétique. Le tableau ci-dessous met en regard les paires à risque de cette leçon.

\begin{center}
\begin{tabularx}{\linewidth}{|X|X|X|X|}
\hline
\rowcolor{poppurpleL} \textbf{Caractère cible} & \textbf{Pinyin / Sens} & \textbf{Caractère perturbateur} & \textbf{Pinyin / Sens + Différence clé} \\
\hline
\zh{遗} & yí — hériter, perdre & \zh{贵} & guì — précieux, cher \\ 
& \textit{Clé \zh{辶} (marche)} & \textit{Clé \zh{贝} (coquillage/argent)} & \textbf{Astuce} : \zh{遗} a le « pas » \zh{辶}, \zh{贵} a les « mains » \zh{贝} qui tiennent le trésor. \\
\hline
\zh{继} & jì — continuer, hériter & \zh{断} & duàn — couper, interrompre \\
& \textit{Clé \zh{纟} (fil)} & \textit{Clé \zh{斤} (hache/poids)} & \textbf{Antonymes} : \zh{纟} file le fil (continuer) ; \zh{斤} la hache coupe (interrompre). \\
\hline
\zh{俗} & sú — coutume, vulgaire & \zh{浴} & yù — bain, se baigner \\
& \textit{Clé \zh{亻} (personne)} & \textit{Clé \zh{氵} (eau)} & \textbf{Astuce} : \zh{俗} concerne les \textbf{gens} (\zh{亻}) ; \zh{浴} concerne l'\textbf{eau} (\zh{氵}). Phonétique commun : \zh{谷} (gǔ). \\
\hline
\end{tabularx}
\end{center}

\begin{attention}[Piège majeur : \zh{长} — deux prononciations, deux sens]
Le caractère \zh{长} est un \textbf{caractère polyphone} (多音字). En Terminale, vous devez maîtriser ses deux lectures principales :

\begin{center}
\begin{tabularx}{\linewidth}{|c|X|X|X|}
\hline
\rowcolor{poporangeL} \textbf{Pinyin} & \textbf{Sens principal} & \textbf{Exemples de la leçon} & \textbf{Autres exemples HSK} \\
\hline
\zh{zhǎng} (3\textsuperscript{e} ton) & \textbf{Âgé, aîné, chef} (adjectif/nom) & \zh{长子} (zhǎngzi, fils aîné), \zh{酋长} (qiúzhǎng, chef), \zh{班长} (bānzhǎng, délégué de classe) & \zh{长辈} (zhǎngbèi, aîné), \zh{村长} (cūnzhǎng, chef de village) \\
\hline
\zh{cháng} (2\textsuperscript{e} ton) & \textbf{Long, longueur} (adj./nom) & \zh{长短} (chángduǎn, longueur), \zh{长江} (Chángjiāng, Yangzi) & \zh{长久} (chángjiǔ, durable), \zh{长度} (chángdù, longueur) \\
\hline
\end{tabularx}
\end{center}

\textbf{Règle mnémotechnique} : \zh{zhǎng} (ton descendant-ascendant) = quelqu'un qui \textbf{grandit} et devient \textbf{chef/aîné}. \zh{cháng} (ton montant) = une ligne qui \textbf{s'allonge}.
\end{attention}

\courssub{Collocations et structures utiles pour l'expression écrite}

Maîtriser ces « blocs de langue » (chunks) vous permettra de rédiger des phrases correctes et idiomatiques dès l'examen.

\begin{center}
\begin{tabularx}{\linewidth}{|X|X|X|}
\hline
\rowcolor{poppinkL} \textbf{Collocation (结构)} & \textbf{Exemple en caractères + pinyin} & \textbf{Traduction} \\
\hline
\zh{继承 + [objet]} & \zh{继承遗产} (jìchéng yíchǎn) & hériter d'un héritage \\
\hline
\zh{法律 + 规定 + [phrase]} & \zh{法律规定男女平等。} (Fǎlǜ guīdìng nánnǚ píngděng.) & La loi stipule l'égalité hommes-femmes. \\
\hline
\zh{尊重 + [objet abstrait]} & \zh{尊重习俗} (zūnzhòng xísú) & respecter les coutumes \\
\hline
\zh{男女平等} (sujet/objet/attribut) & \zh{我们主张男女平等。} (Wǒmen zhǔzhāng nánnǚ píngděng.) & Nous prônons l'égalité hommes-femmes. \\
\hline
\zh{立 + 遗嘱} & \zh{立遗嘱} (lì yízhǔ) & faire un testament / rédiger un testament \\
\hline
\zh{第一顺序 + 继承人} & \zh{第一顺序继承人} (dì yī shùnxù jìchéngrén) & héritiers du premier ordre (réservé) \\
\hline
\end{tabularx}
\end{center}

\begin{exemplebox}[Phrases modèles pour la rédaction]
\begin{enumerate}
\item \zh{喀麦隆的习俗过去倾向于长子继承遗产。} (Kāmàilóng de xísú guòqù qīngxiàng yú zhǎngzi jìchéng yíchǎn.) — Autrefois, la coutume camerounaise penchait vers l'héritage par le fils aîné.
\item \zh{中国现代法律保障寡妇和女儿的继承权。} (Zhōngguó xiàndài fǎlǜ bǎozhàng guǎfu hé nǚ'ér de jìchéngquán.) — La loi moderne chinoise garantit les droits d'héritage des veuves et des filles.
\item \zh{无论城市还是农村，都应尊重法律，实现男女平等。} (Wúlùn chéngshì háishì nóngcūn, dōu yīng zūnzhòng fǎlǜ, shíxiàn nánnǚ píngděng.) — Que ce soit en ville ou à la campagne, on doit respecter la loi et réaliser l'égalité hommes-femmes.
\item \zh{查阅家谱可以了解祖先的历史。} (Cháyuè jiāpǔ kěyǐ liǎojiě zǔxiān de lìshǐ.) — Consulter la généalogie permet de connaître l'histoire des ancêtres.
\end{enumerate}
\end{exemplebox}

\begin{definition}[Expression clé d'examen : \zh{男女平等} (nánnǚ píngděng)]
\textbf{Structure} : \zh{男} (homme) + \zh{女} (femme) + \zh{平等} (égal/égalité). \\
\textbf{Emploi} : Sujet, objet, attribut ou complément. Invariable. \\
\textbf{Contexte} : Politique, droit, sociologie, vie quotidienne. Figure au programme du Bac (thème « Droits et devoirs », « Société »). \\
\textbf{Variantes} : \zh{性别平等} (xìngbié píngděng, égalité des sexes — plus formel/juridique), \zh{男女同权} (nánnǚ tóngquán, mêmes droits pour hommes et femmes).
\end{definition}

\courssub{Exercice résolu : rédaction d'un paragraphe comparatif}

\begin{exoresolu}[Rédiger un paragraphe de 60-80 caractères chinois comparant l'héritage au Cameroun et en Chine]
\textbf{Consigne} : Utilisez au moins 5 mots du glossaire, la structure \zh{法律规定...} et le connecteur \zh{但是} (dànshì) ou \zh{相比之下} (xiāngbǐ zhīxià). Écrivez en caractères, pinyin et français.

\tcblower

\textbf{Proposition de réponse} :

\zh{喀麦隆和中国的继承习俗有异同。喀麦隆过去按习俗，长子继承主要遗产，寡妇地位低。但是现在法律规定男女平等，保护寡妇权利。中国古代也重长子，家谱记录血脉。现代法律规定第一顺序继承人平分遗产，男女平等。总之，两国法律都走向平等。} \\
\zh{Kāmàilóng hé Zhōngguó de jìchéng xísú yǒu yìtóng. Kāmàilóng guòqù àn xísú, zhǎngzi jìchéng zhǔyào yíchǎn, guǎfu dìwèi dī. Dànshì xiànzài fǎlǜ guīdìng nánnǚ píngděng, bǎohù guǎfu quánlì. Zhōngguó gǔdài yě zhòng zhǎngzi, jiāpǔ jìlù xuémài. Xiàndài fǎlǜ guīdìng dì yī shùnxù jìchéngrén píngfēn yíchǎn, nánnǚ píngděng. Zǒngzhī, liǎngguó fǎlǜ dōu zǒuxiàng píngděng.}

« Le Cameroun et la Chine ont des similitudes et des différences dans les coutumes successorales. Autrefois au Cameroun, selon la coutume, le fils aîné héritait de l'héritage principal, le statut de la veuve était bas. Mais maintenant la loi stipule l'égalité hommes-femmes, protégeant les droits des veuves. La Chine ancienne valorisait aussi le fils aîné, la généalogie enregistrait la lignée. La loi moderne stipule que les héritiers du premier ordre partagent l'héritage à parts égales, hommes et femmes égaux. En résumé, les lois des deux pays tendent vers l'égalité. »

\textbf{Analyse didactique} :
\begin{itemize}
\item \textbf{Vocabulaire utilisé (11/12)} : \zh{继承, 习俗, 遗产, 长子, 寡妇, 法律, 规定, 男女平等, 家谱, 血脉, 第一顺序继承人}.
\item \textbf{Structures} : \zh{有异同} (avoir des différences/similitudes), \zh{过去...现在...} (autrefois... maintenant), \zh{法律规定 + proposition}, \zh{总之} (en conclusion).
\item \textbf{Connecteurs logiques} : \zh{但是} (adversatif), \zh{也} (aussi), \zh{总之} (conclusif).
\end{itemize}
\end{exoresolu}

\courssub{Expression orale : Mini-exposé et débat}

\begin{experience}[Mini-exposé comparatif : « Héritage et égalité »]
\textbf{Objectif} : Produire un exposé oral structuré (1-2 min) en chinois, comparer les deux systèmes et exprimer un avis personnel.

\textbf{Consignes} :
\begin{enumerate}
\item \textbf{Préparation (10 min)} : En binôme, préparez un plan en 3 parties : 1. Situation traditionnelle (Cameroun vs Chine) ; 2. Évolution légale moderne ; 3. Votre opinion : la loi suffit-elle pour changer les mentalités ?
\item \textbf{Vocabulaire imposé} : Utilisez au moins 6 mots du glossaire + \zh{虽然...但是...} / \zh{不仅...而且...}.
\item \textbf{Présentation} : Chaque binôme présente devant la classe. Les autres notent les arguments pour le débat.
\end{enumerate}

\textbf{Amorces utiles} :
\begin{itemize}
\item \zh{我认为法律很重要，但是习俗改变得慢。} (Wǒ rènwéi fǎlǜ hěn zhòngyào, dànshì xísú gǎibiàn de màn.)
\item \zh{在中国和喀麦隆，城市和农村的情况不一样。} (Zài Zhōngguó hé Kāmàilóng, chéngshì hé nóngcūn de qíngkuàng bù yīyàng.)
\item \zh{我们应该尊重传统，也要维护平等。} (Wǒmen yīnggāi zūnzhòng chuántǒng, yě yào wéihù píngděng.)
\end{itemize}
\end{experience}

\courssub{Exercices d'entraînement}

\begin{exercice}[Entraînement écrit — Vocabulaire, Caractères, Traduction]
\begin{enumerate}
\item \textbf{Écriture des caractères} : Écrivez 5 fois chacun en respectant l'ordre des traits : \zh{承}, \zh{法}, \zh{家}, \zh{继}, \zh{遗}. Cochez la case quand c'est fait. ☐ ☐ ☐ ☐ ☐
\item \textbf{Choix du classifieur / Bon mot} : Complétez avec le mot juste du glossaire.
\begin{enumerate}
\item 老王去世了，他的\_\_\_\_\_\_由三个儿女分掉了。 (héritage)
\item 我们要\_\_\_\_\_\_少数民族的传统\_\_\_\_\_\_。 (respecter / coutumes)
\item 她是独生女，是父亲唯一的合法\_\_\_\_\_\_。 (héritière)
\item 请律师帮忙\_\_\_\_\_\_一份\_\_\_\_\_\_。 (rédiger / testament)
\end{enumerate}
\item \textbf{Traduction vers le chinois (Thème)} :
\begin{enumerate}
\item La loi chinoise protège les droits des veuves.
\item Selon la coutume, le fils aîné gère les affaires de la famille.
\item Nous devons consulter la généalogie pour connaître nos ancêtres.
\item L'égalité hommes-femmes est un principe juridique fondamental.
\end{enumerate}
\item \textbf{Discrimination de caractères} : Entourez le caractère correct dans chaque phrase.
\begin{enumerate}
\item 这件衣服很贵 / 遗。
\item 我们要继 / 断 承传统文化。
\item 这个风俗 / 浴 很有意思。
\end{enumerate}
\end{enumerate}
\end{exercice}

\begin{corrige}[Exercice Entraînement écrit]
\begin{enumerate}
\item \textbf{Écriture} : Vérifiez l'ordre des traits ci-dessus (\zh{承} 10 traits, \zh{法} 8 traits, \zh{家} 10 traits, \zh{继} 10 traits, \zh{遗} 12 traits). Attention à la clé \zh{纟} pour \zh{继} (3 traits) et \zh{辶} pour \zh{遗} (3 traits en dernier).
\item \textbf{Complétion} :
\begin{enumerate}
\item \zh{遗产} (yíchǎn)
\item \zh{尊重} (zūnzhòng) ... \zh{习俗} (xísú)
\item \zh{继承人} (jìchéngrén)
\item \zh{立} (lì) ... \zh{遗嘱} (yízhǔ)
\end{enumerate}
\item \textbf{Traduction} :
\begin{enumerate}
\item \zh{中国法律保护寡妇的权利。} (Zhōngguó fǎlǜ bǎohù guǎfu de quánlì.)
\item \zh{按照习俗，长子管理家里的事务。} (Ànzhào xísú, zhǎngzi guǎnlǐ jiāli de shìwù.)
\item \zh{我们要查阅家谱，了解祖先。} (Wǒmen yào cháyuè jiāpǔ, liǎojiě zǔxiān.)
\item \zh{男女平等是基本法律原则。} (Nánnǚ píngděng shì jīběn fǎlǜ yuánzé.)
\end{enumerate}
\item \textbf{Discrimination} :
\begin{enumerate}
\item \zh{贵} (guì, cher/précieux — clé \zh{贝})
\item \zh{继} (jì, continuer — clé \zh{纟})
\item \zh{俗} (sú, coutume — clé \zh{亻})
\end{enumerate}
\end{enumerate}
\end{corrige}

\coursec{Réflexion et mini-exposé}

Après avoir étudié le vocabulaire et l'écriture des caractères du thème de l'hérédité (\zh{继承} \emph{jìchéng} « hériter », \zh{遗产} \emph{yíchǎn} « héritage », \zh{传统} \emph{chuántǒng} « tradition », \zh{法律} \emph{fǎlǜ} « loi »), il est temps de \cle{prendre la parole}. Dans cette dernière étape, tu vas apprendre à exprimer ton opinion en chinois, à la justifier avec des arguments ordonnés, puis à présenter devant la classe un mini-exposé de deux minutes : « L'héritage dans ma région ». C'est exactement le type de production orale attendu au baccalauréat.

\courssub{Observer : un élève donne son avis}

\begin{textebox}[Un élève de Terminale donne son avis]
我觉得女儿也应该继承遗产。\\
\emph{Wǒ juéde nǚér yě yīnggāi jìchéng yíchǎn.}\\
« Je pense que les filles aussi devraient hériter. »\\[4pt]
第一，女儿和儿子都是父母的孩子。\\
\emph{Dìyī, nǚér hé érzi dōu shì fùmǔ de háizi.}\\
« Premièrement, les filles et les fils sont tous les enfants de leurs parents. »\\[4pt]
第二，女儿也常常照顾老人。\\
\emph{Dìèr, nǚér yě chángcháng zhàogù lǎorén.}\\
« Deuxièmement, les filles aussi prennent souvent soin des personnes âgées. »\\[4pt]
所以，我认为我们应该尊重传统，也应该遵守法律。\\
\emph{Suǒyǐ, wǒ rènwéi wǒmen yīnggāi zūnzhòng chuántǒng, yě yīnggāi zūnshǒu fǎlǜ.}\\
« Donc, je considère que nous devons respecter les traditions et aussi respecter la loi. »
\end{textebox}

Observe la construction de cette prise de parole : une \cle{opinion} (\zh{我觉得……}), deux \cle{raisons} numérotées (\zh{第一……第二……}), puis une \cle{conclusion} (\zh{所以……}). C'est la trame universelle du mini-exposé.

\courssub{Exprimer son opinion : 觉得 et 认为}

\begin{propriete}[Les verbes d'opinion]
Pour donner un point de vue personnel, le chinois utilise deux verbes principaux, suivis d'une proposition complète :
\begin{itemize}
\item \zh{觉得} \emph{juéde} + phrase : « je pense que, je trouve que » (opinion personnelle, registre courant, oral) ;
\item \zh{认为} \emph{rènwéi} + phrase : « je considère que, j'estime que » (plus formel, adapté à un exposé ou à une dissertation).
\end{itemize}
Structure : \textbf{Sujet + 觉得 / 认为 + proposition}. Contrairement au français, on n'utilise jamais de subjonctif ni de « que » : la proposition suit directement le verbe.
\end{propriete}

\begin{exemplebox}[Opinions sur l'héritage]
\zh{我觉得儿子和女儿应该一样继承。} \emph{Wǒ juéde érzi hé nǚér yīnggāi yīyàng jìchéng.} « Je pense que les fils et les filles devraient hériter de la même façon. »\\
\zh{他认为传统很重要。} \emph{Tā rènwéi chuántǒng hěn zhòngyào.} « Il considère que la tradition est très importante. »\\
\zh{老师觉得我们应该听老人的意见。} \emph{Lǎoshī juéde wǒmen yīnggāi tīng lǎorén de yìjiàn.} « Le professeur pense que nous devrions écouter l'avis des anciens. »\\
\zh{我认为法律应该保护女儿。} \emph{Wǒ rènwéi fǎlǜ yīnggāi bǎohù nǚér.} « J'estime que la loi devrait protéger les filles. »
\end{exemplebox}

\begin{attention}[觉得 ou 认为 ?]
\zh{觉得} exprime une opinion personnelle, un ressenti (« je trouve que ») ; \zh{认为} est plus formel et plus affirmé (« je considère que »). Dans un exposé noté, préfère \zh{认为} pour la conclusion. Erreur fréquente des francophones : traduire « je pense que non » par \zh{我不觉得对} de façon maladroite — dis plutôt \zh{我觉得不对} \emph{Wǒ juéde bú duì} (« je pense que ce n'est pas correct »), la négation se place dans la proposition, pas devant \zh{觉得}.
\end{attention}

\courssub{Ordonner ses arguments}

\begin{methode}[Structurer un exposé en quatre étapes]
\begin{enumerate}
\item \textbf{Annoncer le sujet} : \zh{今天我说一说……} \emph{Jīntiān wǒ shuō yi shuō……} « Aujourd'hui, je vais parler de…… »
\item \textbf{D'abord} : \zh{首先……} \emph{shǒuxiān} — présente ton premier point (par exemple la coutume de ta région).
\item \textbf{Ensuite} : \zh{然后……} \emph{ránhòu} — ajoute le deuxième point (par exemple ce que dit la loi moderne).
\item \textbf{Enfin / conclusion} : \zh{最后……所以……} \emph{zuìhòu…… suǒyǐ……} — « enfin…… donc…… » : donne ton opinion avec \zh{我认为}.
\end{enumerate}
Pour énumérer les raisons à l'intérieur d'une étape, utilise \zh{第一……} \emph{dìyī}, \zh{第二……} \emph{dìèr}, \zh{第三……} \emph{dìsān}.
\end{methode}

\begin{center}
\begin{tabularx}{\linewidth}{|X|X|X|X|}
\hline
\rowcolor{popblueL} Caractères & Pinyin & Fonction & Traduction \\
\hline
我觉得…… & wǒ juéde & opinion (oral) & je pense que…… \\
\hline
我认为…… & wǒ rènwéi & opinion (formel) & je considère que…… \\
\hline
首先 & shǒuxiān & étape 1 & d'abord \\
\hline
然后 & ránhòu & étape 2 & ensuite \\
\hline
最后 & zuìhòu & étape 3 & enfin \\
\hline
第一 / 第二 & dìyī / dìèr & énumération & premièrement / deuxièmement \\
\hline
所以 & suǒyǐ & conclusion & donc, c'est pourquoi \\
\hline
和……一样 & hé…… yīyàng & comparaison & comme……, pareil que…… \\
\hline
但是……不一样 & dànshì…… bù yīyàng & contraste & mais…… ce n'est pas pareil \\
\hline
\end{tabularx}
\end{center}

\begin{popfigure}
\begin{tikzpicture}[
  node distance=1cm and 1.5cm,
  box/.style={draw=popblue, thick, rounded corners, fill=popblueL, align=center, minimum height=0.9cm, inner sep=4pt, text width=2.5cm},
  arr/.style={-{Stealth[length=3mm]}, thick, popdark}]
\node[box, align=center] (op){\zh{观点}\\ \emph{guāndiǎn}\\ opinion};
\node[box, below left=of op, align=center] (r1){\zh{理由一}\\ \emph{lǐyóu yī}\\ raison 1};
\node[box, below right=of op, align=center] (r2){\zh{理由二}\\ \emph{lǐyóu èr}\\ raison 2};
\node[box, below=2cm of op, fill=popgreenL, draw=popgreen, align=center] (cc){\zh{结论}\\ \emph{jiélùn}\\ conclusion};
\draw[arr] (op.south west) -- (r1.north);
\draw[arr] (op.south east) -- (r2.north);
\draw[arr] (r1.south) -- (cc.north west);
\draw[arr] (r2.south) -- (cc.north east);
\end{tikzpicture}
\legende{Organigramme de la prise de parole : une opinion, deux raisons, une conclusion.}
\end{popfigure}

\courssub{Questions de réflexion guidée}

Travaille d'abord par écrit, en 2 ou 3 phrases chinoises par question, avant le débat oral.

\begin{enumerate}
\item \zh{女儿应该不应该和儿子一样继承遗产？为什么？} \emph{Nǚér yīnggāi bu yīnggāi hé érzi yīyàng jìchéng yíchǎn? Wèishénme?} — Les filles doivent-elles hériter comme les fils ? Pourquoi ?
\item \zh{我们应该保留传统，还是只遵守现代法律？} \emph{Wǒmen yīnggāi bǎoliú chuántǒng, háishi zhǐ zūnshǒu xiàndài fǎlǜ?} — Faut-il conserver les coutumes de succession ou appliquer seulement la loi moderne ?
\item \zh{在一个家庭里，老人、叔叔、姑姑有什么作用？} \emph{Zài yí ge jiātíng lǐ, lǎorén, shūshu, gūgu yǒu shénme zuòyòng?} — Dans une succession, quel est le rôle de la famille élargie (anciens, oncles, tantes) ?
\end{enumerate}

\begin{exoresolu}[Répondre à une question de réflexion]
Énoncé : réponds en 3 phrases à la question \zh{女儿应该不应该继承遗产？} avec la trame 觉得 / 第一 / 第二 / 所以.
\tcblower
Solution détaillée. On applique l'organigramme : opinion → raison 1 → raison 2 → conclusion.\\
1) Opinion : \zh{我觉得女儿应该继承遗产。} \emph{Wǒ juéde nǚér yīnggāi jìchéng yíchǎn.} « Je pense que les filles devraient hériter. »\\
2) Raison 1 : \zh{第一，女儿也是父母的孩子。} \emph{Dìyī, nǚér yě shì fùmǔ de háizi.} « Premièrement, les filles aussi sont les enfants de leurs parents. »\\
3) Raison 2 : \zh{第二，法律说儿子和女儿一样。} \emph{Dìèr, fǎlǜ shuō érzi hé nǚér yīyàng.} « Deuxièmement, la loi dit que les fils et les filles sont égaux. »\\
4) Conclusion : \zh{所以，我认为传统和法律应该一起走。} \emph{Suǒyǐ, wǒ rènwéi chuántǒng hé fǎlǜ yīnggāi yìqǐ zǒu.} « Donc, je considère que la tradition et la loi doivent avancer ensemble. »
\end{exoresolu}

\courssub{Le mini-exposé : « L'héritage dans ma région »}

\begin{experience}[Mini-exposé de 2 minutes par binômes]
\begin{enumerate}
\item \textbf{Préparation (10 min)} : par binômes, rédigez 6 à 8 phrases sur le thème \zh{我们地区的继承} \emph{Wǒmen dìqū de jìchéng} « L'héritage dans notre région ». Utilisez obligatoirement : \zh{首先}, \zh{然后}, \zh{最后}, \zh{我觉得} ou \zh{我认为}, et au moins trois mots du vocabulaire de la leçon (\zh{遗产}, \zh{传统}, \zh{法律}, \zh{家庭}).
\item \textbf{Passation (2 min par binôme)} : chaque membre du binôme parle environ une minute.
\item \textbf{Écoute active} : pendant chaque exposé, chaque élève de la classe note sur sa feuille \cle{un point commun} (\zh{一样的地方}) et \cle{une différence} (\zh{不一样的地方}) relevés chez ses camarades, par exemple : \zh{我们一样：都说女儿可以继承。} / \zh{不一样：他们地区儿子先继承。}
\item \textbf{Bilan collectif} : le professeur relève au tableau les points communs et les différences des régions du Cameroun, puis compare avec la Chine.
\end{enumerate}
\end{experience}

\begin{exemplebox}[Début de mini-exposé modèle]
\zh{今天我说一说我们地区的继承。首先，在我们的传统里，儿子常常继承土地和房子。然后，女儿也得到一些东西，比如钱。最后，我觉得这个传统在变化。所以，我认为以后儿子和女儿会一样继承。}\\
\emph{Jīntiān wǒ shuō yi shuō wǒmen dìqū de jìchéng. Shǒuxiān, zài wǒmen de chuántǒng lǐ, érzi chángcháng jìchéng tǔdì hé fángzi. Ránhòu, nǚér yě dédào yìxiē dōngxi, bǐrú qián. Zuìhòu, wǒ juéde zhège chuántǒng zài biànhuà. Suǒyǐ, wǒ rènwéi yǐhòu érzi hé nǚér huì yīyàng jìchéng.}\\
« Aujourd'hui, je parle de l'héritage dans notre région. D'abord, dans notre tradition, les fils héritent souvent des terres et de la maison. Ensuite, les filles reçoivent aussi certaines choses, par exemple de l'argent. Enfin, je trouve que cette tradition est en train de changer. Donc, je considère qu'à l'avenir les fils et les filles hériteront de la même façon. »
\end{exemplebox}

\begin{attention}[Pièges à éviter à l'oral]
\begin{itemize}
\item Ne dis pas \zh{我是觉得} : \zh{觉得} est un verbe, pas besoin de \zh{是} devant. Dis \zh{我觉得}.
\item La structure \zh{因为……所以……} (« parce que…… donc…… ») est correcte et standard (niveau HSK 3). Évite simplement de dire \zh{因为……} sans conclusion ou \zh{……所以……} sans cause explicite dans une phrase complexe ; à l'oral débutant, une phrase simple avec \zh{所以} (« donc ») ou \zh{因为} (« parce que ») suffit souvent.
\item Dans \zh{第一}, prononce bien le quatrième ton de \zh{dì} puis le premier ton de \zh{yī} ; ne confonds pas avec \zh{第一} et un nombre isolé.
\end{itemize}
\end{attention}

\begin{savaistu}[La question socioculturelle au baccalauréat]
Au Bac, les questions socioculturelles demandent souvent de comparer le Cameroun et la Chine en 3 à 5 phrases. Entraîne-toi avec la trame de comparaison : \zh{和……一样} « comme…… » pour le point commun, \zh{但是……不一样} « mais…… ce n'est pas pareil » pour la différence. Exemple : \zh{中国和喀麦隆一样，家庭都很重要。但是继承的方法不一样。} \emph{Zhōngguó hé Kāmàilóng yīyàng, jiātíng dōu hěn zhòngyào. Dànshì jìchéng de fāngfǎ bù yīyàng.} « La Chine et le Cameroun sont pareils : la famille est très importante. Mais les méthodes de succession ne sont pas les mêmes. »
\end{savaistu}

\begin{exercice}[À toi de parler]
1. Complète avec 觉得 ou 认为 : \zh{我\_\_\_\_传统和法律都重要。} (registre formel, conclusion d'exposé).\\
2. Remets dans l'ordre pour faire une conclusion : 所以 / 女儿 / 应该 / 继承 / 我认为 / 遗产.\\
3. Rédige 4 phrases sur « L'héritage dans ma famille » avec 首先、然后、最后、我觉得.
\end{exercice}

\begin{corrige}[Exercice « À toi de parler »]
1. \zh{我认为传统和法律都重要。} — contexte formel (conclusion d'exposé), donc \zh{认为}.\\
2. \zh{所以，我认为女儿应该继承遗产。} \emph{Suǒyǐ, wǒ rènwéi nǚér yīnggāi jìchéng yíchǎn.} « Donc, je considère que les filles devraient hériter. » Ordre : connecteur de conclusion → opinion → sujet → modal → verbe → complément.\\
3. Modèle : \zh{首先，我的爷爷有很多土地。然后，他把土地给了我的爸爸和叔叔。最后，我觉得姑姑也应该得到土地。} \emph{Shǒuxiān, wǒ de yéye yǒu hěnduō tǔdì. Ránhòu, tā bǎ tǔdì gěile wǒ de bàba hé shūshu. Zuìhòu, wǒ juéde gūgu yě yīnggāi dédào tǔdì.} « D'abord, mon grand-père avait beaucoup de terres. Ensuite, il a donné les terres à mon père et à mes oncles. Enfin, je pense que mes tantes aussi devraient recevoir des terres. »
\end{corrige}

\newpage

\coursec{Activité d'intégration}

\courssub{Objectifs de l'activité}

À l'issue de cette activité d'intégration, l'élève sera capable de :
\begin{itemize}
\item mobiliser le \cle{vocabulaire} du thème de l'hérédité et de la succession (遗产, 继承, 传统, 平等…) dans un échange oral réel ;
\item exprimer un \cle{point de vue} simple et le justifier avec la structure 我觉得……因为…… ;
\item comparer des pratiques camerounaises et chinoises avec les structures de comparaison (跟……一样 / 不一样, 比) ;
\item débattre de façon \cle{respectueuse} en chinois simple, en écoutant et en réagissant aux arguments des autres ;
\item rédiger une courte \cle{synthèse écrite} (5 phrases) comparant la position défendue avec la réalité des deux pays.
\end{itemize}

\courssub{Situation}

Le club de chinois du lycée de Nkol-Eton (Yaoundé) organise sa grande journée portes ouvertes. Pour clôturer le module « Vie familiale et intégration sociale », les élèves de Terminale A ont choisi de présenter un \cle{débat parlementaire} en chinois devant les élèves de Première et quelques parents invités. Le thème retenu, proposé par la commission culturelle, est :

\begin{center}
\textbf{« Filles et fils, même héritage ? »} — 女儿和儿子应该有一样的遗产吗？
\end{center}

La classe est divisée en trois groupes : le camp « pour » l'égalité stricte dans la succession (赞成组), le camp « contre » ou favorable au respect des pratiques traditionnelles (反对组), et un groupe de « sages » (长老组), des notables qui écoutent, posent des questions et rendent un avis final. Pour préparer le débat, la commission a distribué à tous le document suivant :

\begin{textebox}[Document de préparation : deux points de vue]
\textbf{观点一：} 在喀麦隆的一些地方，传统上儿子继承土地和房子，女儿结婚以后去丈夫家。很多人觉得这是文化，应该尊重。\\
\emph{Zài Kāmàilóng de yìxiē dìfang, chuántǒng shang érzi jìchéng tǔdì hé fángzi, nǚ'ér jiéhūn yǐhòu qù zhàngfu jiā. Hěn duō rén juéde zhè shì wénhuà, yīnggāi zūnzhòng.}\\
« Dans certaines régions du Cameroun, traditionnellement, les fils héritent de la terre et de la maison ; après le mariage, la fille rejoint la famille de son mari. Beaucoup de gens pensent que c'est la culture et qu'il faut la respecter. »\\[4pt]
\textbf{观点二：} 现在时代变了。女儿也上学，也工作，也照顾父母。中国的法律说，儿子和女儿有一样的继承权。为什么喀麦隆不可以？\\
\emph{Xiànzài shídài biàn le. Nǚ'ér yě shàngxué, yě gōngzuò, yě zhàogù fùmǔ. Zhōngguó de fǎlǜ shuō, érzi hé nǚ'ér yǒu yíyàng de jìchéngquán. Wèishénme Kāmàilóng bù kěyǐ?}\\
« Aujourd'hui, les temps ont changé. Les filles aussi vont à l'école, travaillent et prennent soin de leurs parents. La loi chinoise dit que les fils et les filles ont les mêmes droits à l'héritage. Pourquoi le Cameroun ne le pourrait-il pas ? »
\end{textebox}

\courssub{Consignes}

\textbf{Étape 1 — Comprendre le document (10 min, en groupes).} Lis le document de préparation. Souligne dans chaque point de vue : (a) le fait décrit, (b) l'argument donné. Vérifie dans le tableau de vocabulaire les mots que tu ne connais pas. Chaque groupe pose une question de compréhension au professeur.

\textbf{Étape 2 — Préparer ses arguments (15 min).} Dans ton camp, rédige individuellement \cle{deux arguments} en chinois avec la trame obligatoire :

\begin{center}
我觉得……因为…… \quad \emph{Wǒ juéde... yīnwèi...} \quad « Je pense que... parce que... »
\end{center}

Exemples de départs : 我觉得女儿和儿子应该一样，因为…… / 我觉得传统很重要，因为……. Chaque argument doit contenir au moins \cle{un mot du lexique du thème} (遗产, 继承, 平等, 传统, 法律, 照顾…). Les « sages » préparent, eux, trois questions à poser aux deux camps (用 为什么 / 如果……怎么办？).

\textbf{Étape 3 — Débattre (20 min).} Le débat se déroule entièrement en chinois simple, sous la conduite du chef des sages :
\begin{enumerate}
\item Tour 1 : chaque camp présente ses deux arguments (1 minute par camp).
\item Tour 2 : les sages posent leurs questions ; chaque camp répond.
\item Tour 3 : réactions libres — pour réagir poliment, utilise 我同意…… (je suis d'accord) ou 对不起，我不同意，因为…… (désolé, je ne suis pas d'accord, parce que...).
\item Les sages rendent leur avis final en trois phrases : 我们听了两个组。我们觉得……。我们希望……。
\end{enumerate}

\textbf{Étape 4 — Rédiger la synthèse (15 min, par groupe).} Chaque groupe rédige une \cle{synthèse écrite de 5 phrases} en caractères : (1) rappel de la position du camp ; (2) premier argument ; (3) deuxième argument ; (4) comparaison avec la situation réelle au Cameroun ; (5) comparaison avec la situation en Chine (utilise au moins une fois 跟……一样 / 不一样 ou 比).

\textbf{Étape 5 — Affichage et relecture croisée (10 min).} Les synthèses sont affichées au mur de la classe. Chaque groupe lit celle d'un autre camp et note une phrase qu'il trouve particulièrement bien construite.

\courssub{Ressources à mobiliser}

\textbf{Lexique du thème :}

\begin{center}\begin{tabularx}{\linewidth}{|X|X|X|X|}\hline
\rowcolor{popblueL} Caractères & Pinyin & Nature & Traduction \\ \hline
遗产 & yíchǎn & n. & héritage, patrimoine \\ \hline
继承 & jìchéng & v. & hériter (de) \\ \hline
继承权 & jìchéngquán & n. & droit à l'héritage \\ \hline
传统 & chuántǒng & n. / adj. & tradition ; traditionnel \\ \hline
平等 & píngděng & adj. / n. & égal ; égalité \\ \hline
法律 & fǎlǜ & n. & loi \\ \hline
照顾 & zhàogù & v. & prendre soin de \\ \hline
尊重 & zūnzhòng & v. & respecter \\ \hline
文化 & wénhuà & n. & culture \\ \hline
时代 & shídài & n. & époque \\ \hline
\end{tabularx}\end{center}

\textbf{Structures d'opinion et de débat :}

\begin{center}\begin{tabularx}{\linewidth}{|X|X|X|}\hline
\rowcolor{popblueL} Structure & Exemple & Traduction \\ \hline
我觉得 + phrase & 我觉得女儿也应该继承。 & Je pense que les filles aussi devraient hériter. \\ \hline
……，因为…… & 应该平等，因为女儿也照顾父母。 & Il faut l'égalité, car les filles aussi prennent soin des parents. \\ \hline
我同意 / 我不同意 & 对不起，我不同意。 & Désolé, je ne suis pas d'accord. \\ \hline
\end{tabularx}\end{center}

\textbf{Structures de comparaison :}

\begin{center}\begin{tabularx}{\linewidth}{|X|X|X|}\hline
\rowcolor{popblueL} Structure & Exemple & Traduction \\ \hline
A 跟 B 一样 & 女儿跟儿子一样。 & Les filles sont comme les fils. \\ \hline
A 跟 B 不一样 & 传统跟法律不一样。 & La tradition n'est pas la même chose que la loi. \\ \hline
A 比 B + adj. & 现在比以前好。 & Aujourd'hui c'est mieux qu'avant. \\ \hline
\end{tabularx}\end{center}

\begin{attention}[Pièges à éviter pendant le débat]
\begin{itemize}
\item Ne dis pas \zh{我觉得是……因为} : 因为 introduit une proposition entière, pas seulement un nom. Dis 我觉得应该平等，\textbf{因为女儿也工作}。
\item 一样 s'emploie avec 跟 : \zh{女儿跟儿子一样}, et non \zh{女儿一样儿子}.
\item Attention au ton de 继承 (jìchéng, 4\textsuperscript{e} + 2\textsuperscript{e} ton) : ne le confonds pas avec 机场 (jīchǎng, aéroport) !
\item Dans 为什么喀麦隆不可以？, la négation porte sur 可以 : on ne dit pas \zh{不为什么}.
\end{itemize}
\end{attention}

\courssub{Proposition de production et corrigé commenté}

\begin{exoresolu}[Synthèse modèle du camp « pour » (赞成组)]
Rédige la synthèse écrite de 5 phrases du camp favorable à l'égalité stricte dans la succession, en respectant le plan : position, deux arguments, comparaison avec le Cameroun, comparaison avec la Chine.
\tcblower
\textbf{Production modèle :}

我们组觉得女儿和儿子应该有一样的继承权。\\
\emph{Wǒmen zǔ juéde nǚ'ér hé érzi yīnggāi yǒu yíyàng de jìchéngquán.}\\
« Notre groupe pense que les filles et les fils devraient avoir les mêmes droits à l'héritage. »

第一，女儿跟儿子一样上学、工作，也照顾父母。\\
\emph{Dì-yī, nǚ'ér gēn érzi yíyàng shàngxué, gōngzuò, yě zhàogù fùmǔ.}\\
« Premièrement, les filles, comme les fils, vont à l'école, travaillent et prennent aussi soin de leurs parents. »

第二，平等对家庭好：没有争吵，大家更团结。\\
\emph{Dì-èr, píngděng duì jiātíng hǎo : méiyǒu zhēngchǎo, dàjiā gèng tuánjié.}\\
« Deuxièmement, l'égalité est bonne pour la famille : pas de disputes, tout le monde est plus uni. »

在喀麦隆，很多地方的 traditions — 传统 — 跟这个想法不一样，儿子常常继承土地。\\
\emph{Zài Kāmàilóng, hěn duō dìfang de chuántǒng gēn zhè ge xiǎngfǎ bù yíyàng, érzi chángcháng jìchéng tǔdì.}\\
« Au Cameroun, dans beaucoup de régions, la tradition ne correspond pas à cette idée : ce sont souvent les fils qui héritent de la terre. »

中国的法律跟喀麦隆的传统不一样：它说儿子和女儿一样，我们觉得这样比以前好。\\
\emph{Zhōngguó de fǎlǜ gēn Kāmàilóng de chuántǒng bù yíyàng : tā shuō érzi hé nǚ'ér yíyàng, wǒmen juéde zhèyàng bǐ yǐqián hǎo.}\\
« La loi chinoise diffère de la tradition camerounaise : elle dit que fils et filles sont égaux, et nous pensons que c'est mieux qu'avant. »

\textbf{Commentaire des choix de langue :}
\begin{itemize}
\item La phrase 1 pose la position avec la trame imposée 觉得 + 应该 + 有, et emploie le mot-clé 继承权.
\item Les arguments sont introduits par 第一 / 第二, marqueurs simples et naturels pour structurer un propos oral ou écrit.
\item La phrase 2 applique correctement la structure A 跟 B 一样 + verbes partagés ; l'adverbe 也 (« aussi ») renforce l'argument.
\item La phrase 4 utilise la négation 跟……不一样 pour la comparaison avec la réalité camerounaise, avec 常常 pour nuancer (« souvent », pas « toujours ») — formulation factuelle et respectueuse.
\item La phrase 5 combine deux structures vues dans la leçon (跟……不一样 et 比) et se termine par un avis personnel, ce qui répond exactement à la compétence « exprimer un point de vue simple ».
\end{itemize}
\end{exoresolu}

\courssub{Bilan de l'activité}

Cette activité a permis de réinvestir l'ensemble des acquis de la leçon dans une tâche finale de communication authentique :
\begin{itemize}
\item \textbf{Compréhension de l'écrit} : lecture et exploitation d'un document présentant deux points de vue contrastés sur l'hérédité au Cameroun et en Chine ;
\item \textbf{Lexique} : emploi actif, à l'oral comme à l'écrit, des mots du thème (遗产, 继承, 继承权, 传统, 平等, 法律, 照顾, 尊重) ;
\item \textbf{Grammaire} : mobilisation des structures d'opinion (我觉得……因为……, 我同意 / 我不同意) et de comparaison (跟……一样 / 不一样, 比) ;
\item \textbf{Expression orale} : tenue d'un débat régi par des règles de politesse et de prise de parole, avec un arbitrage par les « sages », qui fait écho au rôle des notables dans les familles camerounaises ;
\item \textbf{Expression écrite} : production d'une synthèse structurée en 5 phrases, affichée et lue par les pairs ;
\item \textbf{Dimension socioculturelle} : comparaison factuelle et respectueuse des pratiques successorales camerounaises et chinoises, sans jugement de valeur sur les traditions.
\end{itemize}

\begin{aretenir}[Ce que je retiens]
Pour débattre en chinois, je sais : donner mon avis avec \zh{我觉得……因为……}, réagir poliment avec \zh{我同意} ou \zh{对不起，我不同意}, comparer avec \zh{跟……一样 / 不一样} et \zh{比}, et structurer un texte court avec \zh{第一、第二}. Je connais le vocabulaire essentiel pour parler d'héritage et de succession au Cameroun et en Chine.
\end{aretenir}

\newpage

\coursec{Méthodes et exercices résolus}

\coursec{Leçon 4 — Éléments socioculturels : processus d’hérédité au Cameroun et en Chine}

\begin{textebox}[Situation de communication : Discussion entre amis]
\zh{小林：阿明，你们喀麦隆的习俗里，女儿能继承房子吗？} \\
\zh{阿明：法律上当然能，法律保护男女平等。但是在很多传统家庭里，长子通常分得最多。} \\
\zh{小林：中国也一样。法律规定儿子女儿一样，但在农村老一辈人还是重男轻女。} \\
\zh{阿明：所以法律和传统之间，还要靠时间慢慢改变。}

Xiǎolín: Āmíng, nǐmen Kāmàilóng de xísú lǐ, nǚ'ér néng jìchéng fángzi ma? \\
Āmíng: Fǎlǜ shàng dāngrán néng, fǎlǜ bǎohù nánnǚ píngděng. Dànshì zài hěnduō chuántǒng jiātíng lǐ, zhǎngzǐ tōngcháng fēn dé zuì duō. \\
Xiǎolín: Zhōngguó yě yíyàng. Fǎlǜ guīdìng érzi nǚ'ér yíyàng, dàn zài nóngcūn lǎo yīdài rén hái shì zhòngnánqīngnǚ. \\
Āmíng: Suǒyǐ fǎlǜ hé chuántǒng zhī jiān, hái yào kào shíjiān mànmàn gǎibiàn.

« Xiao Lin : Amin, dans les coutumes camerounaises, les filles peuvent-elles hériter de la maison ? \\
Amin : Légalement oui, la loi protège l'égalité hommes-femmes. Mais dans beaucoup de familles traditionnelles, l'aîné reçoit généralement la plus grande part. \\
Xiao Lin : En Chine c'est pareil. La loi stipule l'égalité fils-filles, mais à la campagne les anciennes générations privilégient encore les garçons. \\
Amin : Donc entre la loi et la tradition, il faut encore du temps pour changer les choses. »
\end{textebox}

\begin{definition}[Qu'est-ce que l'hérédité ? / 什么是继承？]
L'\cle{hérédité} (\zh{继承} jìchéng) désigne la transmission du patrimoine (biens, terres, titres, nom) d'une génération à l'autre. Elle repose sur deux piliers qui coexistent et parfois s'affrontent :
\begin{itemize}
\item Le \cle{cadre légal} (\zh{法律} fǎlǜ) : lois écrites, code civil, constitution — même droits pour tous les enfants.
\item Le \cle{cadre coutumier} (\zh{传统习惯} chuántǒng xíguàn) : traditions orales, pratiques ancestrales, rôle de l'aîné, lignée patrilinéaire ou matrilinéaire.
\end{itemize}
Au Cameroun comme en Chine, la loi moderne proclame l'égalité, mais les traditions pèsent encore lourd dans le partage réel.
\end{definition}

\courssub{Savoir-faire 1 : Comparer des pratiques camerounaises et chinoises}

\begin{methode}[Comparer deux réalités socioculturelles en chinois]
Pour comparer l'hérédité au Cameroun et en Chine dans un texte ou à l'oral, suis ces étapes :
\begin{enumerate}
\item \cle{Identifier} le point de comparaison : qui hérite ? quoi hérite-t-on (biens, nom, terres, titres) ? selon quelles règles (loi, coutume, tradition) ?
\item \cle{Décrire} chaque réalité séparément avec des phrases simples : Sujet + Verbe + Complément. Exemple : \zh{在喀麦隆，儿子和女儿都可以继承父母的财产。}
\item \cle{Comparer} avec la structure de comparaison : \cle{A + 比 + B + adjectif}, ou la structure de similitude : \cle{A + 和/跟 + B + 一样/不一样 + adjectif}.
\item \cle{Nuancer} avec des connecteurs : \zh{但是} (mais), \zh{不过} (cependant), \zh{也} (aussi), \zh{都} (tous les deux).
\item \cle{Conclure} par une phrase de synthèse avec \zh{所以} (donc) ou \zh{总之} (en somme).
\end{enumerate}
\end{methode}

\begin{propriete}[Structures de comparaison et de similitude]
\begin{center}
\begin{tabularx}{\linewidth}{|X|X|X|}
\hline
\rowcolor{popblueL} Structure & Exemple (Caractères + Pinyin) & Traduction \\
\hline
A 比 B + adj. & 中国的传统习惯 \textbf{比} 喀麦隆的 \textbf{更严格}。 & La tradition chinoise est plus stricte que la camerounaise. \\
\hline
A 比 B + 更/还 + adj. & 在农村，长子 \textbf{比} 别的孩子 \textbf{还重要}。 & À la campagne, l'aîné est encore plus important que les autres. \\
\hline
A 和/跟 B 一样 + adj. & 中国 \textbf{和} 喀麦隆 \textbf{一样重视} 家庭。 & La Chine attache autant d'importance à la famille que le Cameroun. \\
\hline
A 和/跟 B 不一样 & 两个国家的继承习惯 \textbf{不一样}。 & Les coutumes d'héritage des deux pays sont différentes. \\
\hline
A 跟 B 有 \textbf{一样的} + Nom & 女儿 \textbf{跟} 儿子 \textbf{有一样的} 继承权。 & Les filles ont les mêmes droits d'héritage que les fils. \\
\hline
\end{tabularx}
\end{center}
\end{propriete}

\begin{attention}[Pièges de la comparaison]
\begin{itemize}
\item \textbf{Pas de \zh{很} après \zh{比}} : on écrit \zh{比较严格} ou \zh{更严格}, jamais \zh{比……很严格}。
\item \textbf{\zh{一样} + \zh{的} + Nom} : \zh{一样的权利} (les mêmes droits). Sans nom, \zh{一样} seul : \zh{他们一样} (ils sont pareils).
\item \textbf{\zh{跟} plus oral que \zh{和}} : \zh{跟} s'emploie davantage à l'oral, \zh{和} à l'écrit ; les deux sont corrects.
\end{itemize}
\end{attention}


\begin{exoresolu}[Comparaison guidée]
Énoncé : À l'aide des structures \zh{比} et \zh{和……一样}，rédige trois phrases comparant l'hérédité au Cameroun et en Chine à partir de ces faits : (1) en Chine, la loi donne les mêmes droits aux fils et aux filles ; (2) au Cameroun, dans certaines traditions, l'aîné de la famille joue un rôle important ; (3) dans les deux pays, la famille est très importante.
\tcblower
Solution détaillée :
\begin{enumerate}
\item Phrase avec \zh{比} (supériorité) : on choisit un adjectif pertinent (\zh{重要}, \zh{严格}, \zh{传统}). Structure : A + \zh{比} + B + \zh{更/还} + adjectif.
\zh{在喀麦隆的一些传统里，长子比别的孩子更重要。} \\
(Zài Kāmàilóng de yìxiē chuántǒng lǐ, zhǎngzǐ bǐ bié de háizi gèng zhòngyào.) \\
« Dans certaines traditions camerounaises, l'aîné est plus important que les autres enfants. »
\item Phrase avec \zh{一样} (égalité) : structure \zh{有 + 一样的 + Nom}.
\zh{中国的法律规定儿子和女儿有一样的继承权。} \\
(Zhōngguó de fǎlǜ guīdìng érzi hé nǚ'ér yǒu yíyàng de jìchéngquán.) \\
« La loi chinoise stipule que les fils et les filles ont les mêmes droits d'héritage. »
\item Phrase de synthèse (similitude + conclusion) :
\zh{中国和喀麦隆一样重视家庭，所以继承问题在两个国家都很重要。} \\
(Zhōngguó hé Kāmàilóng yíyàng zhòngshì jiātíng, suǒyǐ jìchéng wèntí zài liǎng ge guójiā dōu hěn zhòngyào.) \\
« La Chine et le Cameroun attachent la même importance à la famille, donc la question de l'héritage est importante dans les deux pays. »
\end{enumerate}
Conclusion : une bonne comparaison au Bac = une structure correcte (\zh{比} ou \zh{一样}), un vocabulaire précis du thème (\zh{继承}, \zh{传统}, \zh{法律}) et une conclusion logique avec \zh{所以}.
\end{exoresolu}

\courssub{Savoir-faire 2 : Employer le vocabulaire socioculturel du thème}

\begin{methode}[Mémoriser et utiliser le lexique de l'hérédité]
\begin{enumerate}
\item \cle{Regrouper} les mots par familles sémantiques (mindmap ci-dessous).
\item \cle{Apprendre} chaque mot avec son caractère, son pinyin (tons !) et une phrase exemple, jamais isolément.
\item \cle{Vérifier} les collocations (associations naturelles) : \zh{继承财产} (hériter des biens), \zh{尊重传统} (respecter la tradition), \zh{分给子女} (répartir entre les enfants), \zh{保护权利} (protéger les droits).
\item \cle{Réutiliser} les mots dans tes propres phrases : un mot n'est acquis que s'il est employé en contexte.
\end{enumerate}
\end{methode}

\begin{popfigure}
\begin{tikzpicture}[mindmap, grow cyclic, every node/.style=concept, concept color=popblue,
    level 1/.append style={level distance=4.5cm, sibling angle=90},
    level 2/.append style={level distance=3cm, sibling angle=45}]
\node[concept color=popdark] {Hérédité / 继承}
    child[concept color=popgreen] {node {Famille / 家庭}
        child {node {父母 fùmǔ}}
        child {node {儿子 érzi}}
        child {node {女儿 nǚ'ér}}
        child {node {长子 zhǎngzǐ}}
    }
    child[concept color=poporange] {node {Biens / 财产}
        child {node {房子 fángzi}}
        child {node {土地 tǔdì}}
        child {node {财产 cáichǎn}}
        child {node {一部分 yí bùfen}}
    }
    child[concept color=poppurple] {node {Règles / 规则}
        child {node {法律 fǎlǜ}}
        child {node {传统 chuántǒng}}
        child {node {习惯 xíguàn}}
        child {node {权利 quánlì}}
    }
    child[concept color=poppink] {node {Actions / 动作}
        child {node {继承 jìchéng}}
        child {node {分给 fēngěi}}
        child {node {尊重 zūnzhòng}}
        child {node {保护 bǎohù}}
    };
\end{tikzpicture}
\legende{Familles lexicales du thème « Héritage »}
\end{popfigure}

\begin{center}
\begin{tabularx}{\linewidth}{|X|X|X|X|}
\hline
\rowcolor{popblueL} Caractères & Pinyin & Nature & Traduction \\
\hline
继承 & jìchéng & verbe & hériter \\
\hline
财产 & cáichǎn & nom & biens, patrimoine \\
\hline
传统 & chuántǒng & nom & tradition \\
\hline
法律 & fǎlǜ & nom & loi \\
\hline
习惯 & xíguàn & nom & coutume, habitude \\
\hline
长子 & zhǎngzǐ & nom & fils aîné \\
\hline
权利 & quánlì & nom & droit \\
\hline
公平 & gōngpíng & adjectif & juste, équitable \\
\hline
平等 & píngděng & adjectif/nom & égalité \\
\hline
重男轻女 & zhòng nán qīng nǚ & locution & privilégier les garçons aux filles \\
\hline
代代相传 & dàidài xiāngchuán & verbe & se transmettre de génération en génération \\
\hline
去世 & qùshì & verbe & décéder (registre soutenu) \\
\hline
负责 & fùzé & verbe & être chargé de, responsable de \\
\hline
规定 & guīdìng & verbe/nom & stipuler, règlement \\
\hline
照顾 & zhàogù & verbe & prendre soin de, s'occuper de \\
\hline
\end{tabularx}
\end{center}

\begin{aretenir}[Écriture des caractères clés : ordre des traits et radicaux]
\begin{itemize}
\item \zh{继} (jì) : 10 traits. Radicale \zh{纟} (soie). Ordre : haut → bas (纟), puis \zh{米} (riz), puis \zh{攵} (main tenant un bâton). Sens : continuer, succéder.
\item \zh{承} (chéng) : 8 traits. Radicale \zh{手} (main, forme \zh{扌} à gauche). Ordre : \zh{禾} (céréale) puis \zh{扌}. Sens : recevoir, supporter.
\item \zh{财} (cái) : 7 traits. Radicale \zh{贝} (coquille/valeur). Ordre : \zh{贝} puis \zh{才}. Sens : richesse, argent.
\item \zh{传} (chuán) : 6 traits. Radicale \zh{亻} (homme). Ordre : \zh{亻} puis \zh{专}. Sens : transmettre.
\item \zh{长} (zhǎng/zhǎng/zhǎng) : 4 traits. Radicale \zh{长} (long/chef). Attention : \zh{zhǎng} (aîné, chef) vs \zh{cháng} (long). Trait horizontal long en bas.
\item \zh{律} (lǜ) : 9 traits. Radicale \zh{彳} (pas). Ordre : \zh{彳} puis \zh{聿} (pinceau). Sens : loi, règle.
\end{itemize}
\end{aretenir}

\begin{exoresolu}[Vocabulaire en contexte]
Énoncé : Complète les phrases avec le mot qui convient : \zh{继承 / 财产 / 传统 / 法律 / 公平}.
(1) 中国的 \_\_\_\_ 保护女儿的继承权。
(2) 父母去世以后，孩子们可以 \_\_\_\_ 他们的房子。
(3) 在喀麦隆，很多 \_\_\_\_ 是代代相传的。
(4) 把 \_\_\_\_ 分给所有的孩子，这样比较 \_\_\_\_。
\tcblower
Solution détaillée :
\begin{enumerate}
\item (1) \zh{法律} (fǎlǜ) : le sujet de « protéger les droits » est la loi. Phrase : \zh{中国的法律保护女儿的继承权。}
\item (2) \zh{继承} (jìchéng) : après \zh{可以} (pouvoir) il faut un verbe. \zh{父母去世以后，孩子们可以继承他们的房子。} Justification : \zh{去世} = décéder (soutenu) ; \zh{以后} marque l'antériorité.
\item (3) \zh{传统} (chuántǒng) : \zh{代代相传} (transmis de génération en génération) qualifie les traditions. \zh{在喀麦隆，很多传统是代代相传的。}
\item (4) \zh{财产} puis \zh{公平} : on répartit des biens (\zh{财产}, nom) ; l'adjectif \zh{公平} qualifie la répartition. \zh{把财产分给所有的孩子，这样比较公平。} Justification : construction \zh{把} + objet + verbe ; \zh{比较} + adjectif = « assez / plutôt ».
\end{enumerate}
\end{exoresolu}

\courssub{Savoir-faire 3 : Comprendre un texte socioculturel}

\begin{methode}[Lire et comprendre un texte sur l'hérédité]
\begin{enumerate}
\item \cle{Lire} une première fois sans dictionnaire pour saisir le thème général (quel pays ? quel problème ?).
\item \cle{Repérer} les mots-clés du thème : \zh{继承, 财产, 传统, 法律, 家庭, 长子, 公平}.
\item \cle{Analyser} la structure de chaque phrase : trouve le Sujet, le Verbe, les Compléments ; repère les marqueurs \zh{在……里} (dans), \zh{根据} (selon), \zh{但是/不过} (mais).
\item \cle{Répondre} aux questions en recopiant la structure de la question : question en \zh{吗} → réponse oui/non + phrase ; question en \zh{什么/谁/为什么} → reprendre l'élément demandé.
\item \cle{Reformuler} en français pour vérifier la compréhension globale.
\end{enumerate}
\end{methode}

\begin{textebox}[Texte support : Héritage, loi et tradition]
\zh{在中国，法律规定儿子和女儿有一样的继承权。不过，在一些农村，老传统还很重要：房子和土地常常给儿子。在喀麦隆，情况也差不多：法律承认所有孩子的权利，但是很多家庭按照传统习惯分财产，长子常常得到更多的东西。}

Zhōngguó, fǎlǜ guīdìng érzi hé nǚ'ér yǒu yíyàng de jìchéngquán. Bùguò, zài yìxiē nóngcūn, lǎo chuántǒng hái hěn zhòngyào: fángzi hé tǔdì chángcháng gěi érzi. Zài Kāmàilóng, qíngkuàng yě chàbuduō: fǎlǜ chéngrèn suǒyǒu háizi de quánlì, dànshì hěnduō jiātíng ànzhào chuántǒng xíguàn fēn cáichǎn, zhǎngzǐ chángcháng dédào gèng duō de dōngxi.
\end{textebox}

\begin{exoresolu}[Compréhension de texte]
Énoncé : Lis le texte ci-dessus puis réponds aux questions en chinois.
Questions : (1) 中国的法律说什么？ (2) 在中国的农村，谁常常得到房子和土地？ (3) 喀麦隆的情况和中国的差不多吗？
\tcblower
Solution détaillée :
\begin{enumerate}
\item (1) Question sur la loi chinoise ; réponse dans la 1\textsuperscript{re} phrase. \zh{中国的法律规定儿子和女儿有一样的继承权。} « La loi chinoise stipule que fils et filles ont les mêmes droits d'héritage. »
\item (2) Question en \zh{谁} (qui) : la réponse est un nom de personne. \zh{儿子常常得到房子和土地。} ou \zh{房子和土地常常给儿子。} « Ce sont souvent les fils qui reçoivent la maison et les terres. »
\item (3) Question en \zh{吗} : répondre oui + justification. \zh{差不多。在喀麦隆，法律也承认所有孩子的权利，但是很多家庭按照传统分财产。} « À peu près. Au Cameroun, la loi reconnaît aussi les droits de tous les enfants, mais beaucoup de familles répartissent les biens selon la tradition. »
\end{enumerate}
Conclusion : au Bac, recopie la structure de la question dans ta réponse et cite les éléments du texte : c'est la méthode la plus sûre.
\end{exoresolu}

\courssub{Savoir-faire 4 : Exprimer un point de vue simple (Expression écrite/orale)}

\begin{methode}[Donner son opinion sur l'hérédité]
\begin{enumerate}
\item \cle{Annoncer} ton opinion avec \zh{我觉得} / \zh{我认为} + phrase : \zh{我觉得……} (je pense que…).
\item \cle{Justifier} avec \zh{因为……所以……} : la cause d'abord, la conséquence ensuite (les deux particules se gardent en chinois).
\item \cle{Illustrer} avec un exemple : \zh{比如说} (par exemple), \zh{在我的国家} (dans mon pays).
\item \cle{Conclure} avec \zh{所以} ou \zh{总之} + reformulation courte.
\item \cle{Garder} des phrases courtes : une idée par phrase, ordre Sujet + Verbe + Complément.
\end{enumerate}
\end{methode}

\begin{experience}[Activité orale : Débat « Loi vs Tradition »]
\textbf{Consigne :} En binôme. L'un défend la position « La loi doit primer sur la tradition » (\zh{法律比传统更重要}), l'autre « La tradition doit être respectée » (\zh{传统也要尊重}). Utilisez : \zh{我觉得……因为……比如说……总之……}. Durée : 3 min par binôme. L'enseignant note l'usage de \zh{比/一样}, \zh{因为/所以}, \zh{应该}, vocabulaire du thème.
\end{experience}

\begin{exoresolu}[Mini-exposé écrit]
Énoncé : Rédige un court texte (4 à 6 phrases) en chinois : « Penses-tu que les filles doivent avoir les mêmes droits d'héritage que les fils ? Pourquoi ? »
\tcblower
Solution détaillée — modèle de réponse :

\zh{我觉得女儿应该有和儿子一样的继承权。因为女儿也是父母的孩子，她们也照顾父母，所以分财产的时候应该公平。比如说，在我的国家，很多女儿帮助家庭很多年。如果她们什么也不能继承，这不公平。总之，法律和传统都应该尊重女儿的权利。}

(Wǒ juéde nǚ'ér yīnggāi yǒu hé érzi yíyàng de jìchéngquán. Yīnwèi nǚ'ér yě shì fùmǔ de háizi, tāmen yě zhàogù fùmǔ, suǒyǐ fēn cáichǎn de shíhou yīnggāi gōngpíng. Bǐrú shuō, zài wǒ de guójiā, hěn duō nǚ'ér bāngzhù jiātíng hěn duō nián. Rúguǒ tāmen shénme yě bù néng jìchéng, zhè bù gōngpíng. Zǒngzhī, fǎlǜ hé chuántǒng dōu yīnggāi zūnzhòng nǚ'ér de quánlì.)

« Je pense que les filles devraient avoir les mêmes droits d'héritage que les fils. Parce que les filles sont aussi les enfants de leurs parents et qu'elles prennent aussi soin d'eux, la répartition des biens devrait être équitable. Par exemple, dans mon pays, beaucoup de filles aident leur famille pendant de nombreuses années. Si elles ne peuvent rien hériter, ce n'est pas juste. En somme, la loi et la tradition devraient respecter les droits des filles. »

Analyse grammaticale : \zh{应该} + verbe = devoir ; \zh{因为……所以……} relie cause et conséquence (on garde les deux en chinois) ; \zh{的时候} marque le moment ; \zh{如果……} (si…) introduit l'hypothèse ; \zh{什么也不能} = « ne rien pouvoir » (\zh{也} dans la négation totale) ; \zh{总之} conclut.
\end{exoresolu}

\courssub{Savoir-faire 5 : Traduire un texte socioculturel (Version : Chinois → Français)}

\begin{methode}[Traduire du chinois vers le français]
\begin{enumerate}
\item \cle{Découper} la phrase chinoise en groupes de sens : Sujet / Verbe / Compléments.
\item \cle{Traduire} les compléments de lieu et de temps en premier en français : \zh{在喀麦隆的农村} → « dans les campagnes camerounaises ».
\item \cle{Respecter} l'ordre français (SVO) : le chinois place les compléments (lieu, temps, manière) AVANT le verbe, le français souvent APRÈS.
\item \cle{Choisir} le bon registre : \zh{去世} = « décéder » (formel) ; \zh{按照} = « conformément à / selon » ; \zh{负责} = « être chargé de ».
\item \cle{Relire} en vérifiant que chaque mot chinois a un équivalent et que la phrase française est naturelle.
\end{enumerate}
\end{methode}

\begin{exoresolu}[Version]
Énoncé : Traduis en français : \zh{在很多喀麦隆家庭里，父母去世以后，长子负责分财产，他也得到最多的一部分。}
\tcblower
Solution détaillée :
\begin{enumerate}
\item Découpage : \zh{在很多喀麦隆家庭里} (CL) / \zh{父母去世以后} (CT) / \zh{长子} (Sujet) / \zh{负责分财产} (V + Objet) / \zh{他也得到最多的一部分} (2\textsuperscript{e} proposition).
\item Traduction des groupes : « dans beaucoup de familles camerounaises » ; « après le décès des parents » ; « le fils aîné » ; « est chargé de répartir les biens » ; « il reçoit aussi la plus grande part ».
\item Assemblage : « Dans beaucoup de familles camerounaises, après le décès des parents, le fils aîné est chargé de répartir les biens et il reçoit aussi la plus grande part. »
\item Points de vigilance : \zh{负责} + V = « être chargé de » ; \zh{最多的} = superlatif avec \zh{最} ; \zh{一部分} = « une part » (classificateur \zh{部分}). Ne pas calquer \zh{也} en début de phrase française : le placer correctement (« il reçoit \textbf{aussi} »).
\end{enumerate}
\end{exoresolu}

\begin{exoresolu}[Thème : Français → Chinois]
Énoncé : Traduis en chinois : « Au Cameroun, la loi protège les droits des filles, mais la coutume donne souvent plus aux fils aînés. »
\tcblower
Solution :
\zh{在喀麦隆，法律保护女儿的权利，但是习惯常常给长子更多。} \\
(Zài Kāmàilóng, fǎlǜ bǎohù nǚ'ér de quánlì, dànshì xíguàn chángcháng gěi zhǎngzǐ gèng duō.)
Justification : \zh{在} + lieu en tête ; \zh{但是} pour l'opposition ; \zh{习惯} (coutume) sujet du 2\textsuperscript{e} verbe ; \zh{给} + destinataire + quantité.
\end{exoresolu}

\courssub{Savoir-faire 6 : Traduire un texte socioculturel (Thème : Français → Chinois)}

\begin{methode}[Traduire du français vers le chinois]
\begin{enumerate}
\item \cle{Analyser} la phrase française : identifier le Sujet, le Verbe, les Compléments (lieu, temps, objet).
\item \cle{Réordonner} selon la syntaxe chinoise : \textbf{Sujet + (Lieu/Temps) + Verbe + Objet}. Le Lieu/Temps monte AVANT le verbe.
\item \cle{Choisir} le vocabulaire précis du thème : \zh{继承} (hériter), \zh{财产} (biens), \zh{长子} (aîné), \zh{公平} (équitable).
\item \cle{Vérifier} les mots-outils : \zh{的} (nominalisation), \zh{了} (aspect accompli si action passée), \zh{把} (objet antéposé), classificateurs (\zh{个, 部分}).
\item \cle{Contrôler} les tons et l'écriture des caractères.
\end{enumerate}
\end{methode}

\begin{attention}[Les 5 erreurs qui coûtent des points au Bac]
\begin{enumerate}
\item \cle{Tons faux ou absents} : confondre \zh{继承} (jìchéng, hériter) avec des homophones ; écrire le pinyin sans tons est sanctionné. Vérifie chaque syllabe : \zh{长子} (zhǎngzǐ, aîné) avec 3\textsuperscript{e} ton sur \zh{长}, et non \zh{cháng} (long).
\item \cle{Caractères confondus} : \zh{传统} (tradition, radical \zh{亻}) ≠ \zh{专} (spécial) ; \zh{财产} (biens, clé \zh{贝} valeur) ≠ \zh{材料} (matériau, clé \zh{木} bois). Un trait oublié change le sens.
\item \cle{Ordre des mots calqué du français} : Complément de lieu/temps AVANT le verbe : \zh{在喀麦隆，长子分财产。} Jamais « \zh{长子分财产在喀麦隆} ». Similitude : \zh{跟父母一样} (pareil aux parents), ordre inverse du français.
\item \cle{Mauvaise structure de comparaison} : Après \zh{比}, pas de \zh{很} : \zh{传统比法律更严格} ou adj. seul. Pour l'égalité, n'oublie pas \zh{一样} : \zh{女儿和儿子一样有权利。}
\item \cle{Classificateurs oubliés ou faux} : \zh{一个家庭}, \zh{两个孩子} (jamais \zh{二个}), \zh{一部分}. \zh{两} devant classif., \zh{二} pour compter/nombres.
\end{enumerate}
\end{attention}

\courssub{Entraînement et Évaluation}

\begin{exercice}[Exercices de révision — Leçon 4]
\begin{enumerate}
\item \textbf{Vocabulaire :} Associe chaque définition au mot chinois (caractères + pinyin).
    \begin{enumerate}
        \item Transmettre de génération en génération.
        \item Biens, patrimoine.
        \item Fils aîné.
        \item Juste, équitable.
        \item Décéder (formel).
    \end{enumerate}
\item \textbf{Grammaire :} Transforme les phrases en utilisant la structure demandée.
    \begin{enumerate}
        \item (Comparaison \zh{比}) La tradition camerounaise est stricte. La tradition chinoise est plus stricte. → \zh{中国的传统……}
        \item (Similitude \zh{一样}) Les filles ont des droits. Les fils ont des droits. → \zh{女儿跟儿子……}
        \item (Structure \zh{把}) Le père répartit les biens aux enfants. → \zh{父亲把……}
    \end{enumerate}
\item \textbf{Compréhension :} Relis le texte de la \texttt{textebox} « Héritage, loi et tradition ». Réponds en chinois :
    \begin{enumerate}
        \item 根据课文，中国农村的老传统怎么分房子和土地？
        \item 喀麦隆的家庭怎么分财产？
    \end{enumerate}
\item \textbf{Expression écrite :} Écris 5 phrases en chinois : « Dans ton pays, comment sont partagés les biens après le décès des parents ? Est-ce juste ? » Utilise \zh{因为……所以……}, \zh{应该}, \zh{公平}.
\item \textbf{Traduction (Version) :} Traduis en français : \zh{根据中国的法律，男女平等，女儿有权利继承父母的房子和土地。}
\item \textbf{Traduction (Thème) :} Traduis en chinois : « Bien que la loi soit équitable, beaucoup de familles suivent encore la tradition. » (Indice : \zh{虽然……但是……})
\end{enumerate}
\end{exercice}

\begin{corrige}[Exercices de révision — Leçon 4]
\begin{enumerate}
\item \textbf{Vocabulaire :}
    \begin{enumerate}
        \item \zh{代代相传} (dàidài xiāngchuán)
        \item \zh{财产} (cáichǎn)
        \item \zh{长子} (zhǎngzǐ)
        \item \zh{公平} (gōngpíng)
        \item \zh{去世} (qùshì)
    \end{enumerate}
\item \textbf{Grammaire :}
    \begin{enumerate}
        \item \zh{中国的传统比喀麦隆的更严格。} (ou \zh{比喀麦隆的严格})
        \item \zh{女儿跟儿子有一样的权利。} / \zh{女儿和儿子一样有权利。}
        \item \zh{父亲把财产分给孩子们。}
    \end{enumerate}
\item \textbf{Compréhension :}
    \begin{enumerate}
        \item \zh{房子和土地常常给儿子。}
        \item \zh{很多家庭按照传统习惯分财产，长子常常得到更多的东西。}
    \end{enumerate}
\item \textbf{Expression écrite (modèle) :}
\zh{在我国家，父母去世以后，财产通常分给所有孩子。但是有些地方，长子得到更多。我觉得这样不公平，因为女儿也照顾父母。所以法律应该保护女儿的权利，传统也要改变。}
\item \textbf{Version :} « Selon la loi chinoise, hommes et femmes sont égaux, les filles ont le droit d'hériter de la maison et des terres de leurs parents. »
\item \textbf{Thème :}
\zh{虽然法律很公平，但是很多家庭还是按照传统办事。} (ou \zh{还是遵守传统。})
\end{enumerate}
\end{corrige}

\newpage

\coursec{Exercices}

\courssub{Je vérifie mes connaissances}

\begin{exercice}[QCM : comprendre les mots-clés]
Pour chaque question, entoure la bonne réponse.
\begin{enumerate}
\item \zh{遗产} (yíchǎn) signifie :
\begin{enumerate}
\item la maison \item l'héritage \item le mariage
\end{enumerate}
\item \zh{继承} (jìchéng) signifie :
\begin{enumerate}
\item hériter \item partager \item vendre
\end{enumerate}
\item \zh{女儿} (nǚ'ér) signifie :
\begin{enumerate}
\item la mère \item la fille \item la sœur
\end{enumerate}
\item \zh{法律} (fǎlǜ) signifie :
\begin{enumerate}
\item la tradition \item la famille \item la loi
\end{enumerate}
\end{enumerate}
\end{exercice}

\begin{exercice}[Vrai ou faux ? Justifie ta réponse]
Réponds par vrai (对 duì) ou faux (不对 bú duì), puis justifie en une phrase en français.
\begin{enumerate}
\item En Chine, la loi de succession de 1985 donne des droits égaux aux fils et aux filles.
\item Au Cameroun, toutes les coutumes de succession sont exactement les mêmes dans toutes les régions.
\item \zh{长子} (zhǎngzǐ) désigne le fils aîné.
\item En Chine comme au Cameroun, la tradition joue encore un rôle dans le partage de l'héritage.
\end{enumerate}
\end{exercice}

\begin{exercice}[Vocabulaire : complète le tableau]
Complète le tableau suivant avec la nature du mot (nom, verbe, adjectif) et la traduction française.
\begin{center}
\begin{tabularx}{\linewidth}{|X|X|X|X|}
\hline
\rowcolor{popblueL} Caractères & Pinyin & Nature & Traduction \\ \hline
\zh{传统} & chuántǒng & \ldots & \ldots \\ \hline
\zh{平等} & píngděng & \ldots & \ldots \\ \hline
\zh{分} & fēn & \ldots & \ldots \\ \hline
\zh{习俗} & xísú & \ldots & \ldots \\ \hline
\zh{权利} & quánlì & \ldots & \ldots \\ \hline
\zh{家庭} & jiātíng & \ldots & \ldots \\ \hline
\end{tabularx}
\end{center}
\end{exercice}

\begin{exercice}[Associe les caractères au pinyin]
Relie chaque mot en caractères à son pinyin.
\begin{center}
\begin{tabularx}{\linewidth}{|X|X|}
\hline
\rowcolor{popblueL} Caractères & Pinyin \\ \hline
1. \zh{继承人} & a. fǎlǜ \\ \hline
2. \zh{遗嘱} & b. hòudài \\ \hline
3. \zh{法律} & c. jìchéngrén \\ \hline
4. \zh{后代} & d. yízhǔ \\ \hline
5. \zh{酋长} & e. qiúzhǎng \\ \hline
\end{tabularx}
\end{center}
\end{exercice}

\courssub{Je m'entraîne}

\begin{exercice}[Traduis en français]
Traduis les phrases suivantes en français.
\begin{enumerate}
\item \zh{在中国，儿子和女儿有平等的继承权。}\\ \emph{Zài Zhōngguó, érzi hé nǚ'ér yǒu píngděng de jìchéngquán.}
\item \zh{在喀麦隆，很多家庭一起讨论怎么分遗产。}\\ \emph{Zài Kāmàilóng, hěn duō jiātíng yìqǐ tǎolun zěnme fēn yíchǎn.}
\item \zh{按照传统，长子常常继承父亲的土地。}\\ \emph{Ànzhào chuántǒng, zhǎngzǐ chángcháng jìchéng fùqin de tǔdì.}
\end{enumerate}
\end{exercice}

\begin{exercice}[Texte à trous : les mots-outils]
Complète le texte suivant avec les mots-outils \zh{的}, \zh{和}, \zh{比}, \zh{在}, \zh{了}.
\begin{textebox}[Texte à trous]
以前，\zh{在很多地方，女儿不能继承父亲\_\_\_\_土地。现在，中国\_\_\_\_喀麦隆的法律都说：儿子\_\_\_\_女儿是平等\_\_\_\_。我觉得今天\_\_\_\_情况\_\_\_\_以前好多了。}\\ \emph{Yǐqián, zài hěn duō dìfang, nǚ'ér bù néng jìchéng fùqin \_\_\_\_ tǔdì. Xiànzài, Zhōngguó \_\_\_\_ Kāmàilóng de fǎlǜ dōu shuō: érzi \_\_\_\_ nǚ'ér shì píngděng \_\_\_\_. Wǒ juéde jīntiān \_\_\_\_ qíngkuàng \_\_\_\_ yǐqián hǎo duō le.}\\ « Autrefois, dans beaucoup d'endroits, les filles ne pouvaient pas hériter de la terre de leur père. Aujourd'hui, les lois de la Chine et du Cameroun disent toutes les deux : fils et filles sont égaux. Je trouve que la situation d'aujourd'hui est bien meilleure qu'avant. »
\end{textebox}
\end{exercice}

\begin{exercice}[Remets les mots dans l'ordre]
Remets les mots dans le bon ordre pour former des phrases correctes, puis traduis-les en français.
\begin{enumerate}
\item \zh{继承权 / 女儿 / 有 / 也 / 儿子 / 和 / 一样 / 的}
\item \zh{家庭 / 遗产 / 在喀麦隆 / 一起 / 讨论 / 常常 / 分 / 怎么}
\item \zh{比 / 法律 / 我 / 觉得 / 传统 / 重要 / 更}
\end{enumerate}
\end{exercice}

\begin{exercice}[Transforme les phrases]
Transforme chaque phrase selon la consigne entre parenthèses.
\begin{enumerate}
\item \zh{儿子继承土地。} (mets la phrase au négatif avec \zh{不})
\item \zh{他继承了父亲的房子。} (mets la phrase au négatif du passé avec \zh{没有})
\item \zh{女儿和儿子平等。} (transforme avec la structure \zh{和……一样})
\item \zh{传统很重要。法律很重要。} (fais une seule phrase avec \zh{和……一样})
\end{enumerate}
\end{exercice}

\courssub{Je me prépare au Bac}

\begin{exercice}[Compréhension de texte — type Bac]
Lis le texte suivant, puis réponds aux questions en chinois, par des phrases complètes.
\begin{textebox}[La loi chinoise sur l'héritage]
\zh{一九八五年，中国有了继承法。这个法律说：儿子和女儿有一样的继承权。妻子也可以继承丈夫的遗产。以前，在很多地方，只有儿子能继承父母的财产。现在，法律面前人人平等。但是，在一些农村，传统还是很重要，有些家庭还是按照老习惯分遗产。}\\ \emph{Yī jiǔ bā wǔ nián, Zhōngguó yǒu le jìchéngfǎ. Zhège fǎlǜ shuō: érzi hé nǚ'ér yǒu yíyàng de jìchéngquán. Qīzi yě kěyǐ jìchéng zhàngfu de yíchǎn. Yǐqián, zài hěn duō dìfang, zhǐyǒu érzi néng jìchéng fùmǔ de cáichǎn. Xiànzài, fǎlǜ miànqián rénrén píngděng. Dànshì, zài yìxiē nóngcūn, chuántǒng háishi hěn zhòngyào, yǒuxiē jiātíng háishi ànzhào lǎo xíguàn fēn yíchǎn.}\\ « En 1985, la Chine s'est dotée d'une loi sur les successions. Cette loi dit : les fils et les filles ont les mêmes droits d'héritage. L'épouse peut aussi hériter des biens de son mari. Autrefois, dans beaucoup d'endroits, seuls les fils pouvaient hériter des biens des parents. Aujourd'hui, tous sont égaux devant la loi. Mais dans certaines campagnes, la tradition reste très importante et certaines familles partagent encore l'héritage selon les vieilles coutumes. »
\end{textebox}
\begin{enumerate}
\item \zh{中国哪一年有了继承法？}
\item \zh{根据这个法律，儿子和女儿一样吗？}
\item \zh{妻子可以继承什么？}
\item \zh{以前，谁能继承父母的财产？}
\item \zh{现在在一些农村，有些家庭怎么分遗产？}
\end{enumerate}
\textbf{Question de langue :} relève dans le texte une phrase avec \zh{还是} et traduis-la en français.
\end{exercice}

\begin{exercice}[Thème et version — type Bac]
\textbf{A. Thème.} Traduis en chinois (caractères et pinyin) :
\begin{enumerate}
\item « Au Cameroun, la famille décide ensemble du partage de l'héritage. »
\item « Maintenant, filles et fils sont égaux devant la loi. »
\item « Je pense que la loi est plus importante que la tradition. »
\end{enumerate}
\textbf{B. Version.} Traduis en français le texte suivant :
\begin{textebox}[Version]
\zh{在这个酋长家里，按照习俗，长子继承父亲的土地和房子。其他的孩子得到一些钱和东西。女儿结婚以后，住在丈夫的家里。但是今天，有些年轻人觉得这种习俗应该改变。}\\ \emph{Zài zhège qiúzhǎng jiā lǐ, ànzhào xísú, zhǎngzǐ jìchéng fùqin de tǔdì hé fángzi. Qítā de háizi dédào yìxiē qián hé dōngxi. Nǚ'ér jiéhūn yǐhòu, zhù zài zhàngfu de jiā lǐ. Dànshì jīntiān, yǒuxiē niánqīngrén juéde zhè zhǒng xísú yīnggāi gǎibiàn.}
\end{textebox}
\end{exercice}

\begin{exercice}[Expression écrite guidée — type Bac]
\textbf{Sujet :} \zh{我对遗产继承的看法} (\emph{Wǒ duì yíchǎn jìchéng de kànfǎ} — « Mon point de vue sur l'héritage »).
\par Rédige un court texte de \textbf{5 à 8 phrases} (environ 80 à 120 caractères) en chinois. Tu dois obligatoirement utiliser :
\begin{itemize}
\item la structure d'opinion \zh{我觉得……} ;
\item une comparaison avec \zh{比} ;
\item une égalité avec \zh{和……一样} ;
\item au moins quatre mots du vocabulaire de la leçon (\zh{遗产}, \zh{继承}, \zh{传统}, \zh{法律}, \zh{平等}, \zh{习俗}…).
\end{itemize}
\textbf{Question de réflexion (à intégrer à ton texte ou à traiter à part en 3 phrases) :} \zh{你觉得继承的传统应该改变吗？为什么？} — « Penses-tu que les traditions de succession doivent évoluer ? Justifie en 3 phrases en chinois. »
\end{exercice}

\newpage

\coursec{Corrigés des exercices}

\begin{corrige}[Exercice 1 — QCM : comprendre les mots-clés]
\begin{enumerate}
\item \zh{遗产} (yíchǎn) : \textbf{b. l'héritage}. \zh{遗} évoque ce qui est laissé (par un défunt), \zh{产} les biens (comme dans \zh{财产} cáichǎn, « biens, fortune »).
\item \zh{继承} (jìchéng) : \textbf{a. hériter}. C'est un verbe : \zh{继承遗产} (jìchéng yíchǎn), « hériter d'un héritage ».
\item \zh{女儿} (nǚ'ér) : \textbf{b. la fille}. Attention à ne pas confondre avec \zh{儿子} (érzi, « le fils ») ni \zh{母亲} (mǔqin, « la mère »).
\item \zh{法律} (fǎlǜ) : \textbf{c. la loi}. La « tradition » se dit \zh{传统} (chuántǒng).
\end{enumerate}
\end{corrige}

\begin{corrige}[Exercice 2 — Vrai ou faux ?]
\begin{enumerate}
\item \textbf{Vrai (对 duì).} La loi chinoise sur les successions de 1985 établit que les fils et les filles ont les mêmes droits d'héritage : \zh{儿子和女儿有一样的继承权。}
\item \textbf{Faux (不对 bú duì).} Le Cameroun compte de nombreuses communautés et les coutumes de succession varient selon les régions et les ethnies ; il n'existe pas une coutume unique.
\item \textbf{Vrai (对 duì).} \zh{长} (zhǎng) signifie ici « aîné » et \zh{子} (zǐ) « fils » : \zh{长子} désigne bien le fils aîné. Attention à la prononciation : \zh{长} se lit zhǎng (et non cháng, « long »).
\item \textbf{Vrai (对 duì).} Dans les deux pays, malgré l'existence de lois modernes, la tradition continue d'influencer le partage de l'héritage, surtout dans les zones rurales.
\end{enumerate}
\end{corrige}

\begin{corrige}[Exercice 3 — Vocabulaire : complète le tableau]
\begin{center}
\begin{tabularx}{\linewidth}{|X|X|X|X|}
\hline
\rowcolor{popblueL} Caractères & Pinyin & Nature & Traduction \\ \hline
\zh{传统} & chuántǒng & nom & la tradition \\ \hline
\zh{平等} & píngděng & adjectif / nom & égal ; l'égalité \\ \hline
\zh{分} & fēn & verbe & partager, diviser \\ \hline
\zh{习俗} & xísú & nom & la coutume \\ \hline
\zh{权利} & quánlì & nom & le droit \\ \hline
\zh{家庭} & jiātíng & nom & la famille \\ \hline
\end{tabularx}
\end{center}
\textbf{Points de vigilance :} \zh{分} se lit fēn au sens de « partager » (il se lit fèn dans \zh{部分}, « partie »). \zh{平等} peut fonctionner comme adjectif (\zh{儿子和女儿是平等的}, « fils et filles sont égaux ») ou comme nom (\zh{男女平等}, « l'égalité hommes-femmes »).
\end{corrige}

\begin{corrige}[Exercice 4 — Associe les caractères au pinyin]
\begin{center}
\begin{tabularx}{\linewidth}{|X|X|}
\hline
\rowcolor{popblueL} Caractères & Pinyin \\ \hline
1. \zh{继承人} & \textbf{c.} jìchéngrén (l'héritier) \\ \hline
2. \zh{遗嘱} & \textbf{d.} yízhǔ (le testament) \\ \hline
3. \zh{法律} & \textbf{a.} fǎlǜ (la loi) \\ \hline
4. \zh{后代} & \textbf{b.} hòudài (la descendance) \\ \hline
5. \zh{酋长} & \textbf{e.} qiúzhǎng (le chef traditionnel) \\ \hline
\end{tabularx}
\end{center}
\textbf{Astuce de mémorisation :} \zh{继承人} = \zh{继承} (hériter) + \zh{人} (personne) ; \zh{后代} = \zh{后} (après) + \zh{代} (génération).
\end{corrige}

\begin{corrige}[Exercice 5 — Traduis en français]
\begin{enumerate}
\item \zh{在中国，儿子和女儿有平等的继承权。} — « En Chine, les fils et les filles ont des droits d'héritage égaux. »\\ Structure : \zh{有} + adjectif + \zh{的} + nom (\zh{有平等的继承权}, « avoir des droits égaux »).
\item \zh{在喀麦隆，很多家庭一起讨论怎么分遗产。} — « Au Cameroun, beaucoup de familles discutent ensemble de la manière de partager l'héritage. »\\ \zh{怎么} (zěnme) introduit ici une interrogative indirecte : « comment partager ».
\item \zh{按照传统，长子常常继承父亲的土地。} — « Selon la tradition, le fils aîné hérite souvent de la terre de son père. »\\ \zh{按照} (ànzhào) = « selon, conformément à » ; \zh{父亲的土地} : possession avec \zh{的}.
\end{enumerate}
\end{corrige}

\begin{corrige}[Exercice 6 — Texte à trous : les mots-outils]
Texte complété :
\begin{enumerate}
\item \zh{父亲\textbf{的}土地} : \zh{的} marque la possession (« la terre \emph{du} père »).
\item \zh{中国\textbf{和}喀麦隆} : \zh{和} relie deux noms (« la Chine \emph{et} le Cameroun »).
\item \zh{儿子\textbf{和}女儿} : même emploi de \zh{和} (« fils \emph{et} filles »).
\item \zh{是平等\textbf{的}} : structure \zh{是 + adjectif + 的}, tournure d'identification (« sont égaux »).
\item \zh{今天\textbf{的}情况} : \zh{的} de possession/appartenance (« la situation d'aujourd'hui »).
\item \zh{\textbf{比}以前好多了} : \zh{比} introduit le comparatif (« bien meilleure \emph{qu}'avant »).
\end{enumerate}
\textbf{Point de vigilance :} le mot-outil \zh{了} n'était pas à placer ici ; il apparaît déjà dans \zh{好多了} (hǎo duō le, « bien meilleur »), où il marque le changement d'état. Ne pas confondre avec le \zh{了} d'accompli.
\end{corrige}

\begin{corrige}[Exercice 7 — Remets les mots dans l'ordre]
\begin{enumerate}
\item \zh{女儿和儿子有一样的继承权。}\\ \emph{Nǚ'ér hé érzi yǒu yíyàng de jìchéngquán.}\\ « Les filles et les fils ont les mêmes droits d'héritage. »\\ Structure : A \zh{和} B \zh{有一样的} + nom.
\item \zh{在喀麦隆，家庭常常一起讨论怎么分遗产。}\\ \emph{Zài Kāmàilóng, jiātíng chángcháng yìqǐ tǎolun zěnme fēn yíchǎn.}\\ « Au Cameroun, les familles discutent souvent ensemble de la façon de partager l'héritage. »\\ Ordre : complément de lieu (\zh{在喀麦隆}) en tête, puis sujet, adverbes (\zh{常常一起}) avant le verbe.
\item \zh{我觉得法律比传统更重要。}\\ \emph{Wǒ juéde fǎlǜ bǐ chuántǒng gèng zhòngyào.}\\ « Je pense que la loi est plus importante que la tradition. »\\ Structure du comparatif : A + \zh{比} + B + (\zh{更}) + adjectif.
\end{enumerate}
\end{corrige}

\begin{corrige}[Exercice 8 — Transforme les phrases]
\begin{enumerate}
\item \zh{儿子\textbf{不}继承土地。} (\emph{Érzi bù jìchéng tǔdì.}) — « Le fils n'hérite pas de la terre. »\\ \zh{不} se place directement devant le verbe, pour une négation au présent ou une vérité générale.
\item \zh{他\textbf{没有}继承父亲的房子。} (\emph{Tā méiyǒu jìchéng fùqin de fángzi.}) — « Il n'a pas hérité de la maison de son père. »\\ Pour nier une action passée, on utilise \zh{没有} (ou \zh{没}) et on supprime le \zh{了} d'accompli. Ne jamais écrire \zh{没有了继承}.
\item \zh{女儿\textbf{和}儿子\textbf{一样}平等。} (\emph{Nǚ'ér hé érzi yíyàng píngděng.}) — « Les filles sont égales aux fils. »\\ Structure : A + \zh{和} + B + \zh{一样} + adjectif.
\item \zh{传统和法律一样重要。} (\emph{Chuántǒng hé fǎlǜ yíyàng zhòngyào.}) — « La tradition et la loi sont aussi importantes l'une que l'autre. »
\end{enumerate}
\end{corrige}

\begin{corrige}[Exercice 9 — Compréhension de texte — type Bac]
\textbf{Réponses attendues (phrases complètes en chinois) :}
\begin{enumerate}
\item \zh{中国一九八五年有了继承法。}\\ \emph{Zhōngguó yī jiǔ bā wǔ nián yǒu le jìchéngfǎ.} — « La Chine s'est dotée d'une loi sur les successions en 1985. »
\item \zh{根据这个法律，儿子和女儿一样，他们有一样的继承权。}\\ \emph{Gēnjù zhège fǎlǜ, érzi hé nǚ'ér yíyàng, tāmen yǒu yíyàng de jìchéngquán.} — « Selon cette loi, fils et filles sont égaux : ils ont les mêmes droits d'héritage. »
\item \zh{妻子可以继承丈夫的遗产。}\\ \emph{Qīzi kěyǐ jìchéng zhàngfu de yíchǎn.} — « L'épouse peut hériter des biens de son mari. »
\item \zh{以前，只有儿子能继承父母的财产。}\\ \emph{Yǐqián, zhǐyǒu érzi néng jìchéng fùmǔ de cáichǎn.} — « Autrefois, seuls les fils pouvaient hériter des biens des parents. »
\item \zh{现在在一些农村，有些家庭还是按照老习惯分遗产。}\\ \emph{Xiànzài zài yìxiē nóngcūn, yǒuxiē jiātíng háishi ànzhào lǎo xíguàn fēn yíchǎn.} — « Aujourd'hui, dans certaines campagnes, certaines familles partagent encore l'héritage selon les vieilles coutumes. »
\end{enumerate}
\textbf{Question de langue :} phrase avec \zh{还是} (háishi, « encore, toujours ») :\\
\zh{在一些农村，传统还是很重要。} — « Dans certaines campagnes, la tradition reste très importante. »\\ (Autre possibilité : \zh{有些家庭还是按照老习惯分遗产。} — « Certaines familles partagent encore l'héritage selon les vieilles habitudes. »)

\textbf{Barème indicatif (sur 10) :} 1,5 point par question de compréhension (contenu 1 + exactitude linguistique 0,5), soit 7,5 points ; question de langue 2,5 points (relevé correct 1 + traduction 1,5).

\textbf{Points de vigilance :} répondre par une phrase complète avec sujet et verbe, pas par un mot isolé ; reprendre le vocabulaire du texte sans le recopier mot à mot sur toute la phrase ; attention à \zh{只有} (zhǐyǒu, « seulement ») et à la place de \zh{还是} devant le verbe ou l'adjectif.
\end{corrige}

\begin{corrige}[Exercice 10 — Thème et version — type Bac]
\textbf{A. Thème — propositions de traduction :}
\begin{enumerate}
\item « Au Cameroun, la famille décide ensemble du partage de l'héritage. »\\
\zh{在喀麦隆，家人一起决定怎么分遗产。}\\ \emph{Zài Kāmàilóng, jiārén yìqǐ juédìng zěnme fēn yíchǎn.}\\ Variante acceptée : \zh{在喀麦隆，家庭一起讨论遗产的分配。}
\item « Maintenant, filles et fils sont égaux devant la loi. »\\
\zh{现在，女儿和儿子在法律面前是平等的。}\\ \emph{Xiànzài, nǚ'ér hé érzi zài fǎlǜ miànqián shì píngděng de.}\\ Variante : \zh{现在，法律面前人人平等。}
\item « Je pense que la loi est plus importante que la tradition. »\\
\zh{我觉得法律比传统更重要。}\\ \emph{Wǒ juéde fǎlǜ bǐ chuántǒng gèng zhòngyào.}
\end{enumerate}
\textbf{B. Version — traduction de référence :}\\
\textbf{Texte source (chinois) :}\\
\zh{在这个酋长的家庭里，按照习俗，长子继承父亲的土地和房子。其他孩子得到一点钱和东西。女儿结婚以后住在丈夫家里。但现在，有些年轻人觉得这个习俗应该改变。}\\
\emph{Zài zhège qiúzhǎng de jiātíng lǐ, ànzhào xísú, zhǎngzǐ jìchéng fùqin de tǔdì hé fángzi. Qítā háizi dédào yìdiǎn qián hé dōngxi. Nǚ'ér jiéhūn yǐhòu zhù zài zhàngfu de jiālǐ. Dàn xiànzài, yǒuxiē niánqīngrén juéde zhège xísú yīnggāi gǎibiàn.}

\textbf{Traduction française :}\\
« Dans la famille de ce chef traditionnel, selon la coutume, le fils aîné hérite de la terre et de la maison de son père. Les autres enfants reçoivent un peu d'argent et des biens. Après leur mariage, les filles vivent dans la maison de leur mari. Mais aujourd'hui, certains jeunes pensent que cette coutume devrait changer. »

\textbf{Barème indicatif (sur 10) :} thème 5 points (environ 1,5 par phrase : caractères 0,5, pinyin 0,5, grammaire 0,5) ; version 5 points (fidélité du sens 3, qualité du français 2).

\textbf{Points de vigilance :} « devant la loi » se rend par \zh{在法律面前} ; « décider de » = \zh{决定} + proposition ; dans la version, \zh{结婚以后} = « après le mariage » et \zh{应该改变} = « devrait changer » (ne pas traduire \zh{应该} par « peut »).
\end{corrige}

\begin{corrige}[Exercice 11 — Expression écrite guidée — type Bac]
\textbf{Proposition de rédaction modèle :}
\par
\zh{我觉得遗产继承是一个很重要的问题。以前，在很多地方，只有儿子能继承父母的财产，女儿什么都没有。现在，法律说儿子和女儿有一样的继承权，我觉得这很好。法律比传统更重要，因为法律面前人人平等。但是，我们也应该尊重一些好的传统。我希望每个家庭都能平等地分遗产。}\par
\emph{Wǒ juéde yíchǎn jìchéng shì yí ge hěn zhòngyào de wèntí. Yǐqián, zài hěn duō dìfang, zhǐyǒu érzi néng jìchéng fùmǔ de cáichǎn, nǚ'ér shénme dōu méiyǒu. Xiànzài, fǎlǜ shuō érzi hé nǚ'ér yǒu yíyàng de jìchéngquán, wǒ juéde zhè hěn hǎo. Fǎlǜ bǐ chuántǒng gèng zhòngyào, yīnwèi fǎlǜ miànqián rénrén píngděng. Dànshì, wǒmen yě yīnggāi zūnzhòng yìxiē hǎo de chuántǒng. Wǒ xīwàng měi ge jiātíng dōu néng píngděng de fēn yíchǎn.}\par
« Je pense que l'héritage est une question très importante. Autrefois, dans beaucoup d'endroits, seuls les fils pouvaient hériter des biens des parents, et les filles n'avaient rien. Aujourd'hui, la loi dit que fils et filles ont les mêmes droits d'héritage, et je trouve cela très bien. La loi est plus importante que la tradition, car tous sont égaux devant la loi. Mais nous devons aussi respecter certaines bonnes traditions. J'espère que chaque famille pourra partager l'héritage équitablement. »

\textbf{Vérification des contraintes :} \zh{我觉得} (opinion) ; \zh{法律比传统更重要} (comparaison avec \zh{比}) ; \zh{一样的继承权} (égalité avec \zh{一样}) ; vocabulaire de la leçon : \zh{遗产}, \zh{继承}, \zh{传统}, \zh{法律}, \zh{平等} (5 mots, minimum 4 exigés).

\textbf{Réponse modèle à la question de réflexion :}\\
\zh{我觉得有些传统应该改变。因为女儿和儿子一样，都是父母的孩子。平等的家庭更幸福。}\par
\emph{Wǒ juéde yìxiē chuántǒng yīnggāi gǎibiàn. Yīnwèi nǚ'ér hé érzi yíyàng, dōu shì fùmǔ de háizi. Píngděng de jiātíng gèng xìngfú.}\par
« Je pense que certaines traditions doivent changer. Parce que les filles sont comme les fils : toutes sont les enfants de leurs parents. Une famille égalitaire est plus heureuse. »

\textbf{Barème indicatif (sur 10) :} respect des contraintes imposées (4 structures/vocabulaire) 4 points ; correction grammaticale et caractères 3 points ; richesse du vocabulaire 1,5 point ; cohérence et organisation du texte 1,5 point.

\textbf{Points de vigilance :}
\begin{itemize}
\item Ne pas écrire \zh{我很觉得} : \zh{觉得} ne prend pas \zh{很}.
\item Comparatif : A + \zh{比} + B + adjectif, sans \zh{很} devant l'adjectif (\zh{比传统很重要} est incorrect ; écrire \zh{比传统更重要}).
\item \zh{和……一样} : l'adjectif se place après \zh{一样} (\zh{和儿子一样平等}).
\item Compter ses caractères : rester entre 80 et 120 caractères, phrases courtes et sûres plutôt que phrases longues et fautives.
\end{itemize}
\end{corrige}

\newpage

\coursec{Fiche bilan}

\begin{popfigure}
\begin{tikzpicture}[mindmap, grow cyclic, every node/.style={concept, font=\small},
  root concept/.append style={concept color=popblue, font=\small\bfseries, minimum size=3.2cm},
  level 1/.append style={level distance=3.6cm, sibling angle=60},
  level 1 concept/.append style={font=\footnotesize, minimum size=2.4cm}]
\node[root concept, align=center]{L'hérédité\\ 继承 \emph{jìchéng}}
  child[concept color=poporange]{ node[align=center]{Notion\\ 遗产 \emph{yíchǎn}\\ 传统 \emph{chuántǒng}} }
  child[concept color=popgreen]{ node[align=center]{Cameroun\\ 家族 \emph{jiāzú}\\ 酋长 \emph{qiúzhǎng}\\ 土地 \emph{tǔdì}} }
  child[concept color=poppurple]{ node[align=center]{Chine\\ 姓氏 \emph{xìngshì}\\ 家谱 \emph{jiāpǔ}\\ 儿子 \emph{érzi}} }
  child[concept color=popgold]{ node[align=center]{Comparaison\\ 一样 \emph{yíyàng}\\ 不同 \emph{bùtóng}\\ 比 \emph{bǐ}} }
  child[concept color=poppink]{ node[align=center]{Grammaire\\ 把字句\\ 被字句\\ 不仅……而且……} }
  child[concept color=popblue!60]{ node[align=center]{Point de vue\\ 我认为 \emph{wǒ rènwéi}\\ 我觉得 \emph{wǒ juéde}} };
\end{tikzpicture}
\legende{Carte mentale de la leçon : l'hérédité au Cameroun et en Chine}
\end{popfigure}

\begin{bilan}
\textbf{1. Lexique clé à connaître par cœur}
\begin{center}
\begin{tabularx}{\linewidth}{|X|X|X|X|}
\hline
\rowcolor{popblueL} Caractères & Pinyin & Nature & Traduction \\
\hline
继承 & jìchéng & verbe & hériter, succéder \\
\hline
遗产 & yíchǎn & nom & héritage, patrimoine \\
\hline
继承人 & jìchéngrén & nom & héritier \\
\hline
传统 & chuántǒng & nom / adj. & tradition ; traditionnel \\
\hline
习俗 & xísú & nom & coutume \\
\hline
家族 & jiāzú & nom & famille élargie, clan \\
\hline
姓氏 & xìngshì & nom & nom de famille \\
\hline
家谱 & jiāpǔ & nom & arbre généalogique \\
\hline
酋长 & qiúzhǎng & nom & chef traditionnel \\
\hline
土地 & tǔdì & nom & terre, terrain \\
\hline
财产 & cáichǎn & nom & biens, propriété \\
\hline
女儿 & nǚ'ér & nom & fille \\
\hline
权利 & quánlì & nom & droit \\
\hline
平等 & píngděng & adj. / nom & égal ; égalité \\
\hline
法律 & fǎlǜ & nom & loi \\
\hline
\end{tabularx}
\end{center}

\textbf{2. Structures de grammaire à maîtriser}
\begin{center}
\begin{tabularx}{\linewidth}{|X|X|X|}
\hline
\rowcolor{popblueL} Structure & Exemple & Traduction \\
\hline
A 把 B + verbe + 给 + C & 父亲把土地传给儿子。 \emph{Fùqin bǎ tǔdì chuán gěi érzi.} & Le père transmet la terre au fils. \\
\hline
B 被 A + verbe & 遗产被平均分配。 \emph{Yíchǎn bèi píngjūn fēnpèi.} & L'héritage est réparti équitablement. \\
\hline
A 比 B + adj. & 中国的家谱比喀麦隆的长。 \emph{Zhōngguó de jiāpǔ bǐ Kāmàilóng de cháng.} & Les généalogies chinoises sont plus longues que les camerounaises. \\
\hline
不仅……而且…… & 他不仅继承土地，而且继承名字。 \emph{Tā bùjǐn jìchéng tǔdì, érqiě jìchéng míngzi.} & Il hérite non seulement de la terre, mais aussi du nom. \\
\hline
越来越 + adj. & 女儿的权利越来越多。 \emph{Nǚ'ér de quánlì yuè lái yuè duō.} & Les droits des filles sont de plus en plus nombreux. \\
\hline
\end{tabularx}
\end{center}

\textbf{3. Faits socioculturels à retenir}
\begin{itemize}
\item Au Cameroun : l'héritage (terres, case, biens, parfois le titre de chef) se transmet souvent selon la coutume, en général aux fils ou à la famille élargie ; les pratiques varient selon les régions et les ethnies.
\item En Chine : le nom de famille \zh{姓} passe traditionnellement par le père ; le \zh{家谱} (arbre généalogique) est très valorisé ; le culte des ancêtres \zh{祖先} reste important.
\item Dans les deux pays, la loi moderne reconnaît de plus en plus l'égalité entre filles et fils dans l'héritage.
\end{itemize}

\textbf{4. Méthodes express}
\begin{itemize}
\item \cle{Comparer} : d'abord un point commun \zh{一样} (pareil), puis une différence avec \zh{比} ou \zh{不同}.
\item \cle{Donner un avis} : \zh{我认为……} / \zh{我觉得……} + phrase simple + un exemple concret.
\item \cle{Mini-exposé} : introduction (1 phrase) $\rightarrow$ Cameroun (2 phrases) $\rightarrow$ Chine (2 phrases) $\rightarrow$ comparaison (1 phrase) $\rightarrow$ conclusion personnelle.
\end{itemize}
\end{bilan}

\begin{aretenir}[Checklist avant le Bac]
\begin{itemize}
\item[$\square$] Je connais le sens et le pinyin des 15 mots clés du thème (继承, 遗产, 家谱, 酋长, 平等…).
\item[$\square$] Je sais écrire les caractères 继承, 遗产 et 传统 dans le bon ordre des traits.
\item[$\square$] Je sais construire une phrase avec 把 pour dire « transmettre quelque chose à quelqu'un ».
\item[$\square$] Je sais utiliser le passif avec 被 pour parler de la répartition d'un héritage.
\item[$\square$] Je sais comparer deux pratiques avec 比 et 不同.
\item[$\square$] Je sais dire « non seulement… mais aussi… » avec 不仅……而且…….
\item[$\square$] Je peux citer deux faits sur l'hérédité au Cameroun et deux faits sur la Chine.
\item[$\square$] Je sais exprimer mon point de vue avec 我认为 et 我觉得.
\item[$\square$] Je respecte l'ordre des mots : Sujet + temps/lieu + verbe + complément.
\item[$\square$] Je vérifie les tons du pinyin (jìchéng, yíchǎn, píngděng) dans mes productions écrites.
\item[$\square$] Je peux faire un mini-exposé de 6 à 8 phrases sur la comparaison Cameroun / Chine.
\item[$\square$] Je reste factuel et respectueux quand je compare les deux cultures.
\end{itemize}
\end{aretenir}

\findecours

\end{document}
